% fig1_framework.pdf -- Fig. 1, Section III overview.
%
% v6 (FIGURE_REDESIGN_SPEC 3.1, step S6).  The figure is now ONE message and
% ONE picture grammar: a single left-to-right pipeline from a real SAR chip to
% a real output mask stack, cut into three dotted regions (Stage A / Stage B
% rung 1 / Stage B rungs 2-3).  Squint focus = the consensus-vote block, the
% only heavy SOLID vermillion frame with a filled vermillion ground in the
% whole figure (H10 / SPEC 2.3).
%
% DELETED in v6, each with its rule (SPEC 3.1 "要删除的元素"):
%   1. \apgroup x6 (18 AP numbers), the 0-75 \baxis and the dagger footnote  H1
%   2. the three \draw[brace] and the four error-axis notes                  H4
%   3. the whole "Inside the frozen base" expansion band  -> fig1b_readout   A1
%   4. the five \legendswatch rows at the foot (colour key -> caption)       H4
%   5. the five numbered badges (\badgecar/\badgeops/\badgetrn)              P3
%   6. the second lock glyph and the three \ckpt slabs                       P3
%   7. every 3rd/4th line of every block (hyper-parameters, timings,
%      parameter counts, the \safefrac formula)                          H2/H3
% The only numbers left are counts and one tensor shape: "3 seeds",
% "800x800", the rung indices 1 and 2--3.  No AP, no epoch, no runtime.
%
% REAL DATA (SPEC 2.5).  Both rasters are the actual tensors of one forward
% pass of runs/c0_maskiou_v2/soup_s023/iter_100872.pth on HRSID inshore chip
% P0101_1200_2000_3000_3800.jpg (image_id 1528), dumped by
% carrier/dump_figure_assets.py; provenance in assets/manifest.json.  Both are
% 800 px drawn at 20.0 mm = 1016 dpi, embedded through \assetdir.
%
% Authored at exactly 18.2 x 4.00 cm = 516 x 113.4 pt = \textwidth, 4.55:1
% -> \includegraphics[width=\textwidth] inside figure*[!t], scale 1.0, so
% \footnotesize is a true 8 pt and \scriptsize a true 7 pt on the page.
%
% THREE DOCUMENTED DEVIATIONS from the letter of SPEC 3.1 (all forced by a
% measured collision or by a SPEC 4.2 gate; see the S6 report):
%   D1  dash budget.  SPEC 2.2 allows exactly two dash patterns per figure
%       (trn's 2/1.2 pt and grp's densely dotted); tikz_style_v5's carrier
%       style adds a third (2.6/1.2 pt).  The carrier is folded onto the trn
%       pattern and stays separable by weight (0.5 vs 0.8 pt) and hue
%       (vermillion vs green); its fill goes white so the consensus operator
%       is the only FILLED vermillion block on the rung-1 strip.
%   D2  colour budget (<= 10 distinct vector colours).  The chip frame moves
%       from black!55 to datablue!60, the colour the output tensor pile
%       already uses, so input and output raster share one "data" hue; and
%       the grp frame moves from black!40 to frozgray, which SPEC 2.3 already
%       assigns to chrome.  Net: 13 colours -> 10.
%   D3  two labels do not fit the SPEC 3.1 box widths in 8/7 pt Times
%       (measured, not estimated).  "Boundary refiner" is 59.19 pt = 2.08 cm
%       against a 2.00 cm box, so that block grows to 22.5 mm and the two
%       blocks to its right shift 0.25 cm right (all gaps stay at the spec
%       0.25/0.30 cm and the canvas is unchanged).  "input space, 8 views" is
%       56.59 pt = 1.99 cm against an 18 mm box, so the count is dropped and
%       the sub-label reads "input space", keeping the weight/input/output
%       space parallel that carries the message; "eight-view" is already in
%       the caption.
\expandafter\def\csname ver@epstopdf-base.sty\endcsname{2020/01/24}
\documentclass[tikz,border=0pt]{standalone}
\usepackage{mathptmx}
\usepackage[T1]{fontenc}
\usepackage{amsmath,amssymb}
% =========================================================================
% tikz_style_v5.tex -- DRAFT style set (R4 tooling report, 2026-09-03).
% NOT wired into any figure.  It loads tikz_style.tex unchanged (every v4
% macro, palette entry and glyph is kept) and then OVERRIDES the chrome:
%
%   T1  two type sizes only (8 pt labels / 7 pt sub-labels) and a 6 pt FLOOR
%       for math sub/superscripts -- LaTeX's size10.clo declares
%       \DeclareMathSizes{8}{8}{6}{5} and {7}{7}{5}{5}, so every ^2 in a
%       7/8 pt label was printing at 4.98 pt (figs/verify_fig.py, all three
%       method figures).  IEEE asks for ~9-10 pt at full size; 6 pt is the
%       absolute floor any surveyed venue tolerates (SIGGRAPH guide,
%       Elsevier subscripts, IEEE Standards style).
%   T2  two stroke weights only: 0.5 pt (ordinary) / 0.8 pt (emphasis).
%       0.5 TeX pt = 0.498 PostScript pt in the PDF, which verify_fig.py
%       tolerates; nothing may go below.  (Science: "0.5 pt minimum line
%       weight at final reduced size".)
%   T3  one corner radius (1.5 pt) for every block, one dotted group frame,
%       one arrowhead (Stealth 1.6 x 1.2 mm) at 0.5 pt, emphasis arrows 0.8 pt.
%   T4  four semantic hues + gray, nothing else:
%         opsred   D55E00  training-free operator (our contribution)
%         trngreen 009E73  trained corrector       (optional, ours)
%         datablue 0072B2  data / tensors / carriers
%         warnred  9E2B2B  failure mode / negative control
%         frozgray 6E6E6E  frozen base, chrome, de-emphasis
%       Hue is always doubled by weight/dash/glyph so the grayscale proof
%       (colour online / B&W in print) still separates the roles.
%   T5  PlotNeuralNet-style 3D tensor pics (\pic {tensor=...}, {tensorpile=...})
%       whose geometry is DERIVED from the tensor shape (G1/G2 of
%       tikz_style.tex), not hand-typed; a real raster may sit on the front
%       face.  Anchors NAME-west/-east/... let arrows attach without maths.
%   T6  \legendswatch : one macro for the in-figure legend row.
%
% Usage (one-line change per figure):   % =========================================================================
% tikz_style_v5.tex -- DRAFT style set (R4 tooling report, 2026-09-03).
% NOT wired into any figure.  It loads tikz_style.tex unchanged (every v4
% macro, palette entry and glyph is kept) and then OVERRIDES the chrome:
%
%   T1  two type sizes only (8 pt labels / 7 pt sub-labels) and a 6 pt FLOOR
%       for math sub/superscripts -- LaTeX's size10.clo declares
%       \DeclareMathSizes{8}{8}{6}{5} and {7}{7}{5}{5}, so every ^2 in a
%       7/8 pt label was printing at 4.98 pt (figs/verify_fig.py, all three
%       method figures).  IEEE asks for ~9-10 pt at full size; 6 pt is the
%       absolute floor any surveyed venue tolerates (SIGGRAPH guide,
%       Elsevier subscripts, IEEE Standards style).
%   T2  two stroke weights only: 0.5 pt (ordinary) / 0.8 pt (emphasis).
%       0.5 TeX pt = 0.498 PostScript pt in the PDF, which verify_fig.py
%       tolerates; nothing may go below.  (Science: "0.5 pt minimum line
%       weight at final reduced size".)
%   T3  one corner radius (1.5 pt) for every block, one dotted group frame,
%       one arrowhead (Stealth 1.6 x 1.2 mm) at 0.5 pt, emphasis arrows 0.8 pt.
%   T4  four semantic hues + gray, nothing else:
%         opsred   D55E00  training-free operator (our contribution)
%         trngreen 009E73  trained corrector       (optional, ours)
%         datablue 0072B2  data / tensors / carriers
%         warnred  9E2B2B  failure mode / negative control
%         frozgray 6E6E6E  frozen base, chrome, de-emphasis
%       Hue is always doubled by weight/dash/glyph so the grayscale proof
%       (colour online / B&W in print) still separates the roles.
%   T5  PlotNeuralNet-style 3D tensor pics (\pic {tensor=...}, {tensorpile=...})
%       whose geometry is DERIVED from the tensor shape (G1/G2 of
%       tikz_style.tex), not hand-typed; a real raster may sit on the front
%       face.  Anchors NAME-west/-east/... let arrows attach without maths.
%   T6  \legendswatch : one macro for the in-figure legend row.
%
% Usage (one-line change per figure):   % =========================================================================
% tikz_style_v5.tex -- DRAFT style set (R4 tooling report, 2026-09-03).
% NOT wired into any figure.  It loads tikz_style.tex unchanged (every v4
% macro, palette entry and glyph is kept) and then OVERRIDES the chrome:
%
%   T1  two type sizes only (8 pt labels / 7 pt sub-labels) and a 6 pt FLOOR
%       for math sub/superscripts -- LaTeX's size10.clo declares
%       \DeclareMathSizes{8}{8}{6}{5} and {7}{7}{5}{5}, so every ^2 in a
%       7/8 pt label was printing at 4.98 pt (figs/verify_fig.py, all three
%       method figures).  IEEE asks for ~9-10 pt at full size; 6 pt is the
%       absolute floor any surveyed venue tolerates (SIGGRAPH guide,
%       Elsevier subscripts, IEEE Standards style).
%   T2  two stroke weights only: 0.5 pt (ordinary) / 0.8 pt (emphasis).
%       0.5 TeX pt = 0.498 PostScript pt in the PDF, which verify_fig.py
%       tolerates; nothing may go below.  (Science: "0.5 pt minimum line
%       weight at final reduced size".)
%   T3  one corner radius (1.5 pt) for every block, one dotted group frame,
%       one arrowhead (Stealth 1.6 x 1.2 mm) at 0.5 pt, emphasis arrows 0.8 pt.
%   T4  four semantic hues + gray, nothing else:
%         opsred   D55E00  training-free operator (our contribution)
%         trngreen 009E73  trained corrector       (optional, ours)
%         datablue 0072B2  data / tensors / carriers
%         warnred  9E2B2B  failure mode / negative control
%         frozgray 6E6E6E  frozen base, chrome, de-emphasis
%       Hue is always doubled by weight/dash/glyph so the grayscale proof
%       (colour online / B&W in print) still separates the roles.
%   T5  PlotNeuralNet-style 3D tensor pics (\pic {tensor=...}, {tensorpile=...})
%       whose geometry is DERIVED from the tensor shape (G1/G2 of
%       tikz_style.tex), not hand-typed; a real raster may sit on the front
%       face.  Anchors NAME-west/-east/... let arrows attach without maths.
%   T6  \legendswatch : one macro for the in-figure legend row.
%
% Usage (one-line change per figure):   \input{tikz_style_v5}
% instead of \input{tikz_style}.  Nothing else in the figure has to change;
% the block/arrow style names are the same, only their numbers move.
% =========================================================================
\input{tikz_style}

% ---------------- T1: type ----------------
% Script sizes for the two text sizes used in figures.  Must be in the
% preamble; standalone + \input keeps it there.
\DeclareMathSizes{8}{8}{6}{6}
\DeclareMathSizes{7}{7}{6}{6}

% ---------------- T2/T3/T4: block, frame and arrow chrome ----------------
\tikzset{
  fig/.style={font=\footnotesize, line width=0.5pt,
              >={Stealth[length=1.6mm,width=1.2mm]},
              every node/.append style={align=center, inner sep=2.4pt}},
  % blocks: one radius, two weights; hue = role, dash = redundant channel
  blk/.style   ={rounded corners=1.5pt, inner sep=3pt, align=center},
  ops/.style   ={blk, draw=opsred,   fill=opsred!14,   line width=0.8pt},
  trn/.style   ={blk, draw=trngreen, fill=trngreen!14, line width=0.8pt,
                 dash pattern=on 2pt off 1.2pt},
  froz/.style  ={blk, draw=frozgray, fill=black!7,     line width=0.5pt},
  data/.style  ={blk, draw=datablue, fill=datablue!12, line width=0.5pt},
  plain/.style ={blk, draw=black!45, fill=white,       line width=0.5pt},
  warn/.style  ={blk, draw=warnred,  fill=warnred!8,   line width=0.5pt,
                 dash pattern=on 1pt off 1pt},
  % a carrier = established practice the operator rides on (fig1 v5 local
  % style, promoted here so every figure spells it the same way)
  carrier/.style={blk, draw=opsred, fill=opsred!5, line width=0.5pt,
                 dash pattern=on 2.6pt off 1.2pt},
  gate/.style  ={draw=opsred, fill=opsred!14, line width=0.8pt,
                 diamond, aspect=2.15, inner sep=1pt, align=center},
  grp/.style   ={draw=black!40, densely dotted, line width=0.5pt,
                 rounded corners=2pt, inner sep=0pt},
  grptitle/.style={font=\footnotesize\bfseries, text=black!80, inner sep=1.5pt},
  % type: exactly two sizes; sub-labels at 75% black (contrast >= 4.5:1)
  ttl/.style   ={font=\footnotesize\bfseries},
  lbl/.style   ={font=\footnotesize},
  sub/.style   ={font=\scriptsize, text=black!75},
  subk/.style  ={font=\footnotesize, text=black!80},
  qty/.style   ={font=\footnotesize, text=black!80},
  tsr/.style   ={font=\scriptsize\ttfamily, text=black!75},
  panl/.style  ={font=\footnotesize\bfseries, text=black!85, inner sep=1pt},
  badgecar/.style={badge, draw=opsred, fill=white, text=opsred, line width=0.5pt},
  % arrows: one head, 0.5 pt flow, 0.8 pt emphasis, dashed auxiliary
  flow/.style  ={-{Stealth[length=1.6mm,width=1.2mm]}, draw=flowink,
                 line width=0.5pt},
  flowt/.style ={flow, rounded corners=1.5pt},
  flowh/.style ={-{Stealth[length=1.9mm,width=1.4mm]}, draw=flowink,
                 line width=0.8pt},
  aux/.style   ={-{Stealth[length=1.4mm,width=1.1mm]}, draw=black!45,
                 line width=0.5pt, densely dashed},
  auxw/.style  ={aux, draw=warnred!75},
  rail/.style  ={draw=black!30, line width=0.5pt},
  brace/.style ={decorate, decoration={brace, amplitude=3pt, raise=1pt},
                 draw=black!50, line width=0.5pt},
  zoomray/.style={draw=opsred!55, line width=0.5pt, dash pattern=on 1.3pt off 1.1pt},
  zoombox/.style={draw=opsred, line width=0.5pt},
}

% ---------------- T7: rules inside math must obey the stroke floor --------
% LaTeX draws a \frac bar at the math extension font's default_rule_thickness,
% which SCALES WITH THE TYPE SIZE: 0.4 pt at 10 pt becomes 0.279 pt at the 7 pt
% \scriptsize used for figure sub-labels -- the single sub-0.5 pt stroke that
% verify_fig.py reports in fig1_framework.pdf (a \frac{1}{3} at (166.7,38.3)).
% \genfrac states the rule thickness explicitly, so it survives the floor.
% Measured: \frac -> 0.279 PDF pt, \safefrac -> 0.498 PDF pt (= 0.5 TeX pt).
% The alternative with no rule at all is to write $(1/3)\sum_s\theta_s$.
\newcommand{\safefrac}[2]{\genfrac{}{}{0.5pt}{}{#1}{#2}}

% ---------------- T8: no Type-3 bitmap fonts ----------------------------
% \ttfamily under [T1]{fontenc} with no Type-1 EC typewriter installed makes
% pdflatex fall back to the ec-tt BITMAP font (ectt0800.pk), which is embedded
% as a TYPE 3 font.  Measured on figm_tta.pdf: "F74  Type 3  Custom  emb=yes
% sub=no uni=no" (pdffonts), and figm_tta.log line 655 shows \ectt0800.pk.
% IEEE asks authors to "either convert all text to outlines or embed the
% fonts" and PDF eXpress flags Type 3; the tsr/ style below sets \ttfamily, so
% the substitution must be global.  lmtt is Latin Modern Typewriter: a real
% Type-1 T1 face.  A/B measured: without -> "F39 Type 3", with -> "LMMono8
% Type 1".  Only the typewriter family moves; mathptmx keeps Times for text.
\renewcommand{\ttdefault}{lmtt}

% ---------------- T6: legend row ----------------
% \legendswatch{x}{y}{style}{text} : a 4.2 x 3 mm swatch in <style> with its
% 7 pt label to the right.  Lay the swatches on one baseline y.
\newcommand{\legendswatch}[4]{%
  \draw[#3] (#1,{#2-0.15}) rectangle ({#1+0.42},{#2+0.15});
  \node[sub, anchor=west] at ({#1+0.54},#2) {#4};}

% =========================================================================
% T5: 3D tensor pics (PlotNeuralNet idiom, geometry from the tensor shape)
%
%   \pic (c2) at (x,y) {tensor={px=200, ch=256, img=feat_c2_x4.png,
%                               caption={$256@200^2$}}};
%   \draw[flow] (x-east) -- (c2-west);
%
% keys (all optional; defaults in brackets)
%   px       spatial side in pixels           -> drawn side = zoom*px/400 cm  [200]
%   ch       channel count                    -> depth = 0.04+0.055*log2(ch/256) cm,
%                                                floored at 0.04 cm           [256]
%   depth    explicit depth in cm, overrides ch (hand-sized blocks, UNet)   [ ]
%   zoom     magnification inside a zoom inset (G4)                          [1]
%   color    any xcolor expression; tints are derived from it               [datablue]
%   img      raster placed on the FRONT face, path relative to \assetdir     [ ]
%   hidden   1 = draw the three hidden edges dashed (the PlotNeuralNet look) [0]
%   caption  7 pt text set under the front face                             [ ]
% anchors (prefix = pic name)
%   -west -east          mid-points of the front face's left / right edge (arrows)
%   -westc -eastc        the same at half depth (for rails behind the block)
%   -north -south        mid-points of the front face's top / bottom edge
%   -front -back         centre of the front / back face
%   -nw -ne -se -sw      front-face corners;  -topne  far top-right corner
% Projection: front face in the page plane (so \includegraphics needs no
% transform), depth receding at 30 degrees to the upper right, identical to
% \fmslab, so the two glyph families can share one figure.
% =========================================================================
\tikzset{
  tensor/.pic={
    \tikzset{/tensor/.cd,#1}
    \colorlet{tnc}{\tncolor}
    \pgfmathsetmacro{\tnW}{\tnzoom*(\tnpx)/400}
    \pgfmathsetmacro{\tnH}{0.5*\tnW}
    \edef\tnDepthUser{\tndepth}
    \ifx\tnDepthUser\empty
      \pgfmathsetmacro{\tnD}{\tnzoom*max(0.04, 0.04+0.055*ln((\tnch)/256)/ln(2))}
    \else
      \pgfmathsetmacro{\tnD}{\tndepth}
    \fi
    \begin{scope}[x={(1cm,0cm)}, y={(0cm,1cm)}, z={(-0.866cm,-0.5cm)}]
      % one light source: top brightest, right middle, front darkest
      \fill[tnc!30] (-\tnH,\tnH,0) -- (-\tnH,\tnH,-\tnD) -- (\tnH,\tnH,-\tnD) -- (\tnH,\tnH,0) -- cycle;
      \fill[tnc!20] (\tnH,-\tnH,0) -- (\tnH,-\tnH,-\tnD) -- (\tnH,\tnH,-\tnD) -- (\tnH,\tnH,0) -- cycle;
      \fill[tnc!11] (-\tnH,-\tnH,0) -- (\tnH,-\tnH,0) -- (\tnH,\tnH,0) -- (-\tnH,\tnH,0) -- cycle;
      \edef\tnImgFile{\tnimg}
      \ifx\tnImgFile\empty\else
        \node[inner sep=0pt, outer sep=0pt] at (0,0,0)
          {\includegraphics[width=\tnW cm, height=\tnW cm]{\assetdir/\tnimg}};
      \fi
      \ifnum\tnhidden=1
        \draw[tnc!60, line width=0.5pt, densely dashed]
          (-\tnH,-\tnH,0) -- (-\tnH,-\tnH,-\tnD) -- (\tnH,-\tnH,-\tnD)
          (-\tnH,-\tnH,-\tnD) -- (-\tnH,\tnH,-\tnD);
      \fi
      \draw[tnc!75, line width=0.5pt]
        (-\tnH,-\tnH,0) -- (\tnH,-\tnH,0) -- (\tnH,\tnH,0) -- (-\tnH,\tnH,0) -- cycle
        (-\tnH,\tnH,0) -- (-\tnH,\tnH,-\tnD) -- (\tnH,\tnH,-\tnD) -- (\tnH,\tnH,0)
        (\tnH,\tnH,-\tnD) -- (\tnH,-\tnH,-\tnD) -- (\tnH,-\tnH,0);
      \coordinate (-front) at (0,0,0);
      \coordinate (-back)  at (0,0,-\tnD);
      \coordinate (-west)  at (-\tnH,0,0);
      \coordinate (-east)  at (\tnH,0,0);
      \coordinate (-westc) at (-\tnH,0,-0.5*\tnD);
      \coordinate (-eastc) at (\tnH,0,-0.5*\tnD);
      \coordinate (-north) at (0,\tnH,0);
      \coordinate (-south) at (0,-\tnH,0);
      \coordinate (-nw) at (-\tnH,\tnH,0);  \coordinate (-ne) at (\tnH,\tnH,0);
      \coordinate (-sw) at (-\tnH,-\tnH,0); \coordinate (-se) at (\tnH,-\tnH,0);
      \coordinate (-topne) at (\tnH,\tnH,-\tnD);
    \end{scope}
    \edef\tnCap{\tncaption}
    \ifx\tnCap\empty\else
      \node[sub, anchor=north, inner sep=1.5pt] at (0,-\tnH) {\tncaption};
    \fi
  },
  % a pile of k offset H x W planes (mask stack); front plane may carry a raster
  %   \pic (m) at (x,y) {tensorpile={px=800, ch=3, img=pred_masks.png}};
  % here ch = number of planes drawn.
  tensorpile/.pic={
    \tikzset{/tensor/.cd,#1}
    \colorlet{tnc}{\tncolor}
    \pgfmathsetmacro{\tpW}{\tnzoom*(\tnpx)/400}
    \pgfmathsetmacro{\tpH}{0.5*\tpW}
    \pgfmathsetmacro{\tpO}{0.06*\tnzoom}
    \pgfmathtruncatemacro{\tpN}{\tnch-1}
    \foreach \i in {\tpN,...,0} {
      \pgfmathsetmacro{\d}{\i*\tpO}
      \fill[white]  ({-\tpH+\d},{-\tpH+\d}) rectangle ({\tpH+\d},{\tpH+\d});
      \fill[tnc!8]  ({-\tpH+\d},{-\tpH+\d}) rectangle ({\tpH+\d},{\tpH+\d});
      \ifnum\i=0 \edef\tpImg{\tnimg}\ifx\tpImg\empty\else
        \node[inner sep=0pt, outer sep=0pt] at (0,0)
          {\includegraphics[width=\tpW cm, height=\tpW cm]{\assetdir/\tnimg}};\fi\fi
      \draw[tnc!60, line width=0.5pt] ({-\tpH+\d},{-\tpH+\d}) rectangle ({\tpH+\d},{\tpH+\d});
    }
    \coordinate (-west)  at (-\tpH,0);
    \coordinate (-east)  at ({\tpH+\tpN*\tpO},0);
    \coordinate (-north) at (0,{\tpH+\tpN*\tpO});
    \coordinate (-south) at (0,-\tpH);
    \coordinate (-front) at (0,0);
    \edef\tnCap{\tncaption}
    \ifx\tnCap\empty\else
      \node[sub, anchor=north, inner sep=1.5pt] at (0,-\tpH) {\tncaption};
    \fi
  },
  /tensor/.search also={/tikz},
  /tensor/.cd,
  px/.store in=\tnpx,           px=200,
  ch/.store in=\tnch,           ch=256,
  zoom/.store in=\tnzoom,       zoom=1,
  color/.store in=\tncolor,     color=datablue,
  img/.store in=\tnimg,         img=,
  depth/.store in=\tndepth,     depth=,
  hidden/.store in=\tnhidden,   hidden=0,
  caption/.store in=\tncaption, caption=,
}

% instead of % =========================================================================
% tikz_style.tex — shared visual system for the Section III method figures.
% Usage: \input{tikz_style} in the preamble of a standalone figure.
%
% -------------------------------------------------------------------------
% DESIGN NOTES — conventions extracted from an accepted TGRS manuscript
% (Gao et al., "BBox-free ... SAR ship instance segmentation", TGRS 2024;
%  local copy: scratchpad/gao_bboxfree_tgrs24.pdf).  Observations taken from
%  its Fig. 2 (overall architecture, p. 5), Fig. 5 (FD module, p. 8) and
%  Fig. 6 (conceptual mask panel, p. 8):
%
%  O1  REGION PARTITION.  The architecture figure is carved into 3-4 named
%      regions by DOTTED rectangles with no fill; the region name is set in
%      bold OUTSIDE the rectangle ("Feature Extractor", "Classification
%      Branch", "Localization Branch", "Center Point Extractor").  Nothing
%      is nested more than two levels deep.
%  O2  COLOR CARRIES MEANING, NOT DECORATION.  Only ~5 hues appear, all as
%      pale fills with black text: light gray = ordinary/standard module,
%      salmon-orange = the module the paper contributes (Feature Decoupling
%      Module, Center Point Extractor), yellow = a distinguished operation
%      (Feature Fusion), light blue / light pink = tensors.  Saturated color
%      is reserved for the few glyphs that must be found first.
%  O3  LEGEND PANEL.  Both Fig. 2 and Fig. 5 carry an explicit legend inside
%      the figure (dotted box top-left in Fig. 2, dashed box right in Fig. 5)
%      that decodes arrow colors and single-letter glyph nodes.  The body of
%      the diagram is kept uncluttered by pushing vocabulary into the legend.
%  O4  TENSOR ICONOGRAPHY.  Feature maps are drawn as thin vertical plates,
%      multi-scale features as 3 offset stacked planes, channel vectors as
%      segmented 3-D bars, the backbone as an isometric slab.  Real SAR
%      chips / heat maps / GT overlays are pasted directly into the flow.
%  O5  TENSOR NAMES SIT BESIDE THE GLYPH, NEVER INSIDE IT, in math italic
%      (F_out, F_mask, F_param, F_H_zip ...), placed above-right of the plate.
%  O6  ARROW DISCIPLINE.  Data flow is a thin gray/black line with a small
%      solid triangular head; routing is orthogonal (Manhattan) with a shared
%      trunk when one source fans out to several targets, so lines never
%      cross.  Semantically special transitions get a THICK COLORED BLOCK
%      arrow instead (2x upsample / 2x downsample / connect), i.e. arrow
%      style is itself a semantic channel.  Long skip paths are routed as
%      rails around the outside of the diagram (Fig. 5 top and bottom).
%  O7  TYPOGRAPHY.  One size for block labels, one smaller size for
%      sub-labels, math only in math mode; no stretched or rotated text
%      except a rare 90-degree axis label.
%  O8  ASPECT RATIO.  Double-column method figures are wide and short:
%      Fig. 2 is 530 x 200 pt (2.65:1), Fig. 5 is 530 x 215 pt (2.5:1).
%      They are laid on a grid with equal gutters between parallel branches.
%  O9  CAPTION DIVISION OF LABOR.  The caption names the parts and decodes
%      abbreviations/colors; it does not re-derive the mechanism, which is
%      left to the body text.  The figure is readable without the caption.
%  O10 CONCEPTUAL PANELS (Fig. 6) use dashed sub-panels with a centered
%      label underneath, a pale blue blob for the true object and a pale red
%      overlay for the error region -- schematic, not photorealistic.
%
% -------------------------------------------------------------------------
% WHAT WE ADOPT / WHERE WE DEVIATE
%  + adopt O1 (dotted named regions), O3 (in-figure legend), O4/O5 (tensor
%    glyphs with the symbol outside), O6 (one flow style + Manhattan
%    routing), O8 (wide-short double-column ratios), O10 (schematic panels).
%  ~ deviate on O2: instead of hue-as-identity we bind hue to ROLE in the
%    contribution -- frozen / training-free operator / trained corrector /
%    data -- because the entire argument of this article is which parts are
%    touched and which are not.
%  ~ deviate on the sans-serif block labels of Fig. 2: we set every figure in
%    the body Times face (as the same paper's Fig. 6 does), so figure text
%    and caption text share one face.
%  + every figure is authored at its EXACT printed width (18.2 cm double
%    column = 516 pt, 8.89 cm single column = 252 pt) and included at
%    scale 1, so \footnotesize is a true 8 pt and \scriptsize a true 7 pt on
%    the page in ALL figures.  Use \figcanvasd / \figcanvasc below.
%
% -------------------------------------------------------------------------
% SEMANTIC PALETTE (Okabe-Ito colorblind-safe base; redundant encoding by
% border weight, dash pattern and badge glyph so it survives grayscale):
%   froz  gray   + lock glyph        frozen, not modified by us  (Stage A)
%   ops   vermil + solid heavy rule  our TRAINING-FREE operator  (rung 1)
%   trn   green  + dashed rule       our OPTIONAL TRAINED corrector (2-3)
%   data  blue   + thin rule         data, tensors, I/O
%   warn  red    + dotted rule, small  failure mode / negative control
% =========================================================================

% Real rasters are embedded on the front faces of the tensor glyphs (see
% \realpanel / \fmslab below).  graphicx must be loaded in the preamble, which
% is where this file is \input.
\usepackage{graphicx}
% Single switch for the raster payload: \def\assetdir{assets_light}
% before \begin{document} swaps every embedded raster for its
% print-sized ~600 dpi copy (see figs/make_light_assets.py).  The
% default is the native dump, 965-1016 dpi at the drawn size.
\providecommand{\assetdir}{assets}
\usetikzlibrary{positioning,arrows.meta,fit,backgrounds,calc,shapes.geometric,
                decorations.pathreplacing,decorations.markings,patterns}

% ---------------- palette ----------------
\definecolor{opsred}{HTML}{D55E00}   % Okabe-Ito vermillion  -- training-free ops
\definecolor{trngreen}{HTML}{009E73} % Okabe-Ito bluish green-- trained correctors
\definecolor{datablue}{HTML}{0072B2} % Okabe-Ito blue        -- data / tensors
\definecolor{frozgray}{HTML}{6E6E6E} % neutral               -- frozen base
\definecolor{flowink}{HTML}{262626}  % near black            -- data-flow arrows
\definecolor{warnred}{HTML}{9E2B2B}  % muted brick           -- failure modes
\definecolor{sarink}{HTML}{2B2B33}   % SAR chip background

% ---------------- exact-width canvases ----------------
% \figcanvasd{H} : double column, 18.2 cm = 516 pt = \textwidth  (scale 1.0)
% \figcanvasc{H} : single column,  8.89 cm = 252 pt = \columnwidth
\newcommand{\figcanvasd}[1]{\useasboundingbox (0,0) rectangle (18.2,#1);}
\newcommand{\figcanvasc}[1]{\useasboundingbox (0,0) rectangle (8.89,#1);}

\tikzset{
  % ---- picture-level defaults: ONE font size system for all five figures ----
  fig/.style={font=\footnotesize, line width=0.65pt,
              >={Stealth[length=1.7mm,width=1.35mm]},
              every node/.append style={align=center, inner sep=2.4pt}},
  % ---- blocks: role -> color + redundant border encoding ----
  % Fill tints were raised in v3: at 4-8% the five roles were separated by an
  % RGB distance of only 6.6-21 (out of 441) and all landed on gray 243-248,
  % i.e. the fill channel was consuming ink and returning nothing.  The border
  % weight / dash pattern / badge remain the PRIMARY, grayscale-safe channel.
  ops/.style   ={draw=opsred,   fill=opsred!16,   line width=0.9pt,
                 rounded corners=1.6pt, inner sep=3pt, align=center},
  trn/.style   ={draw=trngreen, fill=trngreen!15, line width=0.9pt,
                 dash pattern=on 1.7pt off 1.1pt,
                 rounded corners=1.6pt, inner sep=3pt, align=center},
  froz/.style  ={draw=frozgray, fill=black!8,     line width=0.7pt,
                 rounded corners=1.6pt, inner sep=3pt, align=center},
  data/.style  ={draw=datablue, fill=datablue!14, line width=0.55pt,
                 rounded corners=1.6pt, inner sep=3pt, align=center},
  plain/.style ={draw=black!45, fill=white,       line width=0.55pt,
                 rounded corners=1.6pt, inner sep=3pt, align=center},
  % v4: radius 1.4 -> 1.6 pt, the one block style that did not share the
  % system's corner radius; weight 0.5 -> 0.6 pt so the failure chips are not
  % the faintest object on the page while being the ones a reader must find.
  % The dotted dash pattern stays as the grayscale-safe channel.
  warn/.style  ={draw=warnred,  fill=warnred!9,   line width=0.6pt,
                 dash pattern=on 1.1pt off 1.1pt,
                 rounded corners=1.6pt, inner sep=3pt, align=center},
  gate/.style  ={draw=opsred,   fill=opsred!16,   line width=0.9pt,
                 diamond, aspect=2.15, inner sep=1pt, align=center},
  % ---- named region (O1): dotted, unfilled, label set outside ----
  grp/.style   ={draw=black!45, densely dotted, line width=0.7pt,
                 rounded corners=2pt, inner sep=0pt},
  grptitle/.style={font=\footnotesize\bfseries, text=black!80, inner sep=1.5pt},
  % ---- step badges (redundant, grayscale-safe stage identity) ----
  badge/.style ={circle, inner sep=0pt, minimum size=3.3mm, text=white,
                 font=\scriptsize\bfseries},
  badgeops/.style={badge, fill=opsred},
  badgetrn/.style={badge, fill=trngreen},
  % ---- typography (O7): exactly three text sizes in the whole figure set ----
  % v3: 7 pt at 55% black was the worst corner of the design space (minimum
  % size AND minimum contrast) -- sub-labels are now 7 pt at 70% black, and
  % anything carrying a QUANTITY is promoted to 8 pt via \qty / \subk.
  ttl/.style   ={font=\footnotesize\bfseries},
  lbl/.style   ={font=\footnotesize},
  sub/.style   ={font=\scriptsize, text=black!70},
  subk/.style  ={font=\footnotesize, text=black!80},
  qty/.style   ={font=\footnotesize, text=black!80},% tensor dims, AP numbers
  tsr/.style ={font=\scriptsize\ttfamily, text=black!70},% B x C x H x W (O5)
  tag/.style   ={font=\scriptsize, inner sep=1.6pt, rounded corners=1pt},
  % panel letters (B11): \footnotesize\bfseries at the panel's top-left
  panl/.style  ={font=\footnotesize\bfseries, text=black!85, inner sep=1pt},
  % ---- arrows (O6): one data-flow style, one auxiliary/optional style ----
  flow/.style  ={-{Stealth[length=1.7mm,width=1.35mm]}, draw=flowink,
                 line width=0.65pt},
  flowt/.style ={-{Stealth[length=1.7mm,width=1.35mm]}, draw=flowink,
                 line width=0.65pt, rounded corners=1.6pt},
  aux/.style   ={-{Stealth[length=1.5mm,width=1.2mm]}, draw=black!45,
                 line width=0.5pt, densely dashed},
  auxw/.style  ={-{Stealth[length=1.5mm,width=1.2mm]}, draw=warnred!75,
                 line width=0.5pt, densely dashed},
  rail/.style  ={draw=black!30, line width=0.5pt},
  brace/.style ={decorate, decoration={brace, amplitude=3.2pt, raise=1pt},
                 draw=black!50, line width=0.5pt},
  % ---- ONE leader style for every magnifier cone in the figure set --------
  % v4: the zoom cones used to be drawn in the generic dashed `aux' style, so
  % a sight line was typographically the same object as an ordinary auxiliary
  % arrow.  A cone is now its own object: it carries the hue of the locator
  % rectangle it starts from, no arrowhead (it is a correspondence, not a
  % flow), and it always runs corner-to-corner from the locator box to the
  % MAGNIFIED RASTER, never to the enclosing panel frame.
  zoomray/.style={draw=opsred!55, line width=0.5pt,
                  dash pattern=on 1.3pt off 1.1pt},
  zoombox/.style={draw=opsred, line width=0.7pt},
}

% =========================================================================
% \hbarrow{x0}{y}{track}{frac}{colour}{weight}{label}
%   One horizontal bar on a printed 0..1 track, with the numeral set at the
%   bar end.  v4: fig1's AP progression, figm_tta (c), figm_fusion's result
%   panel and figm_rerank (b) all drew this same object with four slightly
%   different geometries.  They now share one macro, so the reader learns the
%   idiom once and meets it four times -- a small multiple across figures.
%   The empty track is always drawn, so a short bar still shows its scale.
%
%   DEPRECATED FOR METHOD FIGURES (docs/FIGURE_REDESIGN_SPEC.md S5, items 15
%   and 16, rule H1).  \hbarrow / \hbarrowr / \baxis print an accuracy number
%   inside the figure, which the v6 redraw forbids: fig1_framework, figm_bpr
%   and figm_fusion no longer call them and their numbers live in Tables IV
%   and VI.  The macros are KEPT, not deleted, because the parked single-
%   column data panel of Fig. 5 (figm_rerank, permutation importance) is a
%   data figure, where a measured bar length is evidence rather than a
%   result.  Same status, same reason: warn/, badge/ + badgeops/ + badgetrn/,
%   \hull and \mplane -- no v6 method figure uses any of them.
% =========================================================================
\newcommand{\hbarrow}[7]{%
  \draw[black!15, line width=0.5pt] (#1,#2) -- ({#1+#3},#2);
  \draw[#5, line width=#6] (#1,#2) -- ({#1+(#3)*(#4)},#2);
  % the numeral gets a paper backing so the empty part of the track does not
  % run into it and read as a stray rule after the digits
  \node[qty, anchor=west, inner sep=1.6pt, fill=white]
    at ({#1+(#3)*(#4)},#2) {#7};}

% \hbarrowr{x0}{y}{track}{frac}{colour}{weight}{label}
%   As \hbarrow, but the numeral is set in a fixed column just past the end
%   of the track instead of at the bar tip.  Use it where several groups sit
%   side by side on a SHORT track: it aligns the numerals into columns across
%   the groups, and it stops the leftover empty track of a short bar from
%   printing as a stray rule immediately after the digits -- at 1.25 cm of
%   track a 27.7/75 bar leaves 3 mm of rule past the numeral, which reads as
%   an en dash and can be mistaken for a range.
\newcommand{\hbarrowr}[7]{%
  \draw[black!15, line width=0.5pt] (#1,#2) -- ({#1+#3},#2);
  \draw[#5, line width=#6] (#1,#2) -- ({#1+(#3)*(#4)},#2);
  \node[qty, anchor=west, inner sep=1.6pt] at ({#1+(#3)+0.04},#2) {#7};}

% \baxis{x0}{y}{track}{lo}{hi}{unit} : the tick pair + caption for a \hbarrow
%   group, drawn once per group rather than once per bar.
\newcommand{\baxis}[6]{%
  \draw[black!45, line width=0.5pt] (#1,#2) -- (#1,{#2-0.05})
                                    ({#1+#3},#2) -- ({#1+#3},{#2-0.05});
  \node[sub, anchor=north, inner sep=1pt] at (#1,{#2-0.05}) {#4};
  \node[sub, anchor=north, inner sep=1pt] at ({#1+#3},{#2-0.05}) {#5};
  \node[sub, anchor=west, inner sep=1pt] at ({#1+#3+0.14},{#2-0.16}) {#6};}

% =========================================================================
% GEOMETRIC ISOMORPHISM (v3).  Tensor shape is DRAWN, never only written.
%
%  G1 SPATIAL EXTENT -> DRAWN SIZE, one scale for the whole figure set:
%       1 mm of drawn edge per 40 pixels of tensor extent,
%       i.e. drawn side [cm] = side [px] / 400   ->  \pxc{800} = 2.00 cm.
%       800 -> 20.0 mm | 200 -> 5.00 mm | 100 -> 2.50 mm |  50 -> 1.25 mm
%        25 ->  0.63 mm | 640 -> 16.0 mm | 960 -> 24.0 mm |1088 -> 27.2 mm
%        64 ->  1.60 mm | 128 ->  3.20 mm
%       A stride-4 map is therefore exactly 8x the linear size of a
%       stride-32 map, as it must be.
%  G2 CHANNEL COUNT -> SLAB DEPTH (isometric extrusion, 30 degrees):
%       depth [cm] = 0.04 + 0.055 * log2(C/256)
%       256 -> 0.40 mm | 512 -> 0.95 mm | 1024 -> 1.50 mm | 2048 -> 2.05 mm
%       i.e. +0.55 mm of depth per channel doubling.
%  G3 CATEGORY -> GLYPH SHAPE (readable in grayscale, no label needed):
%       spatial feature map (B,C,H,W) ......... extruded slab   \fmslab
%       token set (N x C, no spatial extent) .. column of bars  \tokencol
%       mask stack (N x H x W) ................ offset planes   \maskstack
%       scalar score .......................... short bar       \scorebar
%       feature vector (d cells) .............. strip of cells  (drawn inline)
%  G4 A LABELLED ZOOM INSET may magnify a sub-scene by an explicit factor z;
%       inside it G1/G2 hold multiplied by z.  Every macro takes z as its
%       last argument (z=1 for the global scale).
% =========================================================================
\newcommand{\pxc}[1]{(#1)/400}     % tensor side in px -> drawn side in cm

% =========================================================================
% G5 REAL DATA (v3).  Wherever a glyph stands for a tensor we actually own,
%     the glyph's FRONT FACE carries that tensor's real raster, dumped by
%     carrier/dump_figure_assets.py with provenance in assets/manifest.json.
%     The vector chrome (frames, arrows, labels, extrusion) stays vector, so
%     type remains selectable and the fonts stay embedded -- the bitmap-panel
%     + vector-chrome architecture Nature's artwork guide prescribes.
%     Resolution guard: every embedded raster must clear 300 dpi at its drawn
%     size and we clear 600 dpi everywhere (chip.png 800 px at 20 mm = 1016
%     dpi).  Feature maps are embedded from the nearest-neighbour _xN variants
%     so the PDF renderer cannot smooth pixels the tensor does not have.
%
% \realpanel{x}{y}{sidepx}{zoom}{file}{framespec} : a real raster, framed and
%     sized by G1 exactly as a schematic glyph of the same tensor would be.
% =========================================================================
\newcommand{\realpanel}[6]{%
  \begin{scope}[shift={(#1,#2)}]
    \pgfmathsetmacro{\rpH}{0.5*#4*\pxc{#3}}
    \pgfmathsetmacro{\rpS}{2*\rpH}
    \node[inner sep=0pt, outer sep=0pt] at (0,0)
      {\includegraphics[width=\rpS cm, height=\rpS cm]{\assetdir/#5}};
    \draw[#6] (-\rpH,-\rpH) rectangle (\rpH,\rpH);
  \end{scope}}

% \realpanelh{...} : the same panel HORIZONTALLY MIRRORED.  Used by the TTA
%     figure so the hflip half of each view pair is a real flip of the real
%     raster, not a redrawn mirror of a schematic.
\newcommand{\realpanelh}[6]{%
  \begin{scope}[shift={(#1,#2)}]
    \pgfmathsetmacro{\rpH}{0.5*#4*\pxc{#3}}
    \pgfmathsetmacro{\rpS}{2*\rpH}
    \node[inner sep=0pt, outer sep=0pt] at (0,0)
      {\reflectbox{\includegraphics[width=\rpS cm, height=\rpS cm]{\assetdir/#5}}};
    \draw[#6] (-\rpH,-\rpH) rectangle (\rpH,\rpH);
  \end{scope}}

% --- \fmslab{x}{y}{sidepx}{C}{colour}{zoom} : a spatial feature map -------
\newcommand{\fmslab}[6]{\fmslabr{#1}{#2}{#3}{#4}{#5}{#6}{}}

% --- \fmslabr{...}{file} : the same slab with a REAL raster on the front
%     face.  Pass an empty 7th argument for the flat-tinted variant.  The
%     extruded top/right faces stay flat tinted, so G2 (depth = channel
%     count) and the single light source (top brightest, right middle,
%     front darkest) are preserved.
\newcommand{\fmslabr}[7]{%
  \begin{scope}[shift={(#1,#2)}]
    \pgfmathsetmacro{\sbH}{0.5*#6*\pxc{#3}}
    \pgfmathsetmacro{\sbD}{#6*(0.04+0.055*ln((#4)/256)/ln(2))}
    \pgfmathsetmacro{\sbX}{0.866*\sbD}
    \pgfmathsetmacro{\sbY}{0.500*\sbD}
    \pgfmathsetmacro{\sbW}{2*\sbH}
    \fill[#5!30] (-\sbH,\sbH) -- (-\sbH+\sbX,\sbH+\sbY)
                 -- (\sbH+\sbX,\sbH+\sbY) -- (\sbH,\sbH) -- cycle;
    \fill[#5!20] (\sbH,\sbH) -- (\sbH+\sbX,\sbH+\sbY)
                 -- (\sbH+\sbX,-\sbH+\sbY) -- (\sbH,-\sbH) -- cycle;
    \fill[#5!11] (-\sbH,-\sbH) rectangle (\sbH,\sbH);
    \edef\fmfile{#7}%
    \ifx\fmfile\empty\else
      \node[inner sep=0pt, outer sep=0pt] at (0,0)
        {\includegraphics[width=\sbW cm, height=\sbW cm]{\assetdir/#7}};
    \fi
    \draw[#5!75, line width=0.5pt]
      (-\sbH,-\sbH) rectangle (\sbH,\sbH)
      (-\sbH,\sbH) -- (-\sbH+\sbX,\sbH+\sbY) -- (\sbH+\sbX,\sbH+\sbY) -- (\sbH,\sbH)
      (\sbH+\sbX,\sbH+\sbY) -- (\sbH+\sbX,-\sbH+\sbY) -- (\sbH,-\sbH);
  \end{scope}}

% --- \tokencol{x}{y}{nbars}{width}{colour} : an N x C token set -----------
%     Deliberately NOT a slab: flat bars, visible gaps, no extrusion, so a
%     query set can never be mistaken for a feature map.
\newcommand{\tokencol}[5]{%
  \begin{scope}[shift={(#1,#2)}]
    \pgfmathsetmacro{\tkP}{0.098}
    \pgfmathsetmacro{\tkT}{0.056}
    \pgfmathsetmacro{\tkW}{0.5*#4}
    \pgfmathtruncatemacro{\tkN}{#3-1}
    \pgfmathsetmacro{\tkY}{0.5*(\tkN*\tkP+\tkT)}
    \foreach \tki in {0,...,\tkN} {
      \pgfmathsetmacro{\tkc}{\tkY-\tki*\tkP}
      \fill[#5!13] (-\tkW,\tkc-\tkT) rectangle (\tkW,\tkc);
      \draw[#5!75, line width=0.5pt] (-\tkW,\tkc-\tkT) rectangle (\tkW,\tkc);
    }
  \end{scope}}

% --- \maskstack{x}{y}{sidepx}{k}{colour}{zoom} : k planes of H x W --------
\newcommand{\maskstack}[6]{\maskstackr{#1}{#2}{#3}{#4}{#5}{#6}{}}

% --- \maskstackr{...}{file} : the same pile with the FRONT plane carrying a
%     real raster; the planes behind stay schematic so the "stack of N"
%     reading survives.  Empty 7th argument = the all-schematic variant.
\newcommand{\maskstackr}[7]{%
  \begin{scope}[shift={(#1,#2)}]
    \pgfmathsetmacro{\msH}{0.5*#6*\pxc{#3}}
    \pgfmathsetmacro{\msO}{0.30*\msH/2.5}
    \pgfmathtruncatemacro{\msN}{#4-1}
    \pgfmathsetmacro{\msS}{\msH/0.92}
    \pgfmathsetmacro{\msW}{2*\msH}
    \edef\msfile{#7}%
    \foreach \msi in {\msN,...,0} {
      \pgfmathsetmacro{\msd}{\msi*\msO}
      \pgfmathtruncatemacro{\msreal}{(\msi==0) ? 1 : 0}
      \edef\mshasimg{0}%
      \ifx\msfile\empty\else\ifnum\msreal=1 \edef\mshasimg{1}\fi\fi
      \begin{scope}[shift={(\msd,\msd)}]
        \fill[white] (-\msH,-\msH) rectangle (\msH,\msH);
        \fill[#5!8]  (-\msH,-\msH) rectangle (\msH,\msH);
        \ifnum\mshasimg=1
          \node[inner sep=0pt, outer sep=0pt] at (0,0)
            {\includegraphics[width=\msW cm, height=\msW cm]{\assetdir/#7}};
        \fi
        \draw[#5!60, line width=0.5pt] (-\msH,-\msH) rectangle (\msH,\msH);
        \ifnum\mshasimg=0
          \begin{scope}[scale=\msS]
            \fill[#5!35] \hull;
            \draw[#5!90, line width=0.5pt] \hull;
          \end{scope}
        \fi
      \end{scope}}
  \end{scope}}

% --- \scorebar{x}{y}{h}{colour} : one scalar, as a short upright bar ------
\newcommand{\scorebar}[4]{%
  \begin{scope}[shift={(#1,#2)}]
    \draw[black!35, line width=0.5pt] (-0.10,0) rectangle (0.10,0.62);
    \fill[#4!55] (-0.10,0) rectangle (0.10,#3);
    \draw[#4!90, line width=0.5pt] (-0.10,0) rectangle (0.10,#3);
  \end{scope}}

% =========================================================================
% Reusable glyphs.  All are TikZ pics: \pic[scale=s] at (x,y) {name};
% =========================================================================

% --- lock: marks any component whose weights we never update -------------
\tikzset{lock/.pic={
  \fill[frozgray] (-0.105,-0.105) rectangle (0.105,0.055);
  \draw[frozgray, line width=0.65pt]
    (-0.060,0.055) -- (-0.060,0.105)
    arc[start angle=180, end angle=0, radius=0.060] -- (0.060,0.055);
  \fill[white] (0,-0.022) circle[radius=0.020];
}}

% --- planes: 3 offset stacked planes = a multi-scale feature pyramid -----
\tikzset{planes/.pic={
  \foreach \i/\o in {2/0.20, 1/0.10, 0/0.00} {
    \fill[datablue!\the\numexpr 10+8*\i\relax]
      (\o-0.24,\o-0.30) rectangle (\o+0.24,\o+0.30);
    \draw[datablue!70, line width=0.5pt]
      (\o-0.24,\o-0.30) rectangle (\o+0.24,\o+0.30);
  }
}}

% --- plate: a single feature plane (thin vertical plate, O4) -------------
\tikzset{plate/.pic={
  \fill[datablue!14] (-0.10,-0.34) rectangle (0.10,0.34);
  \draw[datablue!70, line width=0.5pt] (-0.10,-0.34) rectangle (0.10,0.34);
  \draw[datablue!45, line width=0.5pt] (-0.10,0.20) -- (0.10,0.20);
  \draw[datablue!45, line width=0.5pt] (-0.10,-0.20) -- (0.10,-0.20);
}}

% --- ckpt: an isometric slab = one stored checkpoint / weight vector -----
\tikzset{ckpt/.pic={
  \fill[frozgray!18] (-0.30,-0.13) rectangle (0.24,0.13);
  \fill[frozgray!32] (-0.30,0.13) -- (-0.24,0.21) -- (0.30,0.21) -- (0.24,0.13) -- cycle;
  \fill[frozgray!26] (0.24,-0.13) -- (0.30,-0.05) -- (0.30,0.21) -- (0.24,0.13) -- cycle;
  \draw[frozgray, line width=0.5pt]
    (-0.30,-0.13) rectangle (0.24,0.13);
  \draw[frozgray, line width=0.5pt]
    (-0.30,0.13) -- (-0.24,0.21) -- (0.30,0.21) -- (0.24,0.13)
    (0.30,0.21) -- (0.30,-0.05) -- (0.24,-0.13);
}}

% --- masked: pictogram of an output instance mask (hull on light field) --
\tikzset{maskicon/.pic={
  \fill[datablue!5] (-0.30,-0.30) rectangle (0.30,0.30);
  \draw[datablue!60, line width=0.5pt] (-0.30,-0.30) rectangle (0.30,0.30);
  \begin{scope}[rotate=28]
    \fill[datablue!45] (0,0.05) ellipse[x radius=0.235, y radius=0.062];
    \draw[datablue, line width=0.5pt] (0,0.05) ellipse[x radius=0.235, y radius=0.062];
  \end{scope}
  \begin{scope}[rotate=-12]
    \fill[datablue!45] (0.02,-0.19) ellipse[x radius=0.145, y radius=0.045];
    \draw[datablue, line width=0.5pt] (0.02,-0.19) ellipse[x radius=0.145, y radius=0.045];
  \end{scope}
}}

% =========================================================================
% \sarchip was REMOVED in v3.  It drew two smooth ellipses each ~2.6% of the
% frame on a bokeh-like field; the real HRSID inshore median instance covers
% 0.028% of the chip (13.4 px on 800x800), so the glyph was ~90x too large and
% contradicted the small-target-in-clutter argument the paper makes.  Every
% call site now embeds the REAL chip raster instead (assets/chip.png,
% assets/view_{640,800,960,1088}.png; see the realpanel macro above and
% assets/manifest.json).
% =========================================================================

% =========================================================================
% \hull{}  -- schematic ship-mask contour, ~1.3 cm wide, centred at origin.
% Used by the BPR and fusion figures so every ship outline in the paper is
% literally the same curve.
% =========================================================================
\newcommand{\hull}{plot[smooth cycle, tension=0.72] coordinates
  {(0.72,0.00) (0.34,0.26) (-0.22,0.33) (-0.64,0.17) (-0.62,-0.12)
   (-0.10,-0.30) (0.42,-0.22)}}
.  Nothing else in the figure has to change;
% the block/arrow style names are the same, only their numbers move.
% =========================================================================
% =========================================================================
% tikz_style.tex — shared visual system for the Section III method figures.
% Usage: \input{tikz_style} in the preamble of a standalone figure.
%
% -------------------------------------------------------------------------
% DESIGN NOTES — conventions extracted from an accepted TGRS manuscript
% (Gao et al., "BBox-free ... SAR ship instance segmentation", TGRS 2024;
%  local copy: scratchpad/gao_bboxfree_tgrs24.pdf).  Observations taken from
%  its Fig. 2 (overall architecture, p. 5), Fig. 5 (FD module, p. 8) and
%  Fig. 6 (conceptual mask panel, p. 8):
%
%  O1  REGION PARTITION.  The architecture figure is carved into 3-4 named
%      regions by DOTTED rectangles with no fill; the region name is set in
%      bold OUTSIDE the rectangle ("Feature Extractor", "Classification
%      Branch", "Localization Branch", "Center Point Extractor").  Nothing
%      is nested more than two levels deep.
%  O2  COLOR CARRIES MEANING, NOT DECORATION.  Only ~5 hues appear, all as
%      pale fills with black text: light gray = ordinary/standard module,
%      salmon-orange = the module the paper contributes (Feature Decoupling
%      Module, Center Point Extractor), yellow = a distinguished operation
%      (Feature Fusion), light blue / light pink = tensors.  Saturated color
%      is reserved for the few glyphs that must be found first.
%  O3  LEGEND PANEL.  Both Fig. 2 and Fig. 5 carry an explicit legend inside
%      the figure (dotted box top-left in Fig. 2, dashed box right in Fig. 5)
%      that decodes arrow colors and single-letter glyph nodes.  The body of
%      the diagram is kept uncluttered by pushing vocabulary into the legend.
%  O4  TENSOR ICONOGRAPHY.  Feature maps are drawn as thin vertical plates,
%      multi-scale features as 3 offset stacked planes, channel vectors as
%      segmented 3-D bars, the backbone as an isometric slab.  Real SAR
%      chips / heat maps / GT overlays are pasted directly into the flow.
%  O5  TENSOR NAMES SIT BESIDE THE GLYPH, NEVER INSIDE IT, in math italic
%      (F_out, F_mask, F_param, F_H_zip ...), placed above-right of the plate.
%  O6  ARROW DISCIPLINE.  Data flow is a thin gray/black line with a small
%      solid triangular head; routing is orthogonal (Manhattan) with a shared
%      trunk when one source fans out to several targets, so lines never
%      cross.  Semantically special transitions get a THICK COLORED BLOCK
%      arrow instead (2x upsample / 2x downsample / connect), i.e. arrow
%      style is itself a semantic channel.  Long skip paths are routed as
%      rails around the outside of the diagram (Fig. 5 top and bottom).
%  O7  TYPOGRAPHY.  One size for block labels, one smaller size for
%      sub-labels, math only in math mode; no stretched or rotated text
%      except a rare 90-degree axis label.
%  O8  ASPECT RATIO.  Double-column method figures are wide and short:
%      Fig. 2 is 530 x 200 pt (2.65:1), Fig. 5 is 530 x 215 pt (2.5:1).
%      They are laid on a grid with equal gutters between parallel branches.
%  O9  CAPTION DIVISION OF LABOR.  The caption names the parts and decodes
%      abbreviations/colors; it does not re-derive the mechanism, which is
%      left to the body text.  The figure is readable without the caption.
%  O10 CONCEPTUAL PANELS (Fig. 6) use dashed sub-panels with a centered
%      label underneath, a pale blue blob for the true object and a pale red
%      overlay for the error region -- schematic, not photorealistic.
%
% -------------------------------------------------------------------------
% WHAT WE ADOPT / WHERE WE DEVIATE
%  + adopt O1 (dotted named regions), O3 (in-figure legend), O4/O5 (tensor
%    glyphs with the symbol outside), O6 (one flow style + Manhattan
%    routing), O8 (wide-short double-column ratios), O10 (schematic panels).
%  ~ deviate on O2: instead of hue-as-identity we bind hue to ROLE in the
%    contribution -- frozen / training-free operator / trained corrector /
%    data -- because the entire argument of this article is which parts are
%    touched and which are not.
%  ~ deviate on the sans-serif block labels of Fig. 2: we set every figure in
%    the body Times face (as the same paper's Fig. 6 does), so figure text
%    and caption text share one face.
%  + every figure is authored at its EXACT printed width (18.2 cm double
%    column = 516 pt, 8.89 cm single column = 252 pt) and included at
%    scale 1, so \footnotesize is a true 8 pt and \scriptsize a true 7 pt on
%    the page in ALL figures.  Use \figcanvasd / \figcanvasc below.
%
% -------------------------------------------------------------------------
% SEMANTIC PALETTE (Okabe-Ito colorblind-safe base; redundant encoding by
% border weight, dash pattern and badge glyph so it survives grayscale):
%   froz  gray   + lock glyph        frozen, not modified by us  (Stage A)
%   ops   vermil + solid heavy rule  our TRAINING-FREE operator  (rung 1)
%   trn   green  + dashed rule       our OPTIONAL TRAINED corrector (2-3)
%   data  blue   + thin rule         data, tensors, I/O
%   warn  red    + dotted rule, small  failure mode / negative control
% =========================================================================

% Real rasters are embedded on the front faces of the tensor glyphs (see
% \realpanel / \fmslab below).  graphicx must be loaded in the preamble, which
% is where this file is \input.
\usepackage{graphicx}
% Single switch for the raster payload: \def\assetdir{assets_light}
% before \begin{document} swaps every embedded raster for its
% print-sized ~600 dpi copy (see figs/make_light_assets.py).  The
% default is the native dump, 965-1016 dpi at the drawn size.
\providecommand{\assetdir}{assets}
\usetikzlibrary{positioning,arrows.meta,fit,backgrounds,calc,shapes.geometric,
                decorations.pathreplacing,decorations.markings,patterns}

% ---------------- palette ----------------
\definecolor{opsred}{HTML}{D55E00}   % Okabe-Ito vermillion  -- training-free ops
\definecolor{trngreen}{HTML}{009E73} % Okabe-Ito bluish green-- trained correctors
\definecolor{datablue}{HTML}{0072B2} % Okabe-Ito blue        -- data / tensors
\definecolor{frozgray}{HTML}{6E6E6E} % neutral               -- frozen base
\definecolor{flowink}{HTML}{262626}  % near black            -- data-flow arrows
\definecolor{warnred}{HTML}{9E2B2B}  % muted brick           -- failure modes
\definecolor{sarink}{HTML}{2B2B33}   % SAR chip background

% ---------------- exact-width canvases ----------------
% \figcanvasd{H} : double column, 18.2 cm = 516 pt = \textwidth  (scale 1.0)
% \figcanvasc{H} : single column,  8.89 cm = 252 pt = \columnwidth
\newcommand{\figcanvasd}[1]{\useasboundingbox (0,0) rectangle (18.2,#1);}
\newcommand{\figcanvasc}[1]{\useasboundingbox (0,0) rectangle (8.89,#1);}

\tikzset{
  % ---- picture-level defaults: ONE font size system for all five figures ----
  fig/.style={font=\footnotesize, line width=0.65pt,
              >={Stealth[length=1.7mm,width=1.35mm]},
              every node/.append style={align=center, inner sep=2.4pt}},
  % ---- blocks: role -> color + redundant border encoding ----
  % Fill tints were raised in v3: at 4-8% the five roles were separated by an
  % RGB distance of only 6.6-21 (out of 441) and all landed on gray 243-248,
  % i.e. the fill channel was consuming ink and returning nothing.  The border
  % weight / dash pattern / badge remain the PRIMARY, grayscale-safe channel.
  ops/.style   ={draw=opsred,   fill=opsred!16,   line width=0.9pt,
                 rounded corners=1.6pt, inner sep=3pt, align=center},
  trn/.style   ={draw=trngreen, fill=trngreen!15, line width=0.9pt,
                 dash pattern=on 1.7pt off 1.1pt,
                 rounded corners=1.6pt, inner sep=3pt, align=center},
  froz/.style  ={draw=frozgray, fill=black!8,     line width=0.7pt,
                 rounded corners=1.6pt, inner sep=3pt, align=center},
  data/.style  ={draw=datablue, fill=datablue!14, line width=0.55pt,
                 rounded corners=1.6pt, inner sep=3pt, align=center},
  plain/.style ={draw=black!45, fill=white,       line width=0.55pt,
                 rounded corners=1.6pt, inner sep=3pt, align=center},
  % v4: radius 1.4 -> 1.6 pt, the one block style that did not share the
  % system's corner radius; weight 0.5 -> 0.6 pt so the failure chips are not
  % the faintest object on the page while being the ones a reader must find.
  % The dotted dash pattern stays as the grayscale-safe channel.
  warn/.style  ={draw=warnred,  fill=warnred!9,   line width=0.6pt,
                 dash pattern=on 1.1pt off 1.1pt,
                 rounded corners=1.6pt, inner sep=3pt, align=center},
  gate/.style  ={draw=opsred,   fill=opsred!16,   line width=0.9pt,
                 diamond, aspect=2.15, inner sep=1pt, align=center},
  % ---- named region (O1): dotted, unfilled, label set outside ----
  grp/.style   ={draw=black!45, densely dotted, line width=0.7pt,
                 rounded corners=2pt, inner sep=0pt},
  grptitle/.style={font=\footnotesize\bfseries, text=black!80, inner sep=1.5pt},
  % ---- step badges (redundant, grayscale-safe stage identity) ----
  badge/.style ={circle, inner sep=0pt, minimum size=3.3mm, text=white,
                 font=\scriptsize\bfseries},
  badgeops/.style={badge, fill=opsred},
  badgetrn/.style={badge, fill=trngreen},
  % ---- typography (O7): exactly three text sizes in the whole figure set ----
  % v3: 7 pt at 55% black was the worst corner of the design space (minimum
  % size AND minimum contrast) -- sub-labels are now 7 pt at 70% black, and
  % anything carrying a QUANTITY is promoted to 8 pt via \qty / \subk.
  ttl/.style   ={font=\footnotesize\bfseries},
  lbl/.style   ={font=\footnotesize},
  sub/.style   ={font=\scriptsize, text=black!70},
  subk/.style  ={font=\footnotesize, text=black!80},
  qty/.style   ={font=\footnotesize, text=black!80},% tensor dims, AP numbers
  tsr/.style ={font=\scriptsize\ttfamily, text=black!70},% B x C x H x W (O5)
  tag/.style   ={font=\scriptsize, inner sep=1.6pt, rounded corners=1pt},
  % panel letters (B11): \footnotesize\bfseries at the panel's top-left
  panl/.style  ={font=\footnotesize\bfseries, text=black!85, inner sep=1pt},
  % ---- arrows (O6): one data-flow style, one auxiliary/optional style ----
  flow/.style  ={-{Stealth[length=1.7mm,width=1.35mm]}, draw=flowink,
                 line width=0.65pt},
  flowt/.style ={-{Stealth[length=1.7mm,width=1.35mm]}, draw=flowink,
                 line width=0.65pt, rounded corners=1.6pt},
  aux/.style   ={-{Stealth[length=1.5mm,width=1.2mm]}, draw=black!45,
                 line width=0.5pt, densely dashed},
  auxw/.style  ={-{Stealth[length=1.5mm,width=1.2mm]}, draw=warnred!75,
                 line width=0.5pt, densely dashed},
  rail/.style  ={draw=black!30, line width=0.5pt},
  brace/.style ={decorate, decoration={brace, amplitude=3.2pt, raise=1pt},
                 draw=black!50, line width=0.5pt},
  % ---- ONE leader style for every magnifier cone in the figure set --------
  % v4: the zoom cones used to be drawn in the generic dashed `aux' style, so
  % a sight line was typographically the same object as an ordinary auxiliary
  % arrow.  A cone is now its own object: it carries the hue of the locator
  % rectangle it starts from, no arrowhead (it is a correspondence, not a
  % flow), and it always runs corner-to-corner from the locator box to the
  % MAGNIFIED RASTER, never to the enclosing panel frame.
  zoomray/.style={draw=opsred!55, line width=0.5pt,
                  dash pattern=on 1.3pt off 1.1pt},
  zoombox/.style={draw=opsred, line width=0.7pt},
}

% =========================================================================
% \hbarrow{x0}{y}{track}{frac}{colour}{weight}{label}
%   One horizontal bar on a printed 0..1 track, with the numeral set at the
%   bar end.  v4: fig1's AP progression, figm_tta (c), figm_fusion's result
%   panel and figm_rerank (b) all drew this same object with four slightly
%   different geometries.  They now share one macro, so the reader learns the
%   idiom once and meets it four times -- a small multiple across figures.
%   The empty track is always drawn, so a short bar still shows its scale.
%
%   DEPRECATED FOR METHOD FIGURES (docs/FIGURE_REDESIGN_SPEC.md S5, items 15
%   and 16, rule H1).  \hbarrow / \hbarrowr / \baxis print an accuracy number
%   inside the figure, which the v6 redraw forbids: fig1_framework, figm_bpr
%   and figm_fusion no longer call them and their numbers live in Tables IV
%   and VI.  The macros are KEPT, not deleted, because the parked single-
%   column data panel of Fig. 5 (figm_rerank, permutation importance) is a
%   data figure, where a measured bar length is evidence rather than a
%   result.  Same status, same reason: warn/, badge/ + badgeops/ + badgetrn/,
%   \hull and \mplane -- no v6 method figure uses any of them.
% =========================================================================
\newcommand{\hbarrow}[7]{%
  \draw[black!15, line width=0.5pt] (#1,#2) -- ({#1+#3},#2);
  \draw[#5, line width=#6] (#1,#2) -- ({#1+(#3)*(#4)},#2);
  % the numeral gets a paper backing so the empty part of the track does not
  % run into it and read as a stray rule after the digits
  \node[qty, anchor=west, inner sep=1.6pt, fill=white]
    at ({#1+(#3)*(#4)},#2) {#7};}

% \hbarrowr{x0}{y}{track}{frac}{colour}{weight}{label}
%   As \hbarrow, but the numeral is set in a fixed column just past the end
%   of the track instead of at the bar tip.  Use it where several groups sit
%   side by side on a SHORT track: it aligns the numerals into columns across
%   the groups, and it stops the leftover empty track of a short bar from
%   printing as a stray rule immediately after the digits -- at 1.25 cm of
%   track a 27.7/75 bar leaves 3 mm of rule past the numeral, which reads as
%   an en dash and can be mistaken for a range.
\newcommand{\hbarrowr}[7]{%
  \draw[black!15, line width=0.5pt] (#1,#2) -- ({#1+#3},#2);
  \draw[#5, line width=#6] (#1,#2) -- ({#1+(#3)*(#4)},#2);
  \node[qty, anchor=west, inner sep=1.6pt] at ({#1+(#3)+0.04},#2) {#7};}

% \baxis{x0}{y}{track}{lo}{hi}{unit} : the tick pair + caption for a \hbarrow
%   group, drawn once per group rather than once per bar.
\newcommand{\baxis}[6]{%
  \draw[black!45, line width=0.5pt] (#1,#2) -- (#1,{#2-0.05})
                                    ({#1+#3},#2) -- ({#1+#3},{#2-0.05});
  \node[sub, anchor=north, inner sep=1pt] at (#1,{#2-0.05}) {#4};
  \node[sub, anchor=north, inner sep=1pt] at ({#1+#3},{#2-0.05}) {#5};
  \node[sub, anchor=west, inner sep=1pt] at ({#1+#3+0.14},{#2-0.16}) {#6};}

% =========================================================================
% GEOMETRIC ISOMORPHISM (v3).  Tensor shape is DRAWN, never only written.
%
%  G1 SPATIAL EXTENT -> DRAWN SIZE, one scale for the whole figure set:
%       1 mm of drawn edge per 40 pixels of tensor extent,
%       i.e. drawn side [cm] = side [px] / 400   ->  \pxc{800} = 2.00 cm.
%       800 -> 20.0 mm | 200 -> 5.00 mm | 100 -> 2.50 mm |  50 -> 1.25 mm
%        25 ->  0.63 mm | 640 -> 16.0 mm | 960 -> 24.0 mm |1088 -> 27.2 mm
%        64 ->  1.60 mm | 128 ->  3.20 mm
%       A stride-4 map is therefore exactly 8x the linear size of a
%       stride-32 map, as it must be.
%  G2 CHANNEL COUNT -> SLAB DEPTH (isometric extrusion, 30 degrees):
%       depth [cm] = 0.04 + 0.055 * log2(C/256)
%       256 -> 0.40 mm | 512 -> 0.95 mm | 1024 -> 1.50 mm | 2048 -> 2.05 mm
%       i.e. +0.55 mm of depth per channel doubling.
%  G3 CATEGORY -> GLYPH SHAPE (readable in grayscale, no label needed):
%       spatial feature map (B,C,H,W) ......... extruded slab   \fmslab
%       token set (N x C, no spatial extent) .. column of bars  \tokencol
%       mask stack (N x H x W) ................ offset planes   \maskstack
%       scalar score .......................... short bar       \scorebar
%       feature vector (d cells) .............. strip of cells  (drawn inline)
%  G4 A LABELLED ZOOM INSET may magnify a sub-scene by an explicit factor z;
%       inside it G1/G2 hold multiplied by z.  Every macro takes z as its
%       last argument (z=1 for the global scale).
% =========================================================================
\newcommand{\pxc}[1]{(#1)/400}     % tensor side in px -> drawn side in cm

% =========================================================================
% G5 REAL DATA (v3).  Wherever a glyph stands for a tensor we actually own,
%     the glyph's FRONT FACE carries that tensor's real raster, dumped by
%     carrier/dump_figure_assets.py with provenance in assets/manifest.json.
%     The vector chrome (frames, arrows, labels, extrusion) stays vector, so
%     type remains selectable and the fonts stay embedded -- the bitmap-panel
%     + vector-chrome architecture Nature's artwork guide prescribes.
%     Resolution guard: every embedded raster must clear 300 dpi at its drawn
%     size and we clear 600 dpi everywhere (chip.png 800 px at 20 mm = 1016
%     dpi).  Feature maps are embedded from the nearest-neighbour _xN variants
%     so the PDF renderer cannot smooth pixels the tensor does not have.
%
% \realpanel{x}{y}{sidepx}{zoom}{file}{framespec} : a real raster, framed and
%     sized by G1 exactly as a schematic glyph of the same tensor would be.
% =========================================================================
\newcommand{\realpanel}[6]{%
  \begin{scope}[shift={(#1,#2)}]
    \pgfmathsetmacro{\rpH}{0.5*#4*\pxc{#3}}
    \pgfmathsetmacro{\rpS}{2*\rpH}
    \node[inner sep=0pt, outer sep=0pt] at (0,0)
      {\includegraphics[width=\rpS cm, height=\rpS cm]{\assetdir/#5}};
    \draw[#6] (-\rpH,-\rpH) rectangle (\rpH,\rpH);
  \end{scope}}

% \realpanelh{...} : the same panel HORIZONTALLY MIRRORED.  Used by the TTA
%     figure so the hflip half of each view pair is a real flip of the real
%     raster, not a redrawn mirror of a schematic.
\newcommand{\realpanelh}[6]{%
  \begin{scope}[shift={(#1,#2)}]
    \pgfmathsetmacro{\rpH}{0.5*#4*\pxc{#3}}
    \pgfmathsetmacro{\rpS}{2*\rpH}
    \node[inner sep=0pt, outer sep=0pt] at (0,0)
      {\reflectbox{\includegraphics[width=\rpS cm, height=\rpS cm]{\assetdir/#5}}};
    \draw[#6] (-\rpH,-\rpH) rectangle (\rpH,\rpH);
  \end{scope}}

% --- \fmslab{x}{y}{sidepx}{C}{colour}{zoom} : a spatial feature map -------
\newcommand{\fmslab}[6]{\fmslabr{#1}{#2}{#3}{#4}{#5}{#6}{}}

% --- \fmslabr{...}{file} : the same slab with a REAL raster on the front
%     face.  Pass an empty 7th argument for the flat-tinted variant.  The
%     extruded top/right faces stay flat tinted, so G2 (depth = channel
%     count) and the single light source (top brightest, right middle,
%     front darkest) are preserved.
\newcommand{\fmslabr}[7]{%
  \begin{scope}[shift={(#1,#2)}]
    \pgfmathsetmacro{\sbH}{0.5*#6*\pxc{#3}}
    \pgfmathsetmacro{\sbD}{#6*(0.04+0.055*ln((#4)/256)/ln(2))}
    \pgfmathsetmacro{\sbX}{0.866*\sbD}
    \pgfmathsetmacro{\sbY}{0.500*\sbD}
    \pgfmathsetmacro{\sbW}{2*\sbH}
    \fill[#5!30] (-\sbH,\sbH) -- (-\sbH+\sbX,\sbH+\sbY)
                 -- (\sbH+\sbX,\sbH+\sbY) -- (\sbH,\sbH) -- cycle;
    \fill[#5!20] (\sbH,\sbH) -- (\sbH+\sbX,\sbH+\sbY)
                 -- (\sbH+\sbX,-\sbH+\sbY) -- (\sbH,-\sbH) -- cycle;
    \fill[#5!11] (-\sbH,-\sbH) rectangle (\sbH,\sbH);
    \edef\fmfile{#7}%
    \ifx\fmfile\empty\else
      \node[inner sep=0pt, outer sep=0pt] at (0,0)
        {\includegraphics[width=\sbW cm, height=\sbW cm]{\assetdir/#7}};
    \fi
    \draw[#5!75, line width=0.5pt]
      (-\sbH,-\sbH) rectangle (\sbH,\sbH)
      (-\sbH,\sbH) -- (-\sbH+\sbX,\sbH+\sbY) -- (\sbH+\sbX,\sbH+\sbY) -- (\sbH,\sbH)
      (\sbH+\sbX,\sbH+\sbY) -- (\sbH+\sbX,-\sbH+\sbY) -- (\sbH,-\sbH);
  \end{scope}}

% --- \tokencol{x}{y}{nbars}{width}{colour} : an N x C token set -----------
%     Deliberately NOT a slab: flat bars, visible gaps, no extrusion, so a
%     query set can never be mistaken for a feature map.
\newcommand{\tokencol}[5]{%
  \begin{scope}[shift={(#1,#2)}]
    \pgfmathsetmacro{\tkP}{0.098}
    \pgfmathsetmacro{\tkT}{0.056}
    \pgfmathsetmacro{\tkW}{0.5*#4}
    \pgfmathtruncatemacro{\tkN}{#3-1}
    \pgfmathsetmacro{\tkY}{0.5*(\tkN*\tkP+\tkT)}
    \foreach \tki in {0,...,\tkN} {
      \pgfmathsetmacro{\tkc}{\tkY-\tki*\tkP}
      \fill[#5!13] (-\tkW,\tkc-\tkT) rectangle (\tkW,\tkc);
      \draw[#5!75, line width=0.5pt] (-\tkW,\tkc-\tkT) rectangle (\tkW,\tkc);
    }
  \end{scope}}

% --- \maskstack{x}{y}{sidepx}{k}{colour}{zoom} : k planes of H x W --------
\newcommand{\maskstack}[6]{\maskstackr{#1}{#2}{#3}{#4}{#5}{#6}{}}

% --- \maskstackr{...}{file} : the same pile with the FRONT plane carrying a
%     real raster; the planes behind stay schematic so the "stack of N"
%     reading survives.  Empty 7th argument = the all-schematic variant.
\newcommand{\maskstackr}[7]{%
  \begin{scope}[shift={(#1,#2)}]
    \pgfmathsetmacro{\msH}{0.5*#6*\pxc{#3}}
    \pgfmathsetmacro{\msO}{0.30*\msH/2.5}
    \pgfmathtruncatemacro{\msN}{#4-1}
    \pgfmathsetmacro{\msS}{\msH/0.92}
    \pgfmathsetmacro{\msW}{2*\msH}
    \edef\msfile{#7}%
    \foreach \msi in {\msN,...,0} {
      \pgfmathsetmacro{\msd}{\msi*\msO}
      \pgfmathtruncatemacro{\msreal}{(\msi==0) ? 1 : 0}
      \edef\mshasimg{0}%
      \ifx\msfile\empty\else\ifnum\msreal=1 \edef\mshasimg{1}\fi\fi
      \begin{scope}[shift={(\msd,\msd)}]
        \fill[white] (-\msH,-\msH) rectangle (\msH,\msH);
        \fill[#5!8]  (-\msH,-\msH) rectangle (\msH,\msH);
        \ifnum\mshasimg=1
          \node[inner sep=0pt, outer sep=0pt] at (0,0)
            {\includegraphics[width=\msW cm, height=\msW cm]{\assetdir/#7}};
        \fi
        \draw[#5!60, line width=0.5pt] (-\msH,-\msH) rectangle (\msH,\msH);
        \ifnum\mshasimg=0
          \begin{scope}[scale=\msS]
            \fill[#5!35] \hull;
            \draw[#5!90, line width=0.5pt] \hull;
          \end{scope}
        \fi
      \end{scope}}
  \end{scope}}

% --- \scorebar{x}{y}{h}{colour} : one scalar, as a short upright bar ------
\newcommand{\scorebar}[4]{%
  \begin{scope}[shift={(#1,#2)}]
    \draw[black!35, line width=0.5pt] (-0.10,0) rectangle (0.10,0.62);
    \fill[#4!55] (-0.10,0) rectangle (0.10,#3);
    \draw[#4!90, line width=0.5pt] (-0.10,0) rectangle (0.10,#3);
  \end{scope}}

% =========================================================================
% Reusable glyphs.  All are TikZ pics: \pic[scale=s] at (x,y) {name};
% =========================================================================

% --- lock: marks any component whose weights we never update -------------
\tikzset{lock/.pic={
  \fill[frozgray] (-0.105,-0.105) rectangle (0.105,0.055);
  \draw[frozgray, line width=0.65pt]
    (-0.060,0.055) -- (-0.060,0.105)
    arc[start angle=180, end angle=0, radius=0.060] -- (0.060,0.055);
  \fill[white] (0,-0.022) circle[radius=0.020];
}}

% --- planes: 3 offset stacked planes = a multi-scale feature pyramid -----
\tikzset{planes/.pic={
  \foreach \i/\o in {2/0.20, 1/0.10, 0/0.00} {
    \fill[datablue!\the\numexpr 10+8*\i\relax]
      (\o-0.24,\o-0.30) rectangle (\o+0.24,\o+0.30);
    \draw[datablue!70, line width=0.5pt]
      (\o-0.24,\o-0.30) rectangle (\o+0.24,\o+0.30);
  }
}}

% --- plate: a single feature plane (thin vertical plate, O4) -------------
\tikzset{plate/.pic={
  \fill[datablue!14] (-0.10,-0.34) rectangle (0.10,0.34);
  \draw[datablue!70, line width=0.5pt] (-0.10,-0.34) rectangle (0.10,0.34);
  \draw[datablue!45, line width=0.5pt] (-0.10,0.20) -- (0.10,0.20);
  \draw[datablue!45, line width=0.5pt] (-0.10,-0.20) -- (0.10,-0.20);
}}

% --- ckpt: an isometric slab = one stored checkpoint / weight vector -----
\tikzset{ckpt/.pic={
  \fill[frozgray!18] (-0.30,-0.13) rectangle (0.24,0.13);
  \fill[frozgray!32] (-0.30,0.13) -- (-0.24,0.21) -- (0.30,0.21) -- (0.24,0.13) -- cycle;
  \fill[frozgray!26] (0.24,-0.13) -- (0.30,-0.05) -- (0.30,0.21) -- (0.24,0.13) -- cycle;
  \draw[frozgray, line width=0.5pt]
    (-0.30,-0.13) rectangle (0.24,0.13);
  \draw[frozgray, line width=0.5pt]
    (-0.30,0.13) -- (-0.24,0.21) -- (0.30,0.21) -- (0.24,0.13)
    (0.30,0.21) -- (0.30,-0.05) -- (0.24,-0.13);
}}

% --- masked: pictogram of an output instance mask (hull on light field) --
\tikzset{maskicon/.pic={
  \fill[datablue!5] (-0.30,-0.30) rectangle (0.30,0.30);
  \draw[datablue!60, line width=0.5pt] (-0.30,-0.30) rectangle (0.30,0.30);
  \begin{scope}[rotate=28]
    \fill[datablue!45] (0,0.05) ellipse[x radius=0.235, y radius=0.062];
    \draw[datablue, line width=0.5pt] (0,0.05) ellipse[x radius=0.235, y radius=0.062];
  \end{scope}
  \begin{scope}[rotate=-12]
    \fill[datablue!45] (0.02,-0.19) ellipse[x radius=0.145, y radius=0.045];
    \draw[datablue, line width=0.5pt] (0.02,-0.19) ellipse[x radius=0.145, y radius=0.045];
  \end{scope}
}}

% =========================================================================
% \sarchip was REMOVED in v3.  It drew two smooth ellipses each ~2.6% of the
% frame on a bokeh-like field; the real HRSID inshore median instance covers
% 0.028% of the chip (13.4 px on 800x800), so the glyph was ~90x too large and
% contradicted the small-target-in-clutter argument the paper makes.  Every
% call site now embeds the REAL chip raster instead (assets/chip.png,
% assets/view_{640,800,960,1088}.png; see the realpanel macro above and
% assets/manifest.json).
% =========================================================================

% =========================================================================
% \hull{}  -- schematic ship-mask contour, ~1.3 cm wide, centred at origin.
% Used by the BPR and fusion figures so every ship outline in the paper is
% literally the same curve.
% =========================================================================
\newcommand{\hull}{plot[smooth cycle, tension=0.72] coordinates
  {(0.72,0.00) (0.34,0.26) (-0.22,0.33) (-0.64,0.17) (-0.62,-0.12)
   (-0.10,-0.30) (0.42,-0.22)}}


% ---------------- T1: type ----------------
% Script sizes for the two text sizes used in figures.  Must be in the
% preamble; standalone + \input keeps it there.
\DeclareMathSizes{8}{8}{6}{6}
\DeclareMathSizes{7}{7}{6}{6}

% ---------------- T2/T3/T4: block, frame and arrow chrome ----------------
\tikzset{
  fig/.style={font=\footnotesize, line width=0.5pt,
              >={Stealth[length=1.6mm,width=1.2mm]},
              every node/.append style={align=center, inner sep=2.4pt}},
  % blocks: one radius, two weights; hue = role, dash = redundant channel
  blk/.style   ={rounded corners=1.5pt, inner sep=3pt, align=center},
  ops/.style   ={blk, draw=opsred,   fill=opsred!14,   line width=0.8pt},
  trn/.style   ={blk, draw=trngreen, fill=trngreen!14, line width=0.8pt,
                 dash pattern=on 2pt off 1.2pt},
  froz/.style  ={blk, draw=frozgray, fill=black!7,     line width=0.5pt},
  data/.style  ={blk, draw=datablue, fill=datablue!12, line width=0.5pt},
  plain/.style ={blk, draw=black!45, fill=white,       line width=0.5pt},
  warn/.style  ={blk, draw=warnred,  fill=warnred!8,   line width=0.5pt,
                 dash pattern=on 1pt off 1pt},
  % a carrier = established practice the operator rides on (fig1 v5 local
  % style, promoted here so every figure spells it the same way)
  carrier/.style={blk, draw=opsred, fill=opsred!5, line width=0.5pt,
                 dash pattern=on 2.6pt off 1.2pt},
  gate/.style  ={draw=opsred, fill=opsred!14, line width=0.8pt,
                 diamond, aspect=2.15, inner sep=1pt, align=center},
  grp/.style   ={draw=black!40, densely dotted, line width=0.5pt,
                 rounded corners=2pt, inner sep=0pt},
  grptitle/.style={font=\footnotesize\bfseries, text=black!80, inner sep=1.5pt},
  % type: exactly two sizes; sub-labels at 75% black (contrast >= 4.5:1)
  ttl/.style   ={font=\footnotesize\bfseries},
  lbl/.style   ={font=\footnotesize},
  sub/.style   ={font=\scriptsize, text=black!75},
  subk/.style  ={font=\footnotesize, text=black!80},
  qty/.style   ={font=\footnotesize, text=black!80},
  tsr/.style   ={font=\scriptsize\ttfamily, text=black!75},
  panl/.style  ={font=\footnotesize\bfseries, text=black!85, inner sep=1pt},
  badgecar/.style={badge, draw=opsred, fill=white, text=opsred, line width=0.5pt},
  % arrows: one head, 0.5 pt flow, 0.8 pt emphasis, dashed auxiliary
  flow/.style  ={-{Stealth[length=1.6mm,width=1.2mm]}, draw=flowink,
                 line width=0.5pt},
  flowt/.style ={flow, rounded corners=1.5pt},
  flowh/.style ={-{Stealth[length=1.9mm,width=1.4mm]}, draw=flowink,
                 line width=0.8pt},
  aux/.style   ={-{Stealth[length=1.4mm,width=1.1mm]}, draw=black!45,
                 line width=0.5pt, densely dashed},
  auxw/.style  ={aux, draw=warnred!75},
  rail/.style  ={draw=black!30, line width=0.5pt},
  brace/.style ={decorate, decoration={brace, amplitude=3pt, raise=1pt},
                 draw=black!50, line width=0.5pt},
  zoomray/.style={draw=opsred!55, line width=0.5pt, dash pattern=on 1.3pt off 1.1pt},
  zoombox/.style={draw=opsred, line width=0.5pt},
}

% ---------------- T7: rules inside math must obey the stroke floor --------
% LaTeX draws a \frac bar at the math extension font's default_rule_thickness,
% which SCALES WITH THE TYPE SIZE: 0.4 pt at 10 pt becomes 0.279 pt at the 7 pt
% \scriptsize used for figure sub-labels -- the single sub-0.5 pt stroke that
% verify_fig.py reports in fig1_framework.pdf (a \frac{1}{3} at (166.7,38.3)).
% \genfrac states the rule thickness explicitly, so it survives the floor.
% Measured: \frac -> 0.279 PDF pt, \safefrac -> 0.498 PDF pt (= 0.5 TeX pt).
% The alternative with no rule at all is to write $(1/3)\sum_s\theta_s$.
\newcommand{\safefrac}[2]{\genfrac{}{}{0.5pt}{}{#1}{#2}}

% ---------------- T8: no Type-3 bitmap fonts ----------------------------
% \ttfamily under [T1]{fontenc} with no Type-1 EC typewriter installed makes
% pdflatex fall back to the ec-tt BITMAP font (ectt0800.pk), which is embedded
% as a TYPE 3 font.  Measured on figm_tta.pdf: "F74  Type 3  Custom  emb=yes
% sub=no uni=no" (pdffonts), and figm_tta.log line 655 shows \ectt0800.pk.
% IEEE asks authors to "either convert all text to outlines or embed the
% fonts" and PDF eXpress flags Type 3; the tsr/ style below sets \ttfamily, so
% the substitution must be global.  lmtt is Latin Modern Typewriter: a real
% Type-1 T1 face.  A/B measured: without -> "F39 Type 3", with -> "LMMono8
% Type 1".  Only the typewriter family moves; mathptmx keeps Times for text.
\renewcommand{\ttdefault}{lmtt}

% ---------------- T6: legend row ----------------
% \legendswatch{x}{y}{style}{text} : a 4.2 x 3 mm swatch in <style> with its
% 7 pt label to the right.  Lay the swatches on one baseline y.
\newcommand{\legendswatch}[4]{%
  \draw[#3] (#1,{#2-0.15}) rectangle ({#1+0.42},{#2+0.15});
  \node[sub, anchor=west] at ({#1+0.54},#2) {#4};}

% =========================================================================
% T5: 3D tensor pics (PlotNeuralNet idiom, geometry from the tensor shape)
%
%   \pic (c2) at (x,y) {tensor={px=200, ch=256, img=feat_c2_x4.png,
%                               caption={$256@200^2$}}};
%   \draw[flow] (x-east) -- (c2-west);
%
% keys (all optional; defaults in brackets)
%   px       spatial side in pixels           -> drawn side = zoom*px/400 cm  [200]
%   ch       channel count                    -> depth = 0.04+0.055*log2(ch/256) cm,
%                                                floored at 0.04 cm           [256]
%   depth    explicit depth in cm, overrides ch (hand-sized blocks, UNet)   [ ]
%   zoom     magnification inside a zoom inset (G4)                          [1]
%   color    any xcolor expression; tints are derived from it               [datablue]
%   img      raster placed on the FRONT face, path relative to \assetdir     [ ]
%   hidden   1 = draw the three hidden edges dashed (the PlotNeuralNet look) [0]
%   caption  7 pt text set under the front face                             [ ]
% anchors (prefix = pic name)
%   -west -east          mid-points of the front face's left / right edge (arrows)
%   -westc -eastc        the same at half depth (for rails behind the block)
%   -north -south        mid-points of the front face's top / bottom edge
%   -front -back         centre of the front / back face
%   -nw -ne -se -sw      front-face corners;  -topne  far top-right corner
% Projection: front face in the page plane (so \includegraphics needs no
% transform), depth receding at 30 degrees to the upper right, identical to
% \fmslab, so the two glyph families can share one figure.
% =========================================================================
\tikzset{
  tensor/.pic={
    \tikzset{/tensor/.cd,#1}
    \colorlet{tnc}{\tncolor}
    \pgfmathsetmacro{\tnW}{\tnzoom*(\tnpx)/400}
    \pgfmathsetmacro{\tnH}{0.5*\tnW}
    \edef\tnDepthUser{\tndepth}
    \ifx\tnDepthUser\empty
      \pgfmathsetmacro{\tnD}{\tnzoom*max(0.04, 0.04+0.055*ln((\tnch)/256)/ln(2))}
    \else
      \pgfmathsetmacro{\tnD}{\tndepth}
    \fi
    \begin{scope}[x={(1cm,0cm)}, y={(0cm,1cm)}, z={(-0.866cm,-0.5cm)}]
      % one light source: top brightest, right middle, front darkest
      \fill[tnc!30] (-\tnH,\tnH,0) -- (-\tnH,\tnH,-\tnD) -- (\tnH,\tnH,-\tnD) -- (\tnH,\tnH,0) -- cycle;
      \fill[tnc!20] (\tnH,-\tnH,0) -- (\tnH,-\tnH,-\tnD) -- (\tnH,\tnH,-\tnD) -- (\tnH,\tnH,0) -- cycle;
      \fill[tnc!11] (-\tnH,-\tnH,0) -- (\tnH,-\tnH,0) -- (\tnH,\tnH,0) -- (-\tnH,\tnH,0) -- cycle;
      \edef\tnImgFile{\tnimg}
      \ifx\tnImgFile\empty\else
        \node[inner sep=0pt, outer sep=0pt] at (0,0,0)
          {\includegraphics[width=\tnW cm, height=\tnW cm]{\assetdir/\tnimg}};
      \fi
      \ifnum\tnhidden=1
        \draw[tnc!60, line width=0.5pt, densely dashed]
          (-\tnH,-\tnH,0) -- (-\tnH,-\tnH,-\tnD) -- (\tnH,-\tnH,-\tnD)
          (-\tnH,-\tnH,-\tnD) -- (-\tnH,\tnH,-\tnD);
      \fi
      \draw[tnc!75, line width=0.5pt]
        (-\tnH,-\tnH,0) -- (\tnH,-\tnH,0) -- (\tnH,\tnH,0) -- (-\tnH,\tnH,0) -- cycle
        (-\tnH,\tnH,0) -- (-\tnH,\tnH,-\tnD) -- (\tnH,\tnH,-\tnD) -- (\tnH,\tnH,0)
        (\tnH,\tnH,-\tnD) -- (\tnH,-\tnH,-\tnD) -- (\tnH,-\tnH,0);
      \coordinate (-front) at (0,0,0);
      \coordinate (-back)  at (0,0,-\tnD);
      \coordinate (-west)  at (-\tnH,0,0);
      \coordinate (-east)  at (\tnH,0,0);
      \coordinate (-westc) at (-\tnH,0,-0.5*\tnD);
      \coordinate (-eastc) at (\tnH,0,-0.5*\tnD);
      \coordinate (-north) at (0,\tnH,0);
      \coordinate (-south) at (0,-\tnH,0);
      \coordinate (-nw) at (-\tnH,\tnH,0);  \coordinate (-ne) at (\tnH,\tnH,0);
      \coordinate (-sw) at (-\tnH,-\tnH,0); \coordinate (-se) at (\tnH,-\tnH,0);
      \coordinate (-topne) at (\tnH,\tnH,-\tnD);
    \end{scope}
    \edef\tnCap{\tncaption}
    \ifx\tnCap\empty\else
      \node[sub, anchor=north, inner sep=1.5pt] at (0,-\tnH) {\tncaption};
    \fi
  },
  % a pile of k offset H x W planes (mask stack); front plane may carry a raster
  %   \pic (m) at (x,y) {tensorpile={px=800, ch=3, img=pred_masks.png}};
  % here ch = number of planes drawn.
  tensorpile/.pic={
    \tikzset{/tensor/.cd,#1}
    \colorlet{tnc}{\tncolor}
    \pgfmathsetmacro{\tpW}{\tnzoom*(\tnpx)/400}
    \pgfmathsetmacro{\tpH}{0.5*\tpW}
    \pgfmathsetmacro{\tpO}{0.06*\tnzoom}
    \pgfmathtruncatemacro{\tpN}{\tnch-1}
    \foreach \i in {\tpN,...,0} {
      \pgfmathsetmacro{\d}{\i*\tpO}
      \fill[white]  ({-\tpH+\d},{-\tpH+\d}) rectangle ({\tpH+\d},{\tpH+\d});
      \fill[tnc!8]  ({-\tpH+\d},{-\tpH+\d}) rectangle ({\tpH+\d},{\tpH+\d});
      \ifnum\i=0 \edef\tpImg{\tnimg}\ifx\tpImg\empty\else
        \node[inner sep=0pt, outer sep=0pt] at (0,0)
          {\includegraphics[width=\tpW cm, height=\tpW cm]{\assetdir/\tnimg}};\fi\fi
      \draw[tnc!60, line width=0.5pt] ({-\tpH+\d},{-\tpH+\d}) rectangle ({\tpH+\d},{\tpH+\d});
    }
    \coordinate (-west)  at (-\tpH,0);
    \coordinate (-east)  at ({\tpH+\tpN*\tpO},0);
    \coordinate (-north) at (0,{\tpH+\tpN*\tpO});
    \coordinate (-south) at (0,-\tpH);
    \coordinate (-front) at (0,0);
    \edef\tnCap{\tncaption}
    \ifx\tnCap\empty\else
      \node[sub, anchor=north, inner sep=1.5pt] at (0,-\tpH) {\tncaption};
    \fi
  },
  /tensor/.search also={/tikz},
  /tensor/.cd,
  px/.store in=\tnpx,           px=200,
  ch/.store in=\tnch,           ch=256,
  zoom/.store in=\tnzoom,       zoom=1,
  color/.store in=\tncolor,     color=datablue,
  img/.store in=\tnimg,         img=,
  depth/.store in=\tndepth,     depth=,
  hidden/.store in=\tnhidden,   hidden=0,
  caption/.store in=\tncaption, caption=,
}

% instead of % =========================================================================
% tikz_style.tex — shared visual system for the Section III method figures.
% Usage: % =========================================================================
% tikz_style.tex — shared visual system for the Section III method figures.
% Usage: \input{tikz_style} in the preamble of a standalone figure.
%
% -------------------------------------------------------------------------
% DESIGN NOTES — conventions extracted from an accepted TGRS manuscript
% (Gao et al., "BBox-free ... SAR ship instance segmentation", TGRS 2024;
%  local copy: scratchpad/gao_bboxfree_tgrs24.pdf).  Observations taken from
%  its Fig. 2 (overall architecture, p. 5), Fig. 5 (FD module, p. 8) and
%  Fig. 6 (conceptual mask panel, p. 8):
%
%  O1  REGION PARTITION.  The architecture figure is carved into 3-4 named
%      regions by DOTTED rectangles with no fill; the region name is set in
%      bold OUTSIDE the rectangle ("Feature Extractor", "Classification
%      Branch", "Localization Branch", "Center Point Extractor").  Nothing
%      is nested more than two levels deep.
%  O2  COLOR CARRIES MEANING, NOT DECORATION.  Only ~5 hues appear, all as
%      pale fills with black text: light gray = ordinary/standard module,
%      salmon-orange = the module the paper contributes (Feature Decoupling
%      Module, Center Point Extractor), yellow = a distinguished operation
%      (Feature Fusion), light blue / light pink = tensors.  Saturated color
%      is reserved for the few glyphs that must be found first.
%  O3  LEGEND PANEL.  Both Fig. 2 and Fig. 5 carry an explicit legend inside
%      the figure (dotted box top-left in Fig. 2, dashed box right in Fig. 5)
%      that decodes arrow colors and single-letter glyph nodes.  The body of
%      the diagram is kept uncluttered by pushing vocabulary into the legend.
%  O4  TENSOR ICONOGRAPHY.  Feature maps are drawn as thin vertical plates,
%      multi-scale features as 3 offset stacked planes, channel vectors as
%      segmented 3-D bars, the backbone as an isometric slab.  Real SAR
%      chips / heat maps / GT overlays are pasted directly into the flow.
%  O5  TENSOR NAMES SIT BESIDE THE GLYPH, NEVER INSIDE IT, in math italic
%      (F_out, F_mask, F_param, F_H_zip ...), placed above-right of the plate.
%  O6  ARROW DISCIPLINE.  Data flow is a thin gray/black line with a small
%      solid triangular head; routing is orthogonal (Manhattan) with a shared
%      trunk when one source fans out to several targets, so lines never
%      cross.  Semantically special transitions get a THICK COLORED BLOCK
%      arrow instead (2x upsample / 2x downsample / connect), i.e. arrow
%      style is itself a semantic channel.  Long skip paths are routed as
%      rails around the outside of the diagram (Fig. 5 top and bottom).
%  O7  TYPOGRAPHY.  One size for block labels, one smaller size for
%      sub-labels, math only in math mode; no stretched or rotated text
%      except a rare 90-degree axis label.
%  O8  ASPECT RATIO.  Double-column method figures are wide and short:
%      Fig. 2 is 530 x 200 pt (2.65:1), Fig. 5 is 530 x 215 pt (2.5:1).
%      They are laid on a grid with equal gutters between parallel branches.
%  O9  CAPTION DIVISION OF LABOR.  The caption names the parts and decodes
%      abbreviations/colors; it does not re-derive the mechanism, which is
%      left to the body text.  The figure is readable without the caption.
%  O10 CONCEPTUAL PANELS (Fig. 6) use dashed sub-panels with a centered
%      label underneath, a pale blue blob for the true object and a pale red
%      overlay for the error region -- schematic, not photorealistic.
%
% -------------------------------------------------------------------------
% WHAT WE ADOPT / WHERE WE DEVIATE
%  + adopt O1 (dotted named regions), O3 (in-figure legend), O4/O5 (tensor
%    glyphs with the symbol outside), O6 (one flow style + Manhattan
%    routing), O8 (wide-short double-column ratios), O10 (schematic panels).
%  ~ deviate on O2: instead of hue-as-identity we bind hue to ROLE in the
%    contribution -- frozen / training-free operator / trained corrector /
%    data -- because the entire argument of this article is which parts are
%    touched and which are not.
%  ~ deviate on the sans-serif block labels of Fig. 2: we set every figure in
%    the body Times face (as the same paper's Fig. 6 does), so figure text
%    and caption text share one face.
%  + every figure is authored at its EXACT printed width (18.2 cm double
%    column = 516 pt, 8.89 cm single column = 252 pt) and included at
%    scale 1, so \footnotesize is a true 8 pt and \scriptsize a true 7 pt on
%    the page in ALL figures.  Use \figcanvasd / \figcanvasc below.
%
% -------------------------------------------------------------------------
% SEMANTIC PALETTE (Okabe-Ito colorblind-safe base; redundant encoding by
% border weight, dash pattern and badge glyph so it survives grayscale):
%   froz  gray   + lock glyph        frozen, not modified by us  (Stage A)
%   ops   vermil + solid heavy rule  our TRAINING-FREE operator  (rung 1)
%   trn   green  + dashed rule       our OPTIONAL TRAINED corrector (2-3)
%   data  blue   + thin rule         data, tensors, I/O
%   warn  red    + dotted rule, small  failure mode / negative control
% =========================================================================

% Real rasters are embedded on the front faces of the tensor glyphs (see
% \realpanel / \fmslab below).  graphicx must be loaded in the preamble, which
% is where this file is \input.
\usepackage{graphicx}
% Single switch for the raster payload: \def\assetdir{assets_light}
% before \begin{document} swaps every embedded raster for its
% print-sized ~600 dpi copy (see figs/make_light_assets.py).  The
% default is the native dump, 965-1016 dpi at the drawn size.
\providecommand{\assetdir}{assets}
\usetikzlibrary{positioning,arrows.meta,fit,backgrounds,calc,shapes.geometric,
                decorations.pathreplacing,decorations.markings,patterns}

% ---------------- palette ----------------
\definecolor{opsred}{HTML}{D55E00}   % Okabe-Ito vermillion  -- training-free ops
\definecolor{trngreen}{HTML}{009E73} % Okabe-Ito bluish green-- trained correctors
\definecolor{datablue}{HTML}{0072B2} % Okabe-Ito blue        -- data / tensors
\definecolor{frozgray}{HTML}{6E6E6E} % neutral               -- frozen base
\definecolor{flowink}{HTML}{262626}  % near black            -- data-flow arrows
\definecolor{warnred}{HTML}{9E2B2B}  % muted brick           -- failure modes
\definecolor{sarink}{HTML}{2B2B33}   % SAR chip background

% ---------------- exact-width canvases ----------------
% \figcanvasd{H} : double column, 18.2 cm = 516 pt = \textwidth  (scale 1.0)
% \figcanvasc{H} : single column,  8.89 cm = 252 pt = \columnwidth
\newcommand{\figcanvasd}[1]{\useasboundingbox (0,0) rectangle (18.2,#1);}
\newcommand{\figcanvasc}[1]{\useasboundingbox (0,0) rectangle (8.89,#1);}

\tikzset{
  % ---- picture-level defaults: ONE font size system for all five figures ----
  fig/.style={font=\footnotesize, line width=0.65pt,
              >={Stealth[length=1.7mm,width=1.35mm]},
              every node/.append style={align=center, inner sep=2.4pt}},
  % ---- blocks: role -> color + redundant border encoding ----
  % Fill tints were raised in v3: at 4-8% the five roles were separated by an
  % RGB distance of only 6.6-21 (out of 441) and all landed on gray 243-248,
  % i.e. the fill channel was consuming ink and returning nothing.  The border
  % weight / dash pattern / badge remain the PRIMARY, grayscale-safe channel.
  ops/.style   ={draw=opsred,   fill=opsred!16,   line width=0.9pt,
                 rounded corners=1.6pt, inner sep=3pt, align=center},
  trn/.style   ={draw=trngreen, fill=trngreen!15, line width=0.9pt,
                 dash pattern=on 1.7pt off 1.1pt,
                 rounded corners=1.6pt, inner sep=3pt, align=center},
  froz/.style  ={draw=frozgray, fill=black!8,     line width=0.7pt,
                 rounded corners=1.6pt, inner sep=3pt, align=center},
  data/.style  ={draw=datablue, fill=datablue!14, line width=0.55pt,
                 rounded corners=1.6pt, inner sep=3pt, align=center},
  plain/.style ={draw=black!45, fill=white,       line width=0.55pt,
                 rounded corners=1.6pt, inner sep=3pt, align=center},
  % v4: radius 1.4 -> 1.6 pt, the one block style that did not share the
  % system's corner radius; weight 0.5 -> 0.6 pt so the failure chips are not
  % the faintest object on the page while being the ones a reader must find.
  % The dotted dash pattern stays as the grayscale-safe channel.
  warn/.style  ={draw=warnred,  fill=warnred!9,   line width=0.6pt,
                 dash pattern=on 1.1pt off 1.1pt,
                 rounded corners=1.6pt, inner sep=3pt, align=center},
  gate/.style  ={draw=opsred,   fill=opsred!16,   line width=0.9pt,
                 diamond, aspect=2.15, inner sep=1pt, align=center},
  % ---- named region (O1): dotted, unfilled, label set outside ----
  grp/.style   ={draw=black!45, densely dotted, line width=0.7pt,
                 rounded corners=2pt, inner sep=0pt},
  grptitle/.style={font=\footnotesize\bfseries, text=black!80, inner sep=1.5pt},
  % ---- step badges (redundant, grayscale-safe stage identity) ----
  badge/.style ={circle, inner sep=0pt, minimum size=3.3mm, text=white,
                 font=\scriptsize\bfseries},
  badgeops/.style={badge, fill=opsred},
  badgetrn/.style={badge, fill=trngreen},
  % ---- typography (O7): exactly three text sizes in the whole figure set ----
  % v3: 7 pt at 55% black was the worst corner of the design space (minimum
  % size AND minimum contrast) -- sub-labels are now 7 pt at 70% black, and
  % anything carrying a QUANTITY is promoted to 8 pt via \qty / \subk.
  ttl/.style   ={font=\footnotesize\bfseries},
  lbl/.style   ={font=\footnotesize},
  sub/.style   ={font=\scriptsize, text=black!70},
  subk/.style  ={font=\footnotesize, text=black!80},
  qty/.style   ={font=\footnotesize, text=black!80},% tensor dims, AP numbers
  tsr/.style ={font=\scriptsize\ttfamily, text=black!70},% B x C x H x W (O5)
  tag/.style   ={font=\scriptsize, inner sep=1.6pt, rounded corners=1pt},
  % panel letters (B11): \footnotesize\bfseries at the panel's top-left
  panl/.style  ={font=\footnotesize\bfseries, text=black!85, inner sep=1pt},
  % ---- arrows (O6): one data-flow style, one auxiliary/optional style ----
  flow/.style  ={-{Stealth[length=1.7mm,width=1.35mm]}, draw=flowink,
                 line width=0.65pt},
  flowt/.style ={-{Stealth[length=1.7mm,width=1.35mm]}, draw=flowink,
                 line width=0.65pt, rounded corners=1.6pt},
  aux/.style   ={-{Stealth[length=1.5mm,width=1.2mm]}, draw=black!45,
                 line width=0.5pt, densely dashed},
  auxw/.style  ={-{Stealth[length=1.5mm,width=1.2mm]}, draw=warnred!75,
                 line width=0.5pt, densely dashed},
  rail/.style  ={draw=black!30, line width=0.5pt},
  brace/.style ={decorate, decoration={brace, amplitude=3.2pt, raise=1pt},
                 draw=black!50, line width=0.5pt},
  % ---- ONE leader style for every magnifier cone in the figure set --------
  % v4: the zoom cones used to be drawn in the generic dashed `aux' style, so
  % a sight line was typographically the same object as an ordinary auxiliary
  % arrow.  A cone is now its own object: it carries the hue of the locator
  % rectangle it starts from, no arrowhead (it is a correspondence, not a
  % flow), and it always runs corner-to-corner from the locator box to the
  % MAGNIFIED RASTER, never to the enclosing panel frame.
  zoomray/.style={draw=opsred!55, line width=0.5pt,
                  dash pattern=on 1.3pt off 1.1pt},
  zoombox/.style={draw=opsred, line width=0.7pt},
}

% =========================================================================
% \hbarrow{x0}{y}{track}{frac}{colour}{weight}{label}
%   One horizontal bar on a printed 0..1 track, with the numeral set at the
%   bar end.  v4: fig1's AP progression, figm_tta (c), figm_fusion's result
%   panel and figm_rerank (b) all drew this same object with four slightly
%   different geometries.  They now share one macro, so the reader learns the
%   idiom once and meets it four times -- a small multiple across figures.
%   The empty track is always drawn, so a short bar still shows its scale.
%
%   DEPRECATED FOR METHOD FIGURES (docs/FIGURE_REDESIGN_SPEC.md S5, items 15
%   and 16, rule H1).  \hbarrow / \hbarrowr / \baxis print an accuracy number
%   inside the figure, which the v6 redraw forbids: fig1_framework, figm_bpr
%   and figm_fusion no longer call them and their numbers live in Tables IV
%   and VI.  The macros are KEPT, not deleted, because the parked single-
%   column data panel of Fig. 5 (figm_rerank, permutation importance) is a
%   data figure, where a measured bar length is evidence rather than a
%   result.  Same status, same reason: warn/, badge/ + badgeops/ + badgetrn/,
%   \hull and \mplane -- no v6 method figure uses any of them.
% =========================================================================
\newcommand{\hbarrow}[7]{%
  \draw[black!15, line width=0.5pt] (#1,#2) -- ({#1+#3},#2);
  \draw[#5, line width=#6] (#1,#2) -- ({#1+(#3)*(#4)},#2);
  % the numeral gets a paper backing so the empty part of the track does not
  % run into it and read as a stray rule after the digits
  \node[qty, anchor=west, inner sep=1.6pt, fill=white]
    at ({#1+(#3)*(#4)},#2) {#7};}

% \hbarrowr{x0}{y}{track}{frac}{colour}{weight}{label}
%   As \hbarrow, but the numeral is set in a fixed column just past the end
%   of the track instead of at the bar tip.  Use it where several groups sit
%   side by side on a SHORT track: it aligns the numerals into columns across
%   the groups, and it stops the leftover empty track of a short bar from
%   printing as a stray rule immediately after the digits -- at 1.25 cm of
%   track a 27.7/75 bar leaves 3 mm of rule past the numeral, which reads as
%   an en dash and can be mistaken for a range.
\newcommand{\hbarrowr}[7]{%
  \draw[black!15, line width=0.5pt] (#1,#2) -- ({#1+#3},#2);
  \draw[#5, line width=#6] (#1,#2) -- ({#1+(#3)*(#4)},#2);
  \node[qty, anchor=west, inner sep=1.6pt] at ({#1+(#3)+0.04},#2) {#7};}

% \baxis{x0}{y}{track}{lo}{hi}{unit} : the tick pair + caption for a \hbarrow
%   group, drawn once per group rather than once per bar.
\newcommand{\baxis}[6]{%
  \draw[black!45, line width=0.5pt] (#1,#2) -- (#1,{#2-0.05})
                                    ({#1+#3},#2) -- ({#1+#3},{#2-0.05});
  \node[sub, anchor=north, inner sep=1pt] at (#1,{#2-0.05}) {#4};
  \node[sub, anchor=north, inner sep=1pt] at ({#1+#3},{#2-0.05}) {#5};
  \node[sub, anchor=west, inner sep=1pt] at ({#1+#3+0.14},{#2-0.16}) {#6};}

% =========================================================================
% GEOMETRIC ISOMORPHISM (v3).  Tensor shape is DRAWN, never only written.
%
%  G1 SPATIAL EXTENT -> DRAWN SIZE, one scale for the whole figure set:
%       1 mm of drawn edge per 40 pixels of tensor extent,
%       i.e. drawn side [cm] = side [px] / 400   ->  \pxc{800} = 2.00 cm.
%       800 -> 20.0 mm | 200 -> 5.00 mm | 100 -> 2.50 mm |  50 -> 1.25 mm
%        25 ->  0.63 mm | 640 -> 16.0 mm | 960 -> 24.0 mm |1088 -> 27.2 mm
%        64 ->  1.60 mm | 128 ->  3.20 mm
%       A stride-4 map is therefore exactly 8x the linear size of a
%       stride-32 map, as it must be.
%  G2 CHANNEL COUNT -> SLAB DEPTH (isometric extrusion, 30 degrees):
%       depth [cm] = 0.04 + 0.055 * log2(C/256)
%       256 -> 0.40 mm | 512 -> 0.95 mm | 1024 -> 1.50 mm | 2048 -> 2.05 mm
%       i.e. +0.55 mm of depth per channel doubling.
%  G3 CATEGORY -> GLYPH SHAPE (readable in grayscale, no label needed):
%       spatial feature map (B,C,H,W) ......... extruded slab   \fmslab
%       token set (N x C, no spatial extent) .. column of bars  \tokencol
%       mask stack (N x H x W) ................ offset planes   \maskstack
%       scalar score .......................... short bar       \scorebar
%       feature vector (d cells) .............. strip of cells  (drawn inline)
%  G4 A LABELLED ZOOM INSET may magnify a sub-scene by an explicit factor z;
%       inside it G1/G2 hold multiplied by z.  Every macro takes z as its
%       last argument (z=1 for the global scale).
% =========================================================================
\newcommand{\pxc}[1]{(#1)/400}     % tensor side in px -> drawn side in cm

% =========================================================================
% G5 REAL DATA (v3).  Wherever a glyph stands for a tensor we actually own,
%     the glyph's FRONT FACE carries that tensor's real raster, dumped by
%     carrier/dump_figure_assets.py with provenance in assets/manifest.json.
%     The vector chrome (frames, arrows, labels, extrusion) stays vector, so
%     type remains selectable and the fonts stay embedded -- the bitmap-panel
%     + vector-chrome architecture Nature's artwork guide prescribes.
%     Resolution guard: every embedded raster must clear 300 dpi at its drawn
%     size and we clear 600 dpi everywhere (chip.png 800 px at 20 mm = 1016
%     dpi).  Feature maps are embedded from the nearest-neighbour _xN variants
%     so the PDF renderer cannot smooth pixels the tensor does not have.
%
% \realpanel{x}{y}{sidepx}{zoom}{file}{framespec} : a real raster, framed and
%     sized by G1 exactly as a schematic glyph of the same tensor would be.
% =========================================================================
\newcommand{\realpanel}[6]{%
  \begin{scope}[shift={(#1,#2)}]
    \pgfmathsetmacro{\rpH}{0.5*#4*\pxc{#3}}
    \pgfmathsetmacro{\rpS}{2*\rpH}
    \node[inner sep=0pt, outer sep=0pt] at (0,0)
      {\includegraphics[width=\rpS cm, height=\rpS cm]{\assetdir/#5}};
    \draw[#6] (-\rpH,-\rpH) rectangle (\rpH,\rpH);
  \end{scope}}

% \realpanelh{...} : the same panel HORIZONTALLY MIRRORED.  Used by the TTA
%     figure so the hflip half of each view pair is a real flip of the real
%     raster, not a redrawn mirror of a schematic.
\newcommand{\realpanelh}[6]{%
  \begin{scope}[shift={(#1,#2)}]
    \pgfmathsetmacro{\rpH}{0.5*#4*\pxc{#3}}
    \pgfmathsetmacro{\rpS}{2*\rpH}
    \node[inner sep=0pt, outer sep=0pt] at (0,0)
      {\reflectbox{\includegraphics[width=\rpS cm, height=\rpS cm]{\assetdir/#5}}};
    \draw[#6] (-\rpH,-\rpH) rectangle (\rpH,\rpH);
  \end{scope}}

% --- \fmslab{x}{y}{sidepx}{C}{colour}{zoom} : a spatial feature map -------
\newcommand{\fmslab}[6]{\fmslabr{#1}{#2}{#3}{#4}{#5}{#6}{}}

% --- \fmslabr{...}{file} : the same slab with a REAL raster on the front
%     face.  Pass an empty 7th argument for the flat-tinted variant.  The
%     extruded top/right faces stay flat tinted, so G2 (depth = channel
%     count) and the single light source (top brightest, right middle,
%     front darkest) are preserved.
\newcommand{\fmslabr}[7]{%
  \begin{scope}[shift={(#1,#2)}]
    \pgfmathsetmacro{\sbH}{0.5*#6*\pxc{#3}}
    \pgfmathsetmacro{\sbD}{#6*(0.04+0.055*ln((#4)/256)/ln(2))}
    \pgfmathsetmacro{\sbX}{0.866*\sbD}
    \pgfmathsetmacro{\sbY}{0.500*\sbD}
    \pgfmathsetmacro{\sbW}{2*\sbH}
    \fill[#5!30] (-\sbH,\sbH) -- (-\sbH+\sbX,\sbH+\sbY)
                 -- (\sbH+\sbX,\sbH+\sbY) -- (\sbH,\sbH) -- cycle;
    \fill[#5!20] (\sbH,\sbH) -- (\sbH+\sbX,\sbH+\sbY)
                 -- (\sbH+\sbX,-\sbH+\sbY) -- (\sbH,-\sbH) -- cycle;
    \fill[#5!11] (-\sbH,-\sbH) rectangle (\sbH,\sbH);
    \edef\fmfile{#7}%
    \ifx\fmfile\empty\else
      \node[inner sep=0pt, outer sep=0pt] at (0,0)
        {\includegraphics[width=\sbW cm, height=\sbW cm]{\assetdir/#7}};
    \fi
    \draw[#5!75, line width=0.5pt]
      (-\sbH,-\sbH) rectangle (\sbH,\sbH)
      (-\sbH,\sbH) -- (-\sbH+\sbX,\sbH+\sbY) -- (\sbH+\sbX,\sbH+\sbY) -- (\sbH,\sbH)
      (\sbH+\sbX,\sbH+\sbY) -- (\sbH+\sbX,-\sbH+\sbY) -- (\sbH,-\sbH);
  \end{scope}}

% --- \tokencol{x}{y}{nbars}{width}{colour} : an N x C token set -----------
%     Deliberately NOT a slab: flat bars, visible gaps, no extrusion, so a
%     query set can never be mistaken for a feature map.
\newcommand{\tokencol}[5]{%
  \begin{scope}[shift={(#1,#2)}]
    \pgfmathsetmacro{\tkP}{0.098}
    \pgfmathsetmacro{\tkT}{0.056}
    \pgfmathsetmacro{\tkW}{0.5*#4}
    \pgfmathtruncatemacro{\tkN}{#3-1}
    \pgfmathsetmacro{\tkY}{0.5*(\tkN*\tkP+\tkT)}
    \foreach \tki in {0,...,\tkN} {
      \pgfmathsetmacro{\tkc}{\tkY-\tki*\tkP}
      \fill[#5!13] (-\tkW,\tkc-\tkT) rectangle (\tkW,\tkc);
      \draw[#5!75, line width=0.5pt] (-\tkW,\tkc-\tkT) rectangle (\tkW,\tkc);
    }
  \end{scope}}

% --- \maskstack{x}{y}{sidepx}{k}{colour}{zoom} : k planes of H x W --------
\newcommand{\maskstack}[6]{\maskstackr{#1}{#2}{#3}{#4}{#5}{#6}{}}

% --- \maskstackr{...}{file} : the same pile with the FRONT plane carrying a
%     real raster; the planes behind stay schematic so the "stack of N"
%     reading survives.  Empty 7th argument = the all-schematic variant.
\newcommand{\maskstackr}[7]{%
  \begin{scope}[shift={(#1,#2)}]
    \pgfmathsetmacro{\msH}{0.5*#6*\pxc{#3}}
    \pgfmathsetmacro{\msO}{0.30*\msH/2.5}
    \pgfmathtruncatemacro{\msN}{#4-1}
    \pgfmathsetmacro{\msS}{\msH/0.92}
    \pgfmathsetmacro{\msW}{2*\msH}
    \edef\msfile{#7}%
    \foreach \msi in {\msN,...,0} {
      \pgfmathsetmacro{\msd}{\msi*\msO}
      \pgfmathtruncatemacro{\msreal}{(\msi==0) ? 1 : 0}
      \edef\mshasimg{0}%
      \ifx\msfile\empty\else\ifnum\msreal=1 \edef\mshasimg{1}\fi\fi
      \begin{scope}[shift={(\msd,\msd)}]
        \fill[white] (-\msH,-\msH) rectangle (\msH,\msH);
        \fill[#5!8]  (-\msH,-\msH) rectangle (\msH,\msH);
        \ifnum\mshasimg=1
          \node[inner sep=0pt, outer sep=0pt] at (0,0)
            {\includegraphics[width=\msW cm, height=\msW cm]{\assetdir/#7}};
        \fi
        \draw[#5!60, line width=0.5pt] (-\msH,-\msH) rectangle (\msH,\msH);
        \ifnum\mshasimg=0
          \begin{scope}[scale=\msS]
            \fill[#5!35] \hull;
            \draw[#5!90, line width=0.5pt] \hull;
          \end{scope}
        \fi
      \end{scope}}
  \end{scope}}

% --- \scorebar{x}{y}{h}{colour} : one scalar, as a short upright bar ------
\newcommand{\scorebar}[4]{%
  \begin{scope}[shift={(#1,#2)}]
    \draw[black!35, line width=0.5pt] (-0.10,0) rectangle (0.10,0.62);
    \fill[#4!55] (-0.10,0) rectangle (0.10,#3);
    \draw[#4!90, line width=0.5pt] (-0.10,0) rectangle (0.10,#3);
  \end{scope}}

% =========================================================================
% Reusable glyphs.  All are TikZ pics: \pic[scale=s] at (x,y) {name};
% =========================================================================

% --- lock: marks any component whose weights we never update -------------
\tikzset{lock/.pic={
  \fill[frozgray] (-0.105,-0.105) rectangle (0.105,0.055);
  \draw[frozgray, line width=0.65pt]
    (-0.060,0.055) -- (-0.060,0.105)
    arc[start angle=180, end angle=0, radius=0.060] -- (0.060,0.055);
  \fill[white] (0,-0.022) circle[radius=0.020];
}}

% --- planes: 3 offset stacked planes = a multi-scale feature pyramid -----
\tikzset{planes/.pic={
  \foreach \i/\o in {2/0.20, 1/0.10, 0/0.00} {
    \fill[datablue!\the\numexpr 10+8*\i\relax]
      (\o-0.24,\o-0.30) rectangle (\o+0.24,\o+0.30);
    \draw[datablue!70, line width=0.5pt]
      (\o-0.24,\o-0.30) rectangle (\o+0.24,\o+0.30);
  }
}}

% --- plate: a single feature plane (thin vertical plate, O4) -------------
\tikzset{plate/.pic={
  \fill[datablue!14] (-0.10,-0.34) rectangle (0.10,0.34);
  \draw[datablue!70, line width=0.5pt] (-0.10,-0.34) rectangle (0.10,0.34);
  \draw[datablue!45, line width=0.5pt] (-0.10,0.20) -- (0.10,0.20);
  \draw[datablue!45, line width=0.5pt] (-0.10,-0.20) -- (0.10,-0.20);
}}

% --- ckpt: an isometric slab = one stored checkpoint / weight vector -----
\tikzset{ckpt/.pic={
  \fill[frozgray!18] (-0.30,-0.13) rectangle (0.24,0.13);
  \fill[frozgray!32] (-0.30,0.13) -- (-0.24,0.21) -- (0.30,0.21) -- (0.24,0.13) -- cycle;
  \fill[frozgray!26] (0.24,-0.13) -- (0.30,-0.05) -- (0.30,0.21) -- (0.24,0.13) -- cycle;
  \draw[frozgray, line width=0.5pt]
    (-0.30,-0.13) rectangle (0.24,0.13);
  \draw[frozgray, line width=0.5pt]
    (-0.30,0.13) -- (-0.24,0.21) -- (0.30,0.21) -- (0.24,0.13)
    (0.30,0.21) -- (0.30,-0.05) -- (0.24,-0.13);
}}

% --- masked: pictogram of an output instance mask (hull on light field) --
\tikzset{maskicon/.pic={
  \fill[datablue!5] (-0.30,-0.30) rectangle (0.30,0.30);
  \draw[datablue!60, line width=0.5pt] (-0.30,-0.30) rectangle (0.30,0.30);
  \begin{scope}[rotate=28]
    \fill[datablue!45] (0,0.05) ellipse[x radius=0.235, y radius=0.062];
    \draw[datablue, line width=0.5pt] (0,0.05) ellipse[x radius=0.235, y radius=0.062];
  \end{scope}
  \begin{scope}[rotate=-12]
    \fill[datablue!45] (0.02,-0.19) ellipse[x radius=0.145, y radius=0.045];
    \draw[datablue, line width=0.5pt] (0.02,-0.19) ellipse[x radius=0.145, y radius=0.045];
  \end{scope}
}}

% =========================================================================
% \sarchip was REMOVED in v3.  It drew two smooth ellipses each ~2.6% of the
% frame on a bokeh-like field; the real HRSID inshore median instance covers
% 0.028% of the chip (13.4 px on 800x800), so the glyph was ~90x too large and
% contradicted the small-target-in-clutter argument the paper makes.  Every
% call site now embeds the REAL chip raster instead (assets/chip.png,
% assets/view_{640,800,960,1088}.png; see the realpanel macro above and
% assets/manifest.json).
% =========================================================================

% =========================================================================
% \hull{}  -- schematic ship-mask contour, ~1.3 cm wide, centred at origin.
% Used by the BPR and fusion figures so every ship outline in the paper is
% literally the same curve.
% =========================================================================
\newcommand{\hull}{plot[smooth cycle, tension=0.72] coordinates
  {(0.72,0.00) (0.34,0.26) (-0.22,0.33) (-0.64,0.17) (-0.62,-0.12)
   (-0.10,-0.30) (0.42,-0.22)}}
 in the preamble of a standalone figure.
%
% -------------------------------------------------------------------------
% DESIGN NOTES — conventions extracted from an accepted TGRS manuscript
% (Gao et al., "BBox-free ... SAR ship instance segmentation", TGRS 2024;
%  local copy: scratchpad/gao_bboxfree_tgrs24.pdf).  Observations taken from
%  its Fig. 2 (overall architecture, p. 5), Fig. 5 (FD module, p. 8) and
%  Fig. 6 (conceptual mask panel, p. 8):
%
%  O1  REGION PARTITION.  The architecture figure is carved into 3-4 named
%      regions by DOTTED rectangles with no fill; the region name is set in
%      bold OUTSIDE the rectangle ("Feature Extractor", "Classification
%      Branch", "Localization Branch", "Center Point Extractor").  Nothing
%      is nested more than two levels deep.
%  O2  COLOR CARRIES MEANING, NOT DECORATION.  Only ~5 hues appear, all as
%      pale fills with black text: light gray = ordinary/standard module,
%      salmon-orange = the module the paper contributes (Feature Decoupling
%      Module, Center Point Extractor), yellow = a distinguished operation
%      (Feature Fusion), light blue / light pink = tensors.  Saturated color
%      is reserved for the few glyphs that must be found first.
%  O3  LEGEND PANEL.  Both Fig. 2 and Fig. 5 carry an explicit legend inside
%      the figure (dotted box top-left in Fig. 2, dashed box right in Fig. 5)
%      that decodes arrow colors and single-letter glyph nodes.  The body of
%      the diagram is kept uncluttered by pushing vocabulary into the legend.
%  O4  TENSOR ICONOGRAPHY.  Feature maps are drawn as thin vertical plates,
%      multi-scale features as 3 offset stacked planes, channel vectors as
%      segmented 3-D bars, the backbone as an isometric slab.  Real SAR
%      chips / heat maps / GT overlays are pasted directly into the flow.
%  O5  TENSOR NAMES SIT BESIDE THE GLYPH, NEVER INSIDE IT, in math italic
%      (F_out, F_mask, F_param, F_H_zip ...), placed above-right of the plate.
%  O6  ARROW DISCIPLINE.  Data flow is a thin gray/black line with a small
%      solid triangular head; routing is orthogonal (Manhattan) with a shared
%      trunk when one source fans out to several targets, so lines never
%      cross.  Semantically special transitions get a THICK COLORED BLOCK
%      arrow instead (2x upsample / 2x downsample / connect), i.e. arrow
%      style is itself a semantic channel.  Long skip paths are routed as
%      rails around the outside of the diagram (Fig. 5 top and bottom).
%  O7  TYPOGRAPHY.  One size for block labels, one smaller size for
%      sub-labels, math only in math mode; no stretched or rotated text
%      except a rare 90-degree axis label.
%  O8  ASPECT RATIO.  Double-column method figures are wide and short:
%      Fig. 2 is 530 x 200 pt (2.65:1), Fig. 5 is 530 x 215 pt (2.5:1).
%      They are laid on a grid with equal gutters between parallel branches.
%  O9  CAPTION DIVISION OF LABOR.  The caption names the parts and decodes
%      abbreviations/colors; it does not re-derive the mechanism, which is
%      left to the body text.  The figure is readable without the caption.
%  O10 CONCEPTUAL PANELS (Fig. 6) use dashed sub-panels with a centered
%      label underneath, a pale blue blob for the true object and a pale red
%      overlay for the error region -- schematic, not photorealistic.
%
% -------------------------------------------------------------------------
% WHAT WE ADOPT / WHERE WE DEVIATE
%  + adopt O1 (dotted named regions), O3 (in-figure legend), O4/O5 (tensor
%    glyphs with the symbol outside), O6 (one flow style + Manhattan
%    routing), O8 (wide-short double-column ratios), O10 (schematic panels).
%  ~ deviate on O2: instead of hue-as-identity we bind hue to ROLE in the
%    contribution -- frozen / training-free operator / trained corrector /
%    data -- because the entire argument of this article is which parts are
%    touched and which are not.
%  ~ deviate on the sans-serif block labels of Fig. 2: we set every figure in
%    the body Times face (as the same paper's Fig. 6 does), so figure text
%    and caption text share one face.
%  + every figure is authored at its EXACT printed width (18.2 cm double
%    column = 516 pt, 8.89 cm single column = 252 pt) and included at
%    scale 1, so \footnotesize is a true 8 pt and \scriptsize a true 7 pt on
%    the page in ALL figures.  Use \figcanvasd / \figcanvasc below.
%
% -------------------------------------------------------------------------
% SEMANTIC PALETTE (Okabe-Ito colorblind-safe base; redundant encoding by
% border weight, dash pattern and badge glyph so it survives grayscale):
%   froz  gray   + lock glyph        frozen, not modified by us  (Stage A)
%   ops   vermil + solid heavy rule  our TRAINING-FREE operator  (rung 1)
%   trn   green  + dashed rule       our OPTIONAL TRAINED corrector (2-3)
%   data  blue   + thin rule         data, tensors, I/O
%   warn  red    + dotted rule, small  failure mode / negative control
% =========================================================================

% Real rasters are embedded on the front faces of the tensor glyphs (see
% \realpanel / \fmslab below).  graphicx must be loaded in the preamble, which
% is where this file is \input.
\usepackage{graphicx}
% Single switch for the raster payload: \def\assetdir{assets_light}
% before \begin{document} swaps every embedded raster for its
% print-sized ~600 dpi copy (see figs/make_light_assets.py).  The
% default is the native dump, 965-1016 dpi at the drawn size.
\providecommand{\assetdir}{assets}
\usetikzlibrary{positioning,arrows.meta,fit,backgrounds,calc,shapes.geometric,
                decorations.pathreplacing,decorations.markings,patterns}

% ---------------- palette ----------------
\definecolor{opsred}{HTML}{D55E00}   % Okabe-Ito vermillion  -- training-free ops
\definecolor{trngreen}{HTML}{009E73} % Okabe-Ito bluish green-- trained correctors
\definecolor{datablue}{HTML}{0072B2} % Okabe-Ito blue        -- data / tensors
\definecolor{frozgray}{HTML}{6E6E6E} % neutral               -- frozen base
\definecolor{flowink}{HTML}{262626}  % near black            -- data-flow arrows
\definecolor{warnred}{HTML}{9E2B2B}  % muted brick           -- failure modes
\definecolor{sarink}{HTML}{2B2B33}   % SAR chip background

% ---------------- exact-width canvases ----------------
% \figcanvasd{H} : double column, 18.2 cm = 516 pt = \textwidth  (scale 1.0)
% \figcanvasc{H} : single column,  8.89 cm = 252 pt = \columnwidth
\newcommand{\figcanvasd}[1]{\useasboundingbox (0,0) rectangle (18.2,#1);}
\newcommand{\figcanvasc}[1]{\useasboundingbox (0,0) rectangle (8.89,#1);}

\tikzset{
  % ---- picture-level defaults: ONE font size system for all five figures ----
  fig/.style={font=\footnotesize, line width=0.65pt,
              >={Stealth[length=1.7mm,width=1.35mm]},
              every node/.append style={align=center, inner sep=2.4pt}},
  % ---- blocks: role -> color + redundant border encoding ----
  % Fill tints were raised in v3: at 4-8% the five roles were separated by an
  % RGB distance of only 6.6-21 (out of 441) and all landed on gray 243-248,
  % i.e. the fill channel was consuming ink and returning nothing.  The border
  % weight / dash pattern / badge remain the PRIMARY, grayscale-safe channel.
  ops/.style   ={draw=opsred,   fill=opsred!16,   line width=0.9pt,
                 rounded corners=1.6pt, inner sep=3pt, align=center},
  trn/.style   ={draw=trngreen, fill=trngreen!15, line width=0.9pt,
                 dash pattern=on 1.7pt off 1.1pt,
                 rounded corners=1.6pt, inner sep=3pt, align=center},
  froz/.style  ={draw=frozgray, fill=black!8,     line width=0.7pt,
                 rounded corners=1.6pt, inner sep=3pt, align=center},
  data/.style  ={draw=datablue, fill=datablue!14, line width=0.55pt,
                 rounded corners=1.6pt, inner sep=3pt, align=center},
  plain/.style ={draw=black!45, fill=white,       line width=0.55pt,
                 rounded corners=1.6pt, inner sep=3pt, align=center},
  % v4: radius 1.4 -> 1.6 pt, the one block style that did not share the
  % system's corner radius; weight 0.5 -> 0.6 pt so the failure chips are not
  % the faintest object on the page while being the ones a reader must find.
  % The dotted dash pattern stays as the grayscale-safe channel.
  warn/.style  ={draw=warnred,  fill=warnred!9,   line width=0.6pt,
                 dash pattern=on 1.1pt off 1.1pt,
                 rounded corners=1.6pt, inner sep=3pt, align=center},
  gate/.style  ={draw=opsred,   fill=opsred!16,   line width=0.9pt,
                 diamond, aspect=2.15, inner sep=1pt, align=center},
  % ---- named region (O1): dotted, unfilled, label set outside ----
  grp/.style   ={draw=black!45, densely dotted, line width=0.7pt,
                 rounded corners=2pt, inner sep=0pt},
  grptitle/.style={font=\footnotesize\bfseries, text=black!80, inner sep=1.5pt},
  % ---- step badges (redundant, grayscale-safe stage identity) ----
  badge/.style ={circle, inner sep=0pt, minimum size=3.3mm, text=white,
                 font=\scriptsize\bfseries},
  badgeops/.style={badge, fill=opsred},
  badgetrn/.style={badge, fill=trngreen},
  % ---- typography (O7): exactly three text sizes in the whole figure set ----
  % v3: 7 pt at 55% black was the worst corner of the design space (minimum
  % size AND minimum contrast) -- sub-labels are now 7 pt at 70% black, and
  % anything carrying a QUANTITY is promoted to 8 pt via \qty / \subk.
  ttl/.style   ={font=\footnotesize\bfseries},
  lbl/.style   ={font=\footnotesize},
  sub/.style   ={font=\scriptsize, text=black!70},
  subk/.style  ={font=\footnotesize, text=black!80},
  qty/.style   ={font=\footnotesize, text=black!80},% tensor dims, AP numbers
  tsr/.style ={font=\scriptsize\ttfamily, text=black!70},% B x C x H x W (O5)
  tag/.style   ={font=\scriptsize, inner sep=1.6pt, rounded corners=1pt},
  % panel letters (B11): \footnotesize\bfseries at the panel's top-left
  panl/.style  ={font=\footnotesize\bfseries, text=black!85, inner sep=1pt},
  % ---- arrows (O6): one data-flow style, one auxiliary/optional style ----
  flow/.style  ={-{Stealth[length=1.7mm,width=1.35mm]}, draw=flowink,
                 line width=0.65pt},
  flowt/.style ={-{Stealth[length=1.7mm,width=1.35mm]}, draw=flowink,
                 line width=0.65pt, rounded corners=1.6pt},
  aux/.style   ={-{Stealth[length=1.5mm,width=1.2mm]}, draw=black!45,
                 line width=0.5pt, densely dashed},
  auxw/.style  ={-{Stealth[length=1.5mm,width=1.2mm]}, draw=warnred!75,
                 line width=0.5pt, densely dashed},
  rail/.style  ={draw=black!30, line width=0.5pt},
  brace/.style ={decorate, decoration={brace, amplitude=3.2pt, raise=1pt},
                 draw=black!50, line width=0.5pt},
  % ---- ONE leader style for every magnifier cone in the figure set --------
  % v4: the zoom cones used to be drawn in the generic dashed `aux' style, so
  % a sight line was typographically the same object as an ordinary auxiliary
  % arrow.  A cone is now its own object: it carries the hue of the locator
  % rectangle it starts from, no arrowhead (it is a correspondence, not a
  % flow), and it always runs corner-to-corner from the locator box to the
  % MAGNIFIED RASTER, never to the enclosing panel frame.
  zoomray/.style={draw=opsred!55, line width=0.5pt,
                  dash pattern=on 1.3pt off 1.1pt},
  zoombox/.style={draw=opsred, line width=0.7pt},
}

% =========================================================================
% \hbarrow{x0}{y}{track}{frac}{colour}{weight}{label}
%   One horizontal bar on a printed 0..1 track, with the numeral set at the
%   bar end.  v4: fig1's AP progression, figm_tta (c), figm_fusion's result
%   panel and figm_rerank (b) all drew this same object with four slightly
%   different geometries.  They now share one macro, so the reader learns the
%   idiom once and meets it four times -- a small multiple across figures.
%   The empty track is always drawn, so a short bar still shows its scale.
%
%   DEPRECATED FOR METHOD FIGURES (docs/FIGURE_REDESIGN_SPEC.md S5, items 15
%   and 16, rule H1).  \hbarrow / \hbarrowr / \baxis print an accuracy number
%   inside the figure, which the v6 redraw forbids: fig1_framework, figm_bpr
%   and figm_fusion no longer call them and their numbers live in Tables IV
%   and VI.  The macros are KEPT, not deleted, because the parked single-
%   column data panel of Fig. 5 (figm_rerank, permutation importance) is a
%   data figure, where a measured bar length is evidence rather than a
%   result.  Same status, same reason: warn/, badge/ + badgeops/ + badgetrn/,
%   \hull and \mplane -- no v6 method figure uses any of them.
% =========================================================================
\newcommand{\hbarrow}[7]{%
  \draw[black!15, line width=0.5pt] (#1,#2) -- ({#1+#3},#2);
  \draw[#5, line width=#6] (#1,#2) -- ({#1+(#3)*(#4)},#2);
  % the numeral gets a paper backing so the empty part of the track does not
  % run into it and read as a stray rule after the digits
  \node[qty, anchor=west, inner sep=1.6pt, fill=white]
    at ({#1+(#3)*(#4)},#2) {#7};}

% \hbarrowr{x0}{y}{track}{frac}{colour}{weight}{label}
%   As \hbarrow, but the numeral is set in a fixed column just past the end
%   of the track instead of at the bar tip.  Use it where several groups sit
%   side by side on a SHORT track: it aligns the numerals into columns across
%   the groups, and it stops the leftover empty track of a short bar from
%   printing as a stray rule immediately after the digits -- at 1.25 cm of
%   track a 27.7/75 bar leaves 3 mm of rule past the numeral, which reads as
%   an en dash and can be mistaken for a range.
\newcommand{\hbarrowr}[7]{%
  \draw[black!15, line width=0.5pt] (#1,#2) -- ({#1+#3},#2);
  \draw[#5, line width=#6] (#1,#2) -- ({#1+(#3)*(#4)},#2);
  \node[qty, anchor=west, inner sep=1.6pt] at ({#1+(#3)+0.04},#2) {#7};}

% \baxis{x0}{y}{track}{lo}{hi}{unit} : the tick pair + caption for a \hbarrow
%   group, drawn once per group rather than once per bar.
\newcommand{\baxis}[6]{%
  \draw[black!45, line width=0.5pt] (#1,#2) -- (#1,{#2-0.05})
                                    ({#1+#3},#2) -- ({#1+#3},{#2-0.05});
  \node[sub, anchor=north, inner sep=1pt] at (#1,{#2-0.05}) {#4};
  \node[sub, anchor=north, inner sep=1pt] at ({#1+#3},{#2-0.05}) {#5};
  \node[sub, anchor=west, inner sep=1pt] at ({#1+#3+0.14},{#2-0.16}) {#6};}

% =========================================================================
% GEOMETRIC ISOMORPHISM (v3).  Tensor shape is DRAWN, never only written.
%
%  G1 SPATIAL EXTENT -> DRAWN SIZE, one scale for the whole figure set:
%       1 mm of drawn edge per 40 pixels of tensor extent,
%       i.e. drawn side [cm] = side [px] / 400   ->  \pxc{800} = 2.00 cm.
%       800 -> 20.0 mm | 200 -> 5.00 mm | 100 -> 2.50 mm |  50 -> 1.25 mm
%        25 ->  0.63 mm | 640 -> 16.0 mm | 960 -> 24.0 mm |1088 -> 27.2 mm
%        64 ->  1.60 mm | 128 ->  3.20 mm
%       A stride-4 map is therefore exactly 8x the linear size of a
%       stride-32 map, as it must be.
%  G2 CHANNEL COUNT -> SLAB DEPTH (isometric extrusion, 30 degrees):
%       depth [cm] = 0.04 + 0.055 * log2(C/256)
%       256 -> 0.40 mm | 512 -> 0.95 mm | 1024 -> 1.50 mm | 2048 -> 2.05 mm
%       i.e. +0.55 mm of depth per channel doubling.
%  G3 CATEGORY -> GLYPH SHAPE (readable in grayscale, no label needed):
%       spatial feature map (B,C,H,W) ......... extruded slab   \fmslab
%       token set (N x C, no spatial extent) .. column of bars  \tokencol
%       mask stack (N x H x W) ................ offset planes   \maskstack
%       scalar score .......................... short bar       \scorebar
%       feature vector (d cells) .............. strip of cells  (drawn inline)
%  G4 A LABELLED ZOOM INSET may magnify a sub-scene by an explicit factor z;
%       inside it G1/G2 hold multiplied by z.  Every macro takes z as its
%       last argument (z=1 for the global scale).
% =========================================================================
\newcommand{\pxc}[1]{(#1)/400}     % tensor side in px -> drawn side in cm

% =========================================================================
% G5 REAL DATA (v3).  Wherever a glyph stands for a tensor we actually own,
%     the glyph's FRONT FACE carries that tensor's real raster, dumped by
%     carrier/dump_figure_assets.py with provenance in assets/manifest.json.
%     The vector chrome (frames, arrows, labels, extrusion) stays vector, so
%     type remains selectable and the fonts stay embedded -- the bitmap-panel
%     + vector-chrome architecture Nature's artwork guide prescribes.
%     Resolution guard: every embedded raster must clear 300 dpi at its drawn
%     size and we clear 600 dpi everywhere (chip.png 800 px at 20 mm = 1016
%     dpi).  Feature maps are embedded from the nearest-neighbour _xN variants
%     so the PDF renderer cannot smooth pixels the tensor does not have.
%
% \realpanel{x}{y}{sidepx}{zoom}{file}{framespec} : a real raster, framed and
%     sized by G1 exactly as a schematic glyph of the same tensor would be.
% =========================================================================
\newcommand{\realpanel}[6]{%
  \begin{scope}[shift={(#1,#2)}]
    \pgfmathsetmacro{\rpH}{0.5*#4*\pxc{#3}}
    \pgfmathsetmacro{\rpS}{2*\rpH}
    \node[inner sep=0pt, outer sep=0pt] at (0,0)
      {\includegraphics[width=\rpS cm, height=\rpS cm]{\assetdir/#5}};
    \draw[#6] (-\rpH,-\rpH) rectangle (\rpH,\rpH);
  \end{scope}}

% \realpanelh{...} : the same panel HORIZONTALLY MIRRORED.  Used by the TTA
%     figure so the hflip half of each view pair is a real flip of the real
%     raster, not a redrawn mirror of a schematic.
\newcommand{\realpanelh}[6]{%
  \begin{scope}[shift={(#1,#2)}]
    \pgfmathsetmacro{\rpH}{0.5*#4*\pxc{#3}}
    \pgfmathsetmacro{\rpS}{2*\rpH}
    \node[inner sep=0pt, outer sep=0pt] at (0,0)
      {\reflectbox{\includegraphics[width=\rpS cm, height=\rpS cm]{\assetdir/#5}}};
    \draw[#6] (-\rpH,-\rpH) rectangle (\rpH,\rpH);
  \end{scope}}

% --- \fmslab{x}{y}{sidepx}{C}{colour}{zoom} : a spatial feature map -------
\newcommand{\fmslab}[6]{\fmslabr{#1}{#2}{#3}{#4}{#5}{#6}{}}

% --- \fmslabr{...}{file} : the same slab with a REAL raster on the front
%     face.  Pass an empty 7th argument for the flat-tinted variant.  The
%     extruded top/right faces stay flat tinted, so G2 (depth = channel
%     count) and the single light source (top brightest, right middle,
%     front darkest) are preserved.
\newcommand{\fmslabr}[7]{%
  \begin{scope}[shift={(#1,#2)}]
    \pgfmathsetmacro{\sbH}{0.5*#6*\pxc{#3}}
    \pgfmathsetmacro{\sbD}{#6*(0.04+0.055*ln((#4)/256)/ln(2))}
    \pgfmathsetmacro{\sbX}{0.866*\sbD}
    \pgfmathsetmacro{\sbY}{0.500*\sbD}
    \pgfmathsetmacro{\sbW}{2*\sbH}
    \fill[#5!30] (-\sbH,\sbH) -- (-\sbH+\sbX,\sbH+\sbY)
                 -- (\sbH+\sbX,\sbH+\sbY) -- (\sbH,\sbH) -- cycle;
    \fill[#5!20] (\sbH,\sbH) -- (\sbH+\sbX,\sbH+\sbY)
                 -- (\sbH+\sbX,-\sbH+\sbY) -- (\sbH,-\sbH) -- cycle;
    \fill[#5!11] (-\sbH,-\sbH) rectangle (\sbH,\sbH);
    \edef\fmfile{#7}%
    \ifx\fmfile\empty\else
      \node[inner sep=0pt, outer sep=0pt] at (0,0)
        {\includegraphics[width=\sbW cm, height=\sbW cm]{\assetdir/#7}};
    \fi
    \draw[#5!75, line width=0.5pt]
      (-\sbH,-\sbH) rectangle (\sbH,\sbH)
      (-\sbH,\sbH) -- (-\sbH+\sbX,\sbH+\sbY) -- (\sbH+\sbX,\sbH+\sbY) -- (\sbH,\sbH)
      (\sbH+\sbX,\sbH+\sbY) -- (\sbH+\sbX,-\sbH+\sbY) -- (\sbH,-\sbH);
  \end{scope}}

% --- \tokencol{x}{y}{nbars}{width}{colour} : an N x C token set -----------
%     Deliberately NOT a slab: flat bars, visible gaps, no extrusion, so a
%     query set can never be mistaken for a feature map.
\newcommand{\tokencol}[5]{%
  \begin{scope}[shift={(#1,#2)}]
    \pgfmathsetmacro{\tkP}{0.098}
    \pgfmathsetmacro{\tkT}{0.056}
    \pgfmathsetmacro{\tkW}{0.5*#4}
    \pgfmathtruncatemacro{\tkN}{#3-1}
    \pgfmathsetmacro{\tkY}{0.5*(\tkN*\tkP+\tkT)}
    \foreach \tki in {0,...,\tkN} {
      \pgfmathsetmacro{\tkc}{\tkY-\tki*\tkP}
      \fill[#5!13] (-\tkW,\tkc-\tkT) rectangle (\tkW,\tkc);
      \draw[#5!75, line width=0.5pt] (-\tkW,\tkc-\tkT) rectangle (\tkW,\tkc);
    }
  \end{scope}}

% --- \maskstack{x}{y}{sidepx}{k}{colour}{zoom} : k planes of H x W --------
\newcommand{\maskstack}[6]{\maskstackr{#1}{#2}{#3}{#4}{#5}{#6}{}}

% --- \maskstackr{...}{file} : the same pile with the FRONT plane carrying a
%     real raster; the planes behind stay schematic so the "stack of N"
%     reading survives.  Empty 7th argument = the all-schematic variant.
\newcommand{\maskstackr}[7]{%
  \begin{scope}[shift={(#1,#2)}]
    \pgfmathsetmacro{\msH}{0.5*#6*\pxc{#3}}
    \pgfmathsetmacro{\msO}{0.30*\msH/2.5}
    \pgfmathtruncatemacro{\msN}{#4-1}
    \pgfmathsetmacro{\msS}{\msH/0.92}
    \pgfmathsetmacro{\msW}{2*\msH}
    \edef\msfile{#7}%
    \foreach \msi in {\msN,...,0} {
      \pgfmathsetmacro{\msd}{\msi*\msO}
      \pgfmathtruncatemacro{\msreal}{(\msi==0) ? 1 : 0}
      \edef\mshasimg{0}%
      \ifx\msfile\empty\else\ifnum\msreal=1 \edef\mshasimg{1}\fi\fi
      \begin{scope}[shift={(\msd,\msd)}]
        \fill[white] (-\msH,-\msH) rectangle (\msH,\msH);
        \fill[#5!8]  (-\msH,-\msH) rectangle (\msH,\msH);
        \ifnum\mshasimg=1
          \node[inner sep=0pt, outer sep=0pt] at (0,0)
            {\includegraphics[width=\msW cm, height=\msW cm]{\assetdir/#7}};
        \fi
        \draw[#5!60, line width=0.5pt] (-\msH,-\msH) rectangle (\msH,\msH);
        \ifnum\mshasimg=0
          \begin{scope}[scale=\msS]
            \fill[#5!35] \hull;
            \draw[#5!90, line width=0.5pt] \hull;
          \end{scope}
        \fi
      \end{scope}}
  \end{scope}}

% --- \scorebar{x}{y}{h}{colour} : one scalar, as a short upright bar ------
\newcommand{\scorebar}[4]{%
  \begin{scope}[shift={(#1,#2)}]
    \draw[black!35, line width=0.5pt] (-0.10,0) rectangle (0.10,0.62);
    \fill[#4!55] (-0.10,0) rectangle (0.10,#3);
    \draw[#4!90, line width=0.5pt] (-0.10,0) rectangle (0.10,#3);
  \end{scope}}

% =========================================================================
% Reusable glyphs.  All are TikZ pics: \pic[scale=s] at (x,y) {name};
% =========================================================================

% --- lock: marks any component whose weights we never update -------------
\tikzset{lock/.pic={
  \fill[frozgray] (-0.105,-0.105) rectangle (0.105,0.055);
  \draw[frozgray, line width=0.65pt]
    (-0.060,0.055) -- (-0.060,0.105)
    arc[start angle=180, end angle=0, radius=0.060] -- (0.060,0.055);
  \fill[white] (0,-0.022) circle[radius=0.020];
}}

% --- planes: 3 offset stacked planes = a multi-scale feature pyramid -----
\tikzset{planes/.pic={
  \foreach \i/\o in {2/0.20, 1/0.10, 0/0.00} {
    \fill[datablue!\the\numexpr 10+8*\i\relax]
      (\o-0.24,\o-0.30) rectangle (\o+0.24,\o+0.30);
    \draw[datablue!70, line width=0.5pt]
      (\o-0.24,\o-0.30) rectangle (\o+0.24,\o+0.30);
  }
}}

% --- plate: a single feature plane (thin vertical plate, O4) -------------
\tikzset{plate/.pic={
  \fill[datablue!14] (-0.10,-0.34) rectangle (0.10,0.34);
  \draw[datablue!70, line width=0.5pt] (-0.10,-0.34) rectangle (0.10,0.34);
  \draw[datablue!45, line width=0.5pt] (-0.10,0.20) -- (0.10,0.20);
  \draw[datablue!45, line width=0.5pt] (-0.10,-0.20) -- (0.10,-0.20);
}}

% --- ckpt: an isometric slab = one stored checkpoint / weight vector -----
\tikzset{ckpt/.pic={
  \fill[frozgray!18] (-0.30,-0.13) rectangle (0.24,0.13);
  \fill[frozgray!32] (-0.30,0.13) -- (-0.24,0.21) -- (0.30,0.21) -- (0.24,0.13) -- cycle;
  \fill[frozgray!26] (0.24,-0.13) -- (0.30,-0.05) -- (0.30,0.21) -- (0.24,0.13) -- cycle;
  \draw[frozgray, line width=0.5pt]
    (-0.30,-0.13) rectangle (0.24,0.13);
  \draw[frozgray, line width=0.5pt]
    (-0.30,0.13) -- (-0.24,0.21) -- (0.30,0.21) -- (0.24,0.13)
    (0.30,0.21) -- (0.30,-0.05) -- (0.24,-0.13);
}}

% --- masked: pictogram of an output instance mask (hull on light field) --
\tikzset{maskicon/.pic={
  \fill[datablue!5] (-0.30,-0.30) rectangle (0.30,0.30);
  \draw[datablue!60, line width=0.5pt] (-0.30,-0.30) rectangle (0.30,0.30);
  \begin{scope}[rotate=28]
    \fill[datablue!45] (0,0.05) ellipse[x radius=0.235, y radius=0.062];
    \draw[datablue, line width=0.5pt] (0,0.05) ellipse[x radius=0.235, y radius=0.062];
  \end{scope}
  \begin{scope}[rotate=-12]
    \fill[datablue!45] (0.02,-0.19) ellipse[x radius=0.145, y radius=0.045];
    \draw[datablue, line width=0.5pt] (0.02,-0.19) ellipse[x radius=0.145, y radius=0.045];
  \end{scope}
}}

% =========================================================================
% \sarchip was REMOVED in v3.  It drew two smooth ellipses each ~2.6% of the
% frame on a bokeh-like field; the real HRSID inshore median instance covers
% 0.028% of the chip (13.4 px on 800x800), so the glyph was ~90x too large and
% contradicted the small-target-in-clutter argument the paper makes.  Every
% call site now embeds the REAL chip raster instead (assets/chip.png,
% assets/view_{640,800,960,1088}.png; see the realpanel macro above and
% assets/manifest.json).
% =========================================================================

% =========================================================================
% \hull{}  -- schematic ship-mask contour, ~1.3 cm wide, centred at origin.
% Used by the BPR and fusion figures so every ship outline in the paper is
% literally the same curve.
% =========================================================================
\newcommand{\hull}{plot[smooth cycle, tension=0.72] coordinates
  {(0.72,0.00) (0.34,0.26) (-0.22,0.33) (-0.64,0.17) (-0.62,-0.12)
   (-0.10,-0.30) (0.42,-0.22)}}
.  Nothing else in the figure has to change;
% the block/arrow style names are the same, only their numbers move.
% =========================================================================
% =========================================================================
% tikz_style.tex — shared visual system for the Section III method figures.
% Usage: % =========================================================================
% tikz_style.tex — shared visual system for the Section III method figures.
% Usage: \input{tikz_style} in the preamble of a standalone figure.
%
% -------------------------------------------------------------------------
% DESIGN NOTES — conventions extracted from an accepted TGRS manuscript
% (Gao et al., "BBox-free ... SAR ship instance segmentation", TGRS 2024;
%  local copy: scratchpad/gao_bboxfree_tgrs24.pdf).  Observations taken from
%  its Fig. 2 (overall architecture, p. 5), Fig. 5 (FD module, p. 8) and
%  Fig. 6 (conceptual mask panel, p. 8):
%
%  O1  REGION PARTITION.  The architecture figure is carved into 3-4 named
%      regions by DOTTED rectangles with no fill; the region name is set in
%      bold OUTSIDE the rectangle ("Feature Extractor", "Classification
%      Branch", "Localization Branch", "Center Point Extractor").  Nothing
%      is nested more than two levels deep.
%  O2  COLOR CARRIES MEANING, NOT DECORATION.  Only ~5 hues appear, all as
%      pale fills with black text: light gray = ordinary/standard module,
%      salmon-orange = the module the paper contributes (Feature Decoupling
%      Module, Center Point Extractor), yellow = a distinguished operation
%      (Feature Fusion), light blue / light pink = tensors.  Saturated color
%      is reserved for the few glyphs that must be found first.
%  O3  LEGEND PANEL.  Both Fig. 2 and Fig. 5 carry an explicit legend inside
%      the figure (dotted box top-left in Fig. 2, dashed box right in Fig. 5)
%      that decodes arrow colors and single-letter glyph nodes.  The body of
%      the diagram is kept uncluttered by pushing vocabulary into the legend.
%  O4  TENSOR ICONOGRAPHY.  Feature maps are drawn as thin vertical plates,
%      multi-scale features as 3 offset stacked planes, channel vectors as
%      segmented 3-D bars, the backbone as an isometric slab.  Real SAR
%      chips / heat maps / GT overlays are pasted directly into the flow.
%  O5  TENSOR NAMES SIT BESIDE THE GLYPH, NEVER INSIDE IT, in math italic
%      (F_out, F_mask, F_param, F_H_zip ...), placed above-right of the plate.
%  O6  ARROW DISCIPLINE.  Data flow is a thin gray/black line with a small
%      solid triangular head; routing is orthogonal (Manhattan) with a shared
%      trunk when one source fans out to several targets, so lines never
%      cross.  Semantically special transitions get a THICK COLORED BLOCK
%      arrow instead (2x upsample / 2x downsample / connect), i.e. arrow
%      style is itself a semantic channel.  Long skip paths are routed as
%      rails around the outside of the diagram (Fig. 5 top and bottom).
%  O7  TYPOGRAPHY.  One size for block labels, one smaller size for
%      sub-labels, math only in math mode; no stretched or rotated text
%      except a rare 90-degree axis label.
%  O8  ASPECT RATIO.  Double-column method figures are wide and short:
%      Fig. 2 is 530 x 200 pt (2.65:1), Fig. 5 is 530 x 215 pt (2.5:1).
%      They are laid on a grid with equal gutters between parallel branches.
%  O9  CAPTION DIVISION OF LABOR.  The caption names the parts and decodes
%      abbreviations/colors; it does not re-derive the mechanism, which is
%      left to the body text.  The figure is readable without the caption.
%  O10 CONCEPTUAL PANELS (Fig. 6) use dashed sub-panels with a centered
%      label underneath, a pale blue blob for the true object and a pale red
%      overlay for the error region -- schematic, not photorealistic.
%
% -------------------------------------------------------------------------
% WHAT WE ADOPT / WHERE WE DEVIATE
%  + adopt O1 (dotted named regions), O3 (in-figure legend), O4/O5 (tensor
%    glyphs with the symbol outside), O6 (one flow style + Manhattan
%    routing), O8 (wide-short double-column ratios), O10 (schematic panels).
%  ~ deviate on O2: instead of hue-as-identity we bind hue to ROLE in the
%    contribution -- frozen / training-free operator / trained corrector /
%    data -- because the entire argument of this article is which parts are
%    touched and which are not.
%  ~ deviate on the sans-serif block labels of Fig. 2: we set every figure in
%    the body Times face (as the same paper's Fig. 6 does), so figure text
%    and caption text share one face.
%  + every figure is authored at its EXACT printed width (18.2 cm double
%    column = 516 pt, 8.89 cm single column = 252 pt) and included at
%    scale 1, so \footnotesize is a true 8 pt and \scriptsize a true 7 pt on
%    the page in ALL figures.  Use \figcanvasd / \figcanvasc below.
%
% -------------------------------------------------------------------------
% SEMANTIC PALETTE (Okabe-Ito colorblind-safe base; redundant encoding by
% border weight, dash pattern and badge glyph so it survives grayscale):
%   froz  gray   + lock glyph        frozen, not modified by us  (Stage A)
%   ops   vermil + solid heavy rule  our TRAINING-FREE operator  (rung 1)
%   trn   green  + dashed rule       our OPTIONAL TRAINED corrector (2-3)
%   data  blue   + thin rule         data, tensors, I/O
%   warn  red    + dotted rule, small  failure mode / negative control
% =========================================================================

% Real rasters are embedded on the front faces of the tensor glyphs (see
% \realpanel / \fmslab below).  graphicx must be loaded in the preamble, which
% is where this file is \input.
\usepackage{graphicx}
% Single switch for the raster payload: \def\assetdir{assets_light}
% before \begin{document} swaps every embedded raster for its
% print-sized ~600 dpi copy (see figs/make_light_assets.py).  The
% default is the native dump, 965-1016 dpi at the drawn size.
\providecommand{\assetdir}{assets}
\usetikzlibrary{positioning,arrows.meta,fit,backgrounds,calc,shapes.geometric,
                decorations.pathreplacing,decorations.markings,patterns}

% ---------------- palette ----------------
\definecolor{opsred}{HTML}{D55E00}   % Okabe-Ito vermillion  -- training-free ops
\definecolor{trngreen}{HTML}{009E73} % Okabe-Ito bluish green-- trained correctors
\definecolor{datablue}{HTML}{0072B2} % Okabe-Ito blue        -- data / tensors
\definecolor{frozgray}{HTML}{6E6E6E} % neutral               -- frozen base
\definecolor{flowink}{HTML}{262626}  % near black            -- data-flow arrows
\definecolor{warnred}{HTML}{9E2B2B}  % muted brick           -- failure modes
\definecolor{sarink}{HTML}{2B2B33}   % SAR chip background

% ---------------- exact-width canvases ----------------
% \figcanvasd{H} : double column, 18.2 cm = 516 pt = \textwidth  (scale 1.0)
% \figcanvasc{H} : single column,  8.89 cm = 252 pt = \columnwidth
\newcommand{\figcanvasd}[1]{\useasboundingbox (0,0) rectangle (18.2,#1);}
\newcommand{\figcanvasc}[1]{\useasboundingbox (0,0) rectangle (8.89,#1);}

\tikzset{
  % ---- picture-level defaults: ONE font size system for all five figures ----
  fig/.style={font=\footnotesize, line width=0.65pt,
              >={Stealth[length=1.7mm,width=1.35mm]},
              every node/.append style={align=center, inner sep=2.4pt}},
  % ---- blocks: role -> color + redundant border encoding ----
  % Fill tints were raised in v3: at 4-8% the five roles were separated by an
  % RGB distance of only 6.6-21 (out of 441) and all landed on gray 243-248,
  % i.e. the fill channel was consuming ink and returning nothing.  The border
  % weight / dash pattern / badge remain the PRIMARY, grayscale-safe channel.
  ops/.style   ={draw=opsred,   fill=opsred!16,   line width=0.9pt,
                 rounded corners=1.6pt, inner sep=3pt, align=center},
  trn/.style   ={draw=trngreen, fill=trngreen!15, line width=0.9pt,
                 dash pattern=on 1.7pt off 1.1pt,
                 rounded corners=1.6pt, inner sep=3pt, align=center},
  froz/.style  ={draw=frozgray, fill=black!8,     line width=0.7pt,
                 rounded corners=1.6pt, inner sep=3pt, align=center},
  data/.style  ={draw=datablue, fill=datablue!14, line width=0.55pt,
                 rounded corners=1.6pt, inner sep=3pt, align=center},
  plain/.style ={draw=black!45, fill=white,       line width=0.55pt,
                 rounded corners=1.6pt, inner sep=3pt, align=center},
  % v4: radius 1.4 -> 1.6 pt, the one block style that did not share the
  % system's corner radius; weight 0.5 -> 0.6 pt so the failure chips are not
  % the faintest object on the page while being the ones a reader must find.
  % The dotted dash pattern stays as the grayscale-safe channel.
  warn/.style  ={draw=warnred,  fill=warnred!9,   line width=0.6pt,
                 dash pattern=on 1.1pt off 1.1pt,
                 rounded corners=1.6pt, inner sep=3pt, align=center},
  gate/.style  ={draw=opsred,   fill=opsred!16,   line width=0.9pt,
                 diamond, aspect=2.15, inner sep=1pt, align=center},
  % ---- named region (O1): dotted, unfilled, label set outside ----
  grp/.style   ={draw=black!45, densely dotted, line width=0.7pt,
                 rounded corners=2pt, inner sep=0pt},
  grptitle/.style={font=\footnotesize\bfseries, text=black!80, inner sep=1.5pt},
  % ---- step badges (redundant, grayscale-safe stage identity) ----
  badge/.style ={circle, inner sep=0pt, minimum size=3.3mm, text=white,
                 font=\scriptsize\bfseries},
  badgeops/.style={badge, fill=opsred},
  badgetrn/.style={badge, fill=trngreen},
  % ---- typography (O7): exactly three text sizes in the whole figure set ----
  % v3: 7 pt at 55% black was the worst corner of the design space (minimum
  % size AND minimum contrast) -- sub-labels are now 7 pt at 70% black, and
  % anything carrying a QUANTITY is promoted to 8 pt via \qty / \subk.
  ttl/.style   ={font=\footnotesize\bfseries},
  lbl/.style   ={font=\footnotesize},
  sub/.style   ={font=\scriptsize, text=black!70},
  subk/.style  ={font=\footnotesize, text=black!80},
  qty/.style   ={font=\footnotesize, text=black!80},% tensor dims, AP numbers
  tsr/.style ={font=\scriptsize\ttfamily, text=black!70},% B x C x H x W (O5)
  tag/.style   ={font=\scriptsize, inner sep=1.6pt, rounded corners=1pt},
  % panel letters (B11): \footnotesize\bfseries at the panel's top-left
  panl/.style  ={font=\footnotesize\bfseries, text=black!85, inner sep=1pt},
  % ---- arrows (O6): one data-flow style, one auxiliary/optional style ----
  flow/.style  ={-{Stealth[length=1.7mm,width=1.35mm]}, draw=flowink,
                 line width=0.65pt},
  flowt/.style ={-{Stealth[length=1.7mm,width=1.35mm]}, draw=flowink,
                 line width=0.65pt, rounded corners=1.6pt},
  aux/.style   ={-{Stealth[length=1.5mm,width=1.2mm]}, draw=black!45,
                 line width=0.5pt, densely dashed},
  auxw/.style  ={-{Stealth[length=1.5mm,width=1.2mm]}, draw=warnred!75,
                 line width=0.5pt, densely dashed},
  rail/.style  ={draw=black!30, line width=0.5pt},
  brace/.style ={decorate, decoration={brace, amplitude=3.2pt, raise=1pt},
                 draw=black!50, line width=0.5pt},
  % ---- ONE leader style for every magnifier cone in the figure set --------
  % v4: the zoom cones used to be drawn in the generic dashed `aux' style, so
  % a sight line was typographically the same object as an ordinary auxiliary
  % arrow.  A cone is now its own object: it carries the hue of the locator
  % rectangle it starts from, no arrowhead (it is a correspondence, not a
  % flow), and it always runs corner-to-corner from the locator box to the
  % MAGNIFIED RASTER, never to the enclosing panel frame.
  zoomray/.style={draw=opsred!55, line width=0.5pt,
                  dash pattern=on 1.3pt off 1.1pt},
  zoombox/.style={draw=opsred, line width=0.7pt},
}

% =========================================================================
% \hbarrow{x0}{y}{track}{frac}{colour}{weight}{label}
%   One horizontal bar on a printed 0..1 track, with the numeral set at the
%   bar end.  v4: fig1's AP progression, figm_tta (c), figm_fusion's result
%   panel and figm_rerank (b) all drew this same object with four slightly
%   different geometries.  They now share one macro, so the reader learns the
%   idiom once and meets it four times -- a small multiple across figures.
%   The empty track is always drawn, so a short bar still shows its scale.
%
%   DEPRECATED FOR METHOD FIGURES (docs/FIGURE_REDESIGN_SPEC.md S5, items 15
%   and 16, rule H1).  \hbarrow / \hbarrowr / \baxis print an accuracy number
%   inside the figure, which the v6 redraw forbids: fig1_framework, figm_bpr
%   and figm_fusion no longer call them and their numbers live in Tables IV
%   and VI.  The macros are KEPT, not deleted, because the parked single-
%   column data panel of Fig. 5 (figm_rerank, permutation importance) is a
%   data figure, where a measured bar length is evidence rather than a
%   result.  Same status, same reason: warn/, badge/ + badgeops/ + badgetrn/,
%   \hull and \mplane -- no v6 method figure uses any of them.
% =========================================================================
\newcommand{\hbarrow}[7]{%
  \draw[black!15, line width=0.5pt] (#1,#2) -- ({#1+#3},#2);
  \draw[#5, line width=#6] (#1,#2) -- ({#1+(#3)*(#4)},#2);
  % the numeral gets a paper backing so the empty part of the track does not
  % run into it and read as a stray rule after the digits
  \node[qty, anchor=west, inner sep=1.6pt, fill=white]
    at ({#1+(#3)*(#4)},#2) {#7};}

% \hbarrowr{x0}{y}{track}{frac}{colour}{weight}{label}
%   As \hbarrow, but the numeral is set in a fixed column just past the end
%   of the track instead of at the bar tip.  Use it where several groups sit
%   side by side on a SHORT track: it aligns the numerals into columns across
%   the groups, and it stops the leftover empty track of a short bar from
%   printing as a stray rule immediately after the digits -- at 1.25 cm of
%   track a 27.7/75 bar leaves 3 mm of rule past the numeral, which reads as
%   an en dash and can be mistaken for a range.
\newcommand{\hbarrowr}[7]{%
  \draw[black!15, line width=0.5pt] (#1,#2) -- ({#1+#3},#2);
  \draw[#5, line width=#6] (#1,#2) -- ({#1+(#3)*(#4)},#2);
  \node[qty, anchor=west, inner sep=1.6pt] at ({#1+(#3)+0.04},#2) {#7};}

% \baxis{x0}{y}{track}{lo}{hi}{unit} : the tick pair + caption for a \hbarrow
%   group, drawn once per group rather than once per bar.
\newcommand{\baxis}[6]{%
  \draw[black!45, line width=0.5pt] (#1,#2) -- (#1,{#2-0.05})
                                    ({#1+#3},#2) -- ({#1+#3},{#2-0.05});
  \node[sub, anchor=north, inner sep=1pt] at (#1,{#2-0.05}) {#4};
  \node[sub, anchor=north, inner sep=1pt] at ({#1+#3},{#2-0.05}) {#5};
  \node[sub, anchor=west, inner sep=1pt] at ({#1+#3+0.14},{#2-0.16}) {#6};}

% =========================================================================
% GEOMETRIC ISOMORPHISM (v3).  Tensor shape is DRAWN, never only written.
%
%  G1 SPATIAL EXTENT -> DRAWN SIZE, one scale for the whole figure set:
%       1 mm of drawn edge per 40 pixels of tensor extent,
%       i.e. drawn side [cm] = side [px] / 400   ->  \pxc{800} = 2.00 cm.
%       800 -> 20.0 mm | 200 -> 5.00 mm | 100 -> 2.50 mm |  50 -> 1.25 mm
%        25 ->  0.63 mm | 640 -> 16.0 mm | 960 -> 24.0 mm |1088 -> 27.2 mm
%        64 ->  1.60 mm | 128 ->  3.20 mm
%       A stride-4 map is therefore exactly 8x the linear size of a
%       stride-32 map, as it must be.
%  G2 CHANNEL COUNT -> SLAB DEPTH (isometric extrusion, 30 degrees):
%       depth [cm] = 0.04 + 0.055 * log2(C/256)
%       256 -> 0.40 mm | 512 -> 0.95 mm | 1024 -> 1.50 mm | 2048 -> 2.05 mm
%       i.e. +0.55 mm of depth per channel doubling.
%  G3 CATEGORY -> GLYPH SHAPE (readable in grayscale, no label needed):
%       spatial feature map (B,C,H,W) ......... extruded slab   \fmslab
%       token set (N x C, no spatial extent) .. column of bars  \tokencol
%       mask stack (N x H x W) ................ offset planes   \maskstack
%       scalar score .......................... short bar       \scorebar
%       feature vector (d cells) .............. strip of cells  (drawn inline)
%  G4 A LABELLED ZOOM INSET may magnify a sub-scene by an explicit factor z;
%       inside it G1/G2 hold multiplied by z.  Every macro takes z as its
%       last argument (z=1 for the global scale).
% =========================================================================
\newcommand{\pxc}[1]{(#1)/400}     % tensor side in px -> drawn side in cm

% =========================================================================
% G5 REAL DATA (v3).  Wherever a glyph stands for a tensor we actually own,
%     the glyph's FRONT FACE carries that tensor's real raster, dumped by
%     carrier/dump_figure_assets.py with provenance in assets/manifest.json.
%     The vector chrome (frames, arrows, labels, extrusion) stays vector, so
%     type remains selectable and the fonts stay embedded -- the bitmap-panel
%     + vector-chrome architecture Nature's artwork guide prescribes.
%     Resolution guard: every embedded raster must clear 300 dpi at its drawn
%     size and we clear 600 dpi everywhere (chip.png 800 px at 20 mm = 1016
%     dpi).  Feature maps are embedded from the nearest-neighbour _xN variants
%     so the PDF renderer cannot smooth pixels the tensor does not have.
%
% \realpanel{x}{y}{sidepx}{zoom}{file}{framespec} : a real raster, framed and
%     sized by G1 exactly as a schematic glyph of the same tensor would be.
% =========================================================================
\newcommand{\realpanel}[6]{%
  \begin{scope}[shift={(#1,#2)}]
    \pgfmathsetmacro{\rpH}{0.5*#4*\pxc{#3}}
    \pgfmathsetmacro{\rpS}{2*\rpH}
    \node[inner sep=0pt, outer sep=0pt] at (0,0)
      {\includegraphics[width=\rpS cm, height=\rpS cm]{\assetdir/#5}};
    \draw[#6] (-\rpH,-\rpH) rectangle (\rpH,\rpH);
  \end{scope}}

% \realpanelh{...} : the same panel HORIZONTALLY MIRRORED.  Used by the TTA
%     figure so the hflip half of each view pair is a real flip of the real
%     raster, not a redrawn mirror of a schematic.
\newcommand{\realpanelh}[6]{%
  \begin{scope}[shift={(#1,#2)}]
    \pgfmathsetmacro{\rpH}{0.5*#4*\pxc{#3}}
    \pgfmathsetmacro{\rpS}{2*\rpH}
    \node[inner sep=0pt, outer sep=0pt] at (0,0)
      {\reflectbox{\includegraphics[width=\rpS cm, height=\rpS cm]{\assetdir/#5}}};
    \draw[#6] (-\rpH,-\rpH) rectangle (\rpH,\rpH);
  \end{scope}}

% --- \fmslab{x}{y}{sidepx}{C}{colour}{zoom} : a spatial feature map -------
\newcommand{\fmslab}[6]{\fmslabr{#1}{#2}{#3}{#4}{#5}{#6}{}}

% --- \fmslabr{...}{file} : the same slab with a REAL raster on the front
%     face.  Pass an empty 7th argument for the flat-tinted variant.  The
%     extruded top/right faces stay flat tinted, so G2 (depth = channel
%     count) and the single light source (top brightest, right middle,
%     front darkest) are preserved.
\newcommand{\fmslabr}[7]{%
  \begin{scope}[shift={(#1,#2)}]
    \pgfmathsetmacro{\sbH}{0.5*#6*\pxc{#3}}
    \pgfmathsetmacro{\sbD}{#6*(0.04+0.055*ln((#4)/256)/ln(2))}
    \pgfmathsetmacro{\sbX}{0.866*\sbD}
    \pgfmathsetmacro{\sbY}{0.500*\sbD}
    \pgfmathsetmacro{\sbW}{2*\sbH}
    \fill[#5!30] (-\sbH,\sbH) -- (-\sbH+\sbX,\sbH+\sbY)
                 -- (\sbH+\sbX,\sbH+\sbY) -- (\sbH,\sbH) -- cycle;
    \fill[#5!20] (\sbH,\sbH) -- (\sbH+\sbX,\sbH+\sbY)
                 -- (\sbH+\sbX,-\sbH+\sbY) -- (\sbH,-\sbH) -- cycle;
    \fill[#5!11] (-\sbH,-\sbH) rectangle (\sbH,\sbH);
    \edef\fmfile{#7}%
    \ifx\fmfile\empty\else
      \node[inner sep=0pt, outer sep=0pt] at (0,0)
        {\includegraphics[width=\sbW cm, height=\sbW cm]{\assetdir/#7}};
    \fi
    \draw[#5!75, line width=0.5pt]
      (-\sbH,-\sbH) rectangle (\sbH,\sbH)
      (-\sbH,\sbH) -- (-\sbH+\sbX,\sbH+\sbY) -- (\sbH+\sbX,\sbH+\sbY) -- (\sbH,\sbH)
      (\sbH+\sbX,\sbH+\sbY) -- (\sbH+\sbX,-\sbH+\sbY) -- (\sbH,-\sbH);
  \end{scope}}

% --- \tokencol{x}{y}{nbars}{width}{colour} : an N x C token set -----------
%     Deliberately NOT a slab: flat bars, visible gaps, no extrusion, so a
%     query set can never be mistaken for a feature map.
\newcommand{\tokencol}[5]{%
  \begin{scope}[shift={(#1,#2)}]
    \pgfmathsetmacro{\tkP}{0.098}
    \pgfmathsetmacro{\tkT}{0.056}
    \pgfmathsetmacro{\tkW}{0.5*#4}
    \pgfmathtruncatemacro{\tkN}{#3-1}
    \pgfmathsetmacro{\tkY}{0.5*(\tkN*\tkP+\tkT)}
    \foreach \tki in {0,...,\tkN} {
      \pgfmathsetmacro{\tkc}{\tkY-\tki*\tkP}
      \fill[#5!13] (-\tkW,\tkc-\tkT) rectangle (\tkW,\tkc);
      \draw[#5!75, line width=0.5pt] (-\tkW,\tkc-\tkT) rectangle (\tkW,\tkc);
    }
  \end{scope}}

% --- \maskstack{x}{y}{sidepx}{k}{colour}{zoom} : k planes of H x W --------
\newcommand{\maskstack}[6]{\maskstackr{#1}{#2}{#3}{#4}{#5}{#6}{}}

% --- \maskstackr{...}{file} : the same pile with the FRONT plane carrying a
%     real raster; the planes behind stay schematic so the "stack of N"
%     reading survives.  Empty 7th argument = the all-schematic variant.
\newcommand{\maskstackr}[7]{%
  \begin{scope}[shift={(#1,#2)}]
    \pgfmathsetmacro{\msH}{0.5*#6*\pxc{#3}}
    \pgfmathsetmacro{\msO}{0.30*\msH/2.5}
    \pgfmathtruncatemacro{\msN}{#4-1}
    \pgfmathsetmacro{\msS}{\msH/0.92}
    \pgfmathsetmacro{\msW}{2*\msH}
    \edef\msfile{#7}%
    \foreach \msi in {\msN,...,0} {
      \pgfmathsetmacro{\msd}{\msi*\msO}
      \pgfmathtruncatemacro{\msreal}{(\msi==0) ? 1 : 0}
      \edef\mshasimg{0}%
      \ifx\msfile\empty\else\ifnum\msreal=1 \edef\mshasimg{1}\fi\fi
      \begin{scope}[shift={(\msd,\msd)}]
        \fill[white] (-\msH,-\msH) rectangle (\msH,\msH);
        \fill[#5!8]  (-\msH,-\msH) rectangle (\msH,\msH);
        \ifnum\mshasimg=1
          \node[inner sep=0pt, outer sep=0pt] at (0,0)
            {\includegraphics[width=\msW cm, height=\msW cm]{\assetdir/#7}};
        \fi
        \draw[#5!60, line width=0.5pt] (-\msH,-\msH) rectangle (\msH,\msH);
        \ifnum\mshasimg=0
          \begin{scope}[scale=\msS]
            \fill[#5!35] \hull;
            \draw[#5!90, line width=0.5pt] \hull;
          \end{scope}
        \fi
      \end{scope}}
  \end{scope}}

% --- \scorebar{x}{y}{h}{colour} : one scalar, as a short upright bar ------
\newcommand{\scorebar}[4]{%
  \begin{scope}[shift={(#1,#2)}]
    \draw[black!35, line width=0.5pt] (-0.10,0) rectangle (0.10,0.62);
    \fill[#4!55] (-0.10,0) rectangle (0.10,#3);
    \draw[#4!90, line width=0.5pt] (-0.10,0) rectangle (0.10,#3);
  \end{scope}}

% =========================================================================
% Reusable glyphs.  All are TikZ pics: \pic[scale=s] at (x,y) {name};
% =========================================================================

% --- lock: marks any component whose weights we never update -------------
\tikzset{lock/.pic={
  \fill[frozgray] (-0.105,-0.105) rectangle (0.105,0.055);
  \draw[frozgray, line width=0.65pt]
    (-0.060,0.055) -- (-0.060,0.105)
    arc[start angle=180, end angle=0, radius=0.060] -- (0.060,0.055);
  \fill[white] (0,-0.022) circle[radius=0.020];
}}

% --- planes: 3 offset stacked planes = a multi-scale feature pyramid -----
\tikzset{planes/.pic={
  \foreach \i/\o in {2/0.20, 1/0.10, 0/0.00} {
    \fill[datablue!\the\numexpr 10+8*\i\relax]
      (\o-0.24,\o-0.30) rectangle (\o+0.24,\o+0.30);
    \draw[datablue!70, line width=0.5pt]
      (\o-0.24,\o-0.30) rectangle (\o+0.24,\o+0.30);
  }
}}

% --- plate: a single feature plane (thin vertical plate, O4) -------------
\tikzset{plate/.pic={
  \fill[datablue!14] (-0.10,-0.34) rectangle (0.10,0.34);
  \draw[datablue!70, line width=0.5pt] (-0.10,-0.34) rectangle (0.10,0.34);
  \draw[datablue!45, line width=0.5pt] (-0.10,0.20) -- (0.10,0.20);
  \draw[datablue!45, line width=0.5pt] (-0.10,-0.20) -- (0.10,-0.20);
}}

% --- ckpt: an isometric slab = one stored checkpoint / weight vector -----
\tikzset{ckpt/.pic={
  \fill[frozgray!18] (-0.30,-0.13) rectangle (0.24,0.13);
  \fill[frozgray!32] (-0.30,0.13) -- (-0.24,0.21) -- (0.30,0.21) -- (0.24,0.13) -- cycle;
  \fill[frozgray!26] (0.24,-0.13) -- (0.30,-0.05) -- (0.30,0.21) -- (0.24,0.13) -- cycle;
  \draw[frozgray, line width=0.5pt]
    (-0.30,-0.13) rectangle (0.24,0.13);
  \draw[frozgray, line width=0.5pt]
    (-0.30,0.13) -- (-0.24,0.21) -- (0.30,0.21) -- (0.24,0.13)
    (0.30,0.21) -- (0.30,-0.05) -- (0.24,-0.13);
}}

% --- masked: pictogram of an output instance mask (hull on light field) --
\tikzset{maskicon/.pic={
  \fill[datablue!5] (-0.30,-0.30) rectangle (0.30,0.30);
  \draw[datablue!60, line width=0.5pt] (-0.30,-0.30) rectangle (0.30,0.30);
  \begin{scope}[rotate=28]
    \fill[datablue!45] (0,0.05) ellipse[x radius=0.235, y radius=0.062];
    \draw[datablue, line width=0.5pt] (0,0.05) ellipse[x radius=0.235, y radius=0.062];
  \end{scope}
  \begin{scope}[rotate=-12]
    \fill[datablue!45] (0.02,-0.19) ellipse[x radius=0.145, y radius=0.045];
    \draw[datablue, line width=0.5pt] (0.02,-0.19) ellipse[x radius=0.145, y radius=0.045];
  \end{scope}
}}

% =========================================================================
% \sarchip was REMOVED in v3.  It drew two smooth ellipses each ~2.6% of the
% frame on a bokeh-like field; the real HRSID inshore median instance covers
% 0.028% of the chip (13.4 px on 800x800), so the glyph was ~90x too large and
% contradicted the small-target-in-clutter argument the paper makes.  Every
% call site now embeds the REAL chip raster instead (assets/chip.png,
% assets/view_{640,800,960,1088}.png; see the realpanel macro above and
% assets/manifest.json).
% =========================================================================

% =========================================================================
% \hull{}  -- schematic ship-mask contour, ~1.3 cm wide, centred at origin.
% Used by the BPR and fusion figures so every ship outline in the paper is
% literally the same curve.
% =========================================================================
\newcommand{\hull}{plot[smooth cycle, tension=0.72] coordinates
  {(0.72,0.00) (0.34,0.26) (-0.22,0.33) (-0.64,0.17) (-0.62,-0.12)
   (-0.10,-0.30) (0.42,-0.22)}}
 in the preamble of a standalone figure.
%
% -------------------------------------------------------------------------
% DESIGN NOTES — conventions extracted from an accepted TGRS manuscript
% (Gao et al., "BBox-free ... SAR ship instance segmentation", TGRS 2024;
%  local copy: scratchpad/gao_bboxfree_tgrs24.pdf).  Observations taken from
%  its Fig. 2 (overall architecture, p. 5), Fig. 5 (FD module, p. 8) and
%  Fig. 6 (conceptual mask panel, p. 8):
%
%  O1  REGION PARTITION.  The architecture figure is carved into 3-4 named
%      regions by DOTTED rectangles with no fill; the region name is set in
%      bold OUTSIDE the rectangle ("Feature Extractor", "Classification
%      Branch", "Localization Branch", "Center Point Extractor").  Nothing
%      is nested more than two levels deep.
%  O2  COLOR CARRIES MEANING, NOT DECORATION.  Only ~5 hues appear, all as
%      pale fills with black text: light gray = ordinary/standard module,
%      salmon-orange = the module the paper contributes (Feature Decoupling
%      Module, Center Point Extractor), yellow = a distinguished operation
%      (Feature Fusion), light blue / light pink = tensors.  Saturated color
%      is reserved for the few glyphs that must be found first.
%  O3  LEGEND PANEL.  Both Fig. 2 and Fig. 5 carry an explicit legend inside
%      the figure (dotted box top-left in Fig. 2, dashed box right in Fig. 5)
%      that decodes arrow colors and single-letter glyph nodes.  The body of
%      the diagram is kept uncluttered by pushing vocabulary into the legend.
%  O4  TENSOR ICONOGRAPHY.  Feature maps are drawn as thin vertical plates,
%      multi-scale features as 3 offset stacked planes, channel vectors as
%      segmented 3-D bars, the backbone as an isometric slab.  Real SAR
%      chips / heat maps / GT overlays are pasted directly into the flow.
%  O5  TENSOR NAMES SIT BESIDE THE GLYPH, NEVER INSIDE IT, in math italic
%      (F_out, F_mask, F_param, F_H_zip ...), placed above-right of the plate.
%  O6  ARROW DISCIPLINE.  Data flow is a thin gray/black line with a small
%      solid triangular head; routing is orthogonal (Manhattan) with a shared
%      trunk when one source fans out to several targets, so lines never
%      cross.  Semantically special transitions get a THICK COLORED BLOCK
%      arrow instead (2x upsample / 2x downsample / connect), i.e. arrow
%      style is itself a semantic channel.  Long skip paths are routed as
%      rails around the outside of the diagram (Fig. 5 top and bottom).
%  O7  TYPOGRAPHY.  One size for block labels, one smaller size for
%      sub-labels, math only in math mode; no stretched or rotated text
%      except a rare 90-degree axis label.
%  O8  ASPECT RATIO.  Double-column method figures are wide and short:
%      Fig. 2 is 530 x 200 pt (2.65:1), Fig. 5 is 530 x 215 pt (2.5:1).
%      They are laid on a grid with equal gutters between parallel branches.
%  O9  CAPTION DIVISION OF LABOR.  The caption names the parts and decodes
%      abbreviations/colors; it does not re-derive the mechanism, which is
%      left to the body text.  The figure is readable without the caption.
%  O10 CONCEPTUAL PANELS (Fig. 6) use dashed sub-panels with a centered
%      label underneath, a pale blue blob for the true object and a pale red
%      overlay for the error region -- schematic, not photorealistic.
%
% -------------------------------------------------------------------------
% WHAT WE ADOPT / WHERE WE DEVIATE
%  + adopt O1 (dotted named regions), O3 (in-figure legend), O4/O5 (tensor
%    glyphs with the symbol outside), O6 (one flow style + Manhattan
%    routing), O8 (wide-short double-column ratios), O10 (schematic panels).
%  ~ deviate on O2: instead of hue-as-identity we bind hue to ROLE in the
%    contribution -- frozen / training-free operator / trained corrector /
%    data -- because the entire argument of this article is which parts are
%    touched and which are not.
%  ~ deviate on the sans-serif block labels of Fig. 2: we set every figure in
%    the body Times face (as the same paper's Fig. 6 does), so figure text
%    and caption text share one face.
%  + every figure is authored at its EXACT printed width (18.2 cm double
%    column = 516 pt, 8.89 cm single column = 252 pt) and included at
%    scale 1, so \footnotesize is a true 8 pt and \scriptsize a true 7 pt on
%    the page in ALL figures.  Use \figcanvasd / \figcanvasc below.
%
% -------------------------------------------------------------------------
% SEMANTIC PALETTE (Okabe-Ito colorblind-safe base; redundant encoding by
% border weight, dash pattern and badge glyph so it survives grayscale):
%   froz  gray   + lock glyph        frozen, not modified by us  (Stage A)
%   ops   vermil + solid heavy rule  our TRAINING-FREE operator  (rung 1)
%   trn   green  + dashed rule       our OPTIONAL TRAINED corrector (2-3)
%   data  blue   + thin rule         data, tensors, I/O
%   warn  red    + dotted rule, small  failure mode / negative control
% =========================================================================

% Real rasters are embedded on the front faces of the tensor glyphs (see
% \realpanel / \fmslab below).  graphicx must be loaded in the preamble, which
% is where this file is \input.
\usepackage{graphicx}
% Single switch for the raster payload: \def\assetdir{assets_light}
% before \begin{document} swaps every embedded raster for its
% print-sized ~600 dpi copy (see figs/make_light_assets.py).  The
% default is the native dump, 965-1016 dpi at the drawn size.
\providecommand{\assetdir}{assets}
\usetikzlibrary{positioning,arrows.meta,fit,backgrounds,calc,shapes.geometric,
                decorations.pathreplacing,decorations.markings,patterns}

% ---------------- palette ----------------
\definecolor{opsred}{HTML}{D55E00}   % Okabe-Ito vermillion  -- training-free ops
\definecolor{trngreen}{HTML}{009E73} % Okabe-Ito bluish green-- trained correctors
\definecolor{datablue}{HTML}{0072B2} % Okabe-Ito blue        -- data / tensors
\definecolor{frozgray}{HTML}{6E6E6E} % neutral               -- frozen base
\definecolor{flowink}{HTML}{262626}  % near black            -- data-flow arrows
\definecolor{warnred}{HTML}{9E2B2B}  % muted brick           -- failure modes
\definecolor{sarink}{HTML}{2B2B33}   % SAR chip background

% ---------------- exact-width canvases ----------------
% \figcanvasd{H} : double column, 18.2 cm = 516 pt = \textwidth  (scale 1.0)
% \figcanvasc{H} : single column,  8.89 cm = 252 pt = \columnwidth
\newcommand{\figcanvasd}[1]{\useasboundingbox (0,0) rectangle (18.2,#1);}
\newcommand{\figcanvasc}[1]{\useasboundingbox (0,0) rectangle (8.89,#1);}

\tikzset{
  % ---- picture-level defaults: ONE font size system for all five figures ----
  fig/.style={font=\footnotesize, line width=0.65pt,
              >={Stealth[length=1.7mm,width=1.35mm]},
              every node/.append style={align=center, inner sep=2.4pt}},
  % ---- blocks: role -> color + redundant border encoding ----
  % Fill tints were raised in v3: at 4-8% the five roles were separated by an
  % RGB distance of only 6.6-21 (out of 441) and all landed on gray 243-248,
  % i.e. the fill channel was consuming ink and returning nothing.  The border
  % weight / dash pattern / badge remain the PRIMARY, grayscale-safe channel.
  ops/.style   ={draw=opsred,   fill=opsred!16,   line width=0.9pt,
                 rounded corners=1.6pt, inner sep=3pt, align=center},
  trn/.style   ={draw=trngreen, fill=trngreen!15, line width=0.9pt,
                 dash pattern=on 1.7pt off 1.1pt,
                 rounded corners=1.6pt, inner sep=3pt, align=center},
  froz/.style  ={draw=frozgray, fill=black!8,     line width=0.7pt,
                 rounded corners=1.6pt, inner sep=3pt, align=center},
  data/.style  ={draw=datablue, fill=datablue!14, line width=0.55pt,
                 rounded corners=1.6pt, inner sep=3pt, align=center},
  plain/.style ={draw=black!45, fill=white,       line width=0.55pt,
                 rounded corners=1.6pt, inner sep=3pt, align=center},
  % v4: radius 1.4 -> 1.6 pt, the one block style that did not share the
  % system's corner radius; weight 0.5 -> 0.6 pt so the failure chips are not
  % the faintest object on the page while being the ones a reader must find.
  % The dotted dash pattern stays as the grayscale-safe channel.
  warn/.style  ={draw=warnred,  fill=warnred!9,   line width=0.6pt,
                 dash pattern=on 1.1pt off 1.1pt,
                 rounded corners=1.6pt, inner sep=3pt, align=center},
  gate/.style  ={draw=opsred,   fill=opsred!16,   line width=0.9pt,
                 diamond, aspect=2.15, inner sep=1pt, align=center},
  % ---- named region (O1): dotted, unfilled, label set outside ----
  grp/.style   ={draw=black!45, densely dotted, line width=0.7pt,
                 rounded corners=2pt, inner sep=0pt},
  grptitle/.style={font=\footnotesize\bfseries, text=black!80, inner sep=1.5pt},
  % ---- step badges (redundant, grayscale-safe stage identity) ----
  badge/.style ={circle, inner sep=0pt, minimum size=3.3mm, text=white,
                 font=\scriptsize\bfseries},
  badgeops/.style={badge, fill=opsred},
  badgetrn/.style={badge, fill=trngreen},
  % ---- typography (O7): exactly three text sizes in the whole figure set ----
  % v3: 7 pt at 55% black was the worst corner of the design space (minimum
  % size AND minimum contrast) -- sub-labels are now 7 pt at 70% black, and
  % anything carrying a QUANTITY is promoted to 8 pt via \qty / \subk.
  ttl/.style   ={font=\footnotesize\bfseries},
  lbl/.style   ={font=\footnotesize},
  sub/.style   ={font=\scriptsize, text=black!70},
  subk/.style  ={font=\footnotesize, text=black!80},
  qty/.style   ={font=\footnotesize, text=black!80},% tensor dims, AP numbers
  tsr/.style ={font=\scriptsize\ttfamily, text=black!70},% B x C x H x W (O5)
  tag/.style   ={font=\scriptsize, inner sep=1.6pt, rounded corners=1pt},
  % panel letters (B11): \footnotesize\bfseries at the panel's top-left
  panl/.style  ={font=\footnotesize\bfseries, text=black!85, inner sep=1pt},
  % ---- arrows (O6): one data-flow style, one auxiliary/optional style ----
  flow/.style  ={-{Stealth[length=1.7mm,width=1.35mm]}, draw=flowink,
                 line width=0.65pt},
  flowt/.style ={-{Stealth[length=1.7mm,width=1.35mm]}, draw=flowink,
                 line width=0.65pt, rounded corners=1.6pt},
  aux/.style   ={-{Stealth[length=1.5mm,width=1.2mm]}, draw=black!45,
                 line width=0.5pt, densely dashed},
  auxw/.style  ={-{Stealth[length=1.5mm,width=1.2mm]}, draw=warnred!75,
                 line width=0.5pt, densely dashed},
  rail/.style  ={draw=black!30, line width=0.5pt},
  brace/.style ={decorate, decoration={brace, amplitude=3.2pt, raise=1pt},
                 draw=black!50, line width=0.5pt},
  % ---- ONE leader style for every magnifier cone in the figure set --------
  % v4: the zoom cones used to be drawn in the generic dashed `aux' style, so
  % a sight line was typographically the same object as an ordinary auxiliary
  % arrow.  A cone is now its own object: it carries the hue of the locator
  % rectangle it starts from, no arrowhead (it is a correspondence, not a
  % flow), and it always runs corner-to-corner from the locator box to the
  % MAGNIFIED RASTER, never to the enclosing panel frame.
  zoomray/.style={draw=opsred!55, line width=0.5pt,
                  dash pattern=on 1.3pt off 1.1pt},
  zoombox/.style={draw=opsred, line width=0.7pt},
}

% =========================================================================
% \hbarrow{x0}{y}{track}{frac}{colour}{weight}{label}
%   One horizontal bar on a printed 0..1 track, with the numeral set at the
%   bar end.  v4: fig1's AP progression, figm_tta (c), figm_fusion's result
%   panel and figm_rerank (b) all drew this same object with four slightly
%   different geometries.  They now share one macro, so the reader learns the
%   idiom once and meets it four times -- a small multiple across figures.
%   The empty track is always drawn, so a short bar still shows its scale.
%
%   DEPRECATED FOR METHOD FIGURES (docs/FIGURE_REDESIGN_SPEC.md S5, items 15
%   and 16, rule H1).  \hbarrow / \hbarrowr / \baxis print an accuracy number
%   inside the figure, which the v6 redraw forbids: fig1_framework, figm_bpr
%   and figm_fusion no longer call them and their numbers live in Tables IV
%   and VI.  The macros are KEPT, not deleted, because the parked single-
%   column data panel of Fig. 5 (figm_rerank, permutation importance) is a
%   data figure, where a measured bar length is evidence rather than a
%   result.  Same status, same reason: warn/, badge/ + badgeops/ + badgetrn/,
%   \hull and \mplane -- no v6 method figure uses any of them.
% =========================================================================
\newcommand{\hbarrow}[7]{%
  \draw[black!15, line width=0.5pt] (#1,#2) -- ({#1+#3},#2);
  \draw[#5, line width=#6] (#1,#2) -- ({#1+(#3)*(#4)},#2);
  % the numeral gets a paper backing so the empty part of the track does not
  % run into it and read as a stray rule after the digits
  \node[qty, anchor=west, inner sep=1.6pt, fill=white]
    at ({#1+(#3)*(#4)},#2) {#7};}

% \hbarrowr{x0}{y}{track}{frac}{colour}{weight}{label}
%   As \hbarrow, but the numeral is set in a fixed column just past the end
%   of the track instead of at the bar tip.  Use it where several groups sit
%   side by side on a SHORT track: it aligns the numerals into columns across
%   the groups, and it stops the leftover empty track of a short bar from
%   printing as a stray rule immediately after the digits -- at 1.25 cm of
%   track a 27.7/75 bar leaves 3 mm of rule past the numeral, which reads as
%   an en dash and can be mistaken for a range.
\newcommand{\hbarrowr}[7]{%
  \draw[black!15, line width=0.5pt] (#1,#2) -- ({#1+#3},#2);
  \draw[#5, line width=#6] (#1,#2) -- ({#1+(#3)*(#4)},#2);
  \node[qty, anchor=west, inner sep=1.6pt] at ({#1+(#3)+0.04},#2) {#7};}

% \baxis{x0}{y}{track}{lo}{hi}{unit} : the tick pair + caption for a \hbarrow
%   group, drawn once per group rather than once per bar.
\newcommand{\baxis}[6]{%
  \draw[black!45, line width=0.5pt] (#1,#2) -- (#1,{#2-0.05})
                                    ({#1+#3},#2) -- ({#1+#3},{#2-0.05});
  \node[sub, anchor=north, inner sep=1pt] at (#1,{#2-0.05}) {#4};
  \node[sub, anchor=north, inner sep=1pt] at ({#1+#3},{#2-0.05}) {#5};
  \node[sub, anchor=west, inner sep=1pt] at ({#1+#3+0.14},{#2-0.16}) {#6};}

% =========================================================================
% GEOMETRIC ISOMORPHISM (v3).  Tensor shape is DRAWN, never only written.
%
%  G1 SPATIAL EXTENT -> DRAWN SIZE, one scale for the whole figure set:
%       1 mm of drawn edge per 40 pixels of tensor extent,
%       i.e. drawn side [cm] = side [px] / 400   ->  \pxc{800} = 2.00 cm.
%       800 -> 20.0 mm | 200 -> 5.00 mm | 100 -> 2.50 mm |  50 -> 1.25 mm
%        25 ->  0.63 mm | 640 -> 16.0 mm | 960 -> 24.0 mm |1088 -> 27.2 mm
%        64 ->  1.60 mm | 128 ->  3.20 mm
%       A stride-4 map is therefore exactly 8x the linear size of a
%       stride-32 map, as it must be.
%  G2 CHANNEL COUNT -> SLAB DEPTH (isometric extrusion, 30 degrees):
%       depth [cm] = 0.04 + 0.055 * log2(C/256)
%       256 -> 0.40 mm | 512 -> 0.95 mm | 1024 -> 1.50 mm | 2048 -> 2.05 mm
%       i.e. +0.55 mm of depth per channel doubling.
%  G3 CATEGORY -> GLYPH SHAPE (readable in grayscale, no label needed):
%       spatial feature map (B,C,H,W) ......... extruded slab   \fmslab
%       token set (N x C, no spatial extent) .. column of bars  \tokencol
%       mask stack (N x H x W) ................ offset planes   \maskstack
%       scalar score .......................... short bar       \scorebar
%       feature vector (d cells) .............. strip of cells  (drawn inline)
%  G4 A LABELLED ZOOM INSET may magnify a sub-scene by an explicit factor z;
%       inside it G1/G2 hold multiplied by z.  Every macro takes z as its
%       last argument (z=1 for the global scale).
% =========================================================================
\newcommand{\pxc}[1]{(#1)/400}     % tensor side in px -> drawn side in cm

% =========================================================================
% G5 REAL DATA (v3).  Wherever a glyph stands for a tensor we actually own,
%     the glyph's FRONT FACE carries that tensor's real raster, dumped by
%     carrier/dump_figure_assets.py with provenance in assets/manifest.json.
%     The vector chrome (frames, arrows, labels, extrusion) stays vector, so
%     type remains selectable and the fonts stay embedded -- the bitmap-panel
%     + vector-chrome architecture Nature's artwork guide prescribes.
%     Resolution guard: every embedded raster must clear 300 dpi at its drawn
%     size and we clear 600 dpi everywhere (chip.png 800 px at 20 mm = 1016
%     dpi).  Feature maps are embedded from the nearest-neighbour _xN variants
%     so the PDF renderer cannot smooth pixels the tensor does not have.
%
% \realpanel{x}{y}{sidepx}{zoom}{file}{framespec} : a real raster, framed and
%     sized by G1 exactly as a schematic glyph of the same tensor would be.
% =========================================================================
\newcommand{\realpanel}[6]{%
  \begin{scope}[shift={(#1,#2)}]
    \pgfmathsetmacro{\rpH}{0.5*#4*\pxc{#3}}
    \pgfmathsetmacro{\rpS}{2*\rpH}
    \node[inner sep=0pt, outer sep=0pt] at (0,0)
      {\includegraphics[width=\rpS cm, height=\rpS cm]{\assetdir/#5}};
    \draw[#6] (-\rpH,-\rpH) rectangle (\rpH,\rpH);
  \end{scope}}

% \realpanelh{...} : the same panel HORIZONTALLY MIRRORED.  Used by the TTA
%     figure so the hflip half of each view pair is a real flip of the real
%     raster, not a redrawn mirror of a schematic.
\newcommand{\realpanelh}[6]{%
  \begin{scope}[shift={(#1,#2)}]
    \pgfmathsetmacro{\rpH}{0.5*#4*\pxc{#3}}
    \pgfmathsetmacro{\rpS}{2*\rpH}
    \node[inner sep=0pt, outer sep=0pt] at (0,0)
      {\reflectbox{\includegraphics[width=\rpS cm, height=\rpS cm]{\assetdir/#5}}};
    \draw[#6] (-\rpH,-\rpH) rectangle (\rpH,\rpH);
  \end{scope}}

% --- \fmslab{x}{y}{sidepx}{C}{colour}{zoom} : a spatial feature map -------
\newcommand{\fmslab}[6]{\fmslabr{#1}{#2}{#3}{#4}{#5}{#6}{}}

% --- \fmslabr{...}{file} : the same slab with a REAL raster on the front
%     face.  Pass an empty 7th argument for the flat-tinted variant.  The
%     extruded top/right faces stay flat tinted, so G2 (depth = channel
%     count) and the single light source (top brightest, right middle,
%     front darkest) are preserved.
\newcommand{\fmslabr}[7]{%
  \begin{scope}[shift={(#1,#2)}]
    \pgfmathsetmacro{\sbH}{0.5*#6*\pxc{#3}}
    \pgfmathsetmacro{\sbD}{#6*(0.04+0.055*ln((#4)/256)/ln(2))}
    \pgfmathsetmacro{\sbX}{0.866*\sbD}
    \pgfmathsetmacro{\sbY}{0.500*\sbD}
    \pgfmathsetmacro{\sbW}{2*\sbH}
    \fill[#5!30] (-\sbH,\sbH) -- (-\sbH+\sbX,\sbH+\sbY)
                 -- (\sbH+\sbX,\sbH+\sbY) -- (\sbH,\sbH) -- cycle;
    \fill[#5!20] (\sbH,\sbH) -- (\sbH+\sbX,\sbH+\sbY)
                 -- (\sbH+\sbX,-\sbH+\sbY) -- (\sbH,-\sbH) -- cycle;
    \fill[#5!11] (-\sbH,-\sbH) rectangle (\sbH,\sbH);
    \edef\fmfile{#7}%
    \ifx\fmfile\empty\else
      \node[inner sep=0pt, outer sep=0pt] at (0,0)
        {\includegraphics[width=\sbW cm, height=\sbW cm]{\assetdir/#7}};
    \fi
    \draw[#5!75, line width=0.5pt]
      (-\sbH,-\sbH) rectangle (\sbH,\sbH)
      (-\sbH,\sbH) -- (-\sbH+\sbX,\sbH+\sbY) -- (\sbH+\sbX,\sbH+\sbY) -- (\sbH,\sbH)
      (\sbH+\sbX,\sbH+\sbY) -- (\sbH+\sbX,-\sbH+\sbY) -- (\sbH,-\sbH);
  \end{scope}}

% --- \tokencol{x}{y}{nbars}{width}{colour} : an N x C token set -----------
%     Deliberately NOT a slab: flat bars, visible gaps, no extrusion, so a
%     query set can never be mistaken for a feature map.
\newcommand{\tokencol}[5]{%
  \begin{scope}[shift={(#1,#2)}]
    \pgfmathsetmacro{\tkP}{0.098}
    \pgfmathsetmacro{\tkT}{0.056}
    \pgfmathsetmacro{\tkW}{0.5*#4}
    \pgfmathtruncatemacro{\tkN}{#3-1}
    \pgfmathsetmacro{\tkY}{0.5*(\tkN*\tkP+\tkT)}
    \foreach \tki in {0,...,\tkN} {
      \pgfmathsetmacro{\tkc}{\tkY-\tki*\tkP}
      \fill[#5!13] (-\tkW,\tkc-\tkT) rectangle (\tkW,\tkc);
      \draw[#5!75, line width=0.5pt] (-\tkW,\tkc-\tkT) rectangle (\tkW,\tkc);
    }
  \end{scope}}

% --- \maskstack{x}{y}{sidepx}{k}{colour}{zoom} : k planes of H x W --------
\newcommand{\maskstack}[6]{\maskstackr{#1}{#2}{#3}{#4}{#5}{#6}{}}

% --- \maskstackr{...}{file} : the same pile with the FRONT plane carrying a
%     real raster; the planes behind stay schematic so the "stack of N"
%     reading survives.  Empty 7th argument = the all-schematic variant.
\newcommand{\maskstackr}[7]{%
  \begin{scope}[shift={(#1,#2)}]
    \pgfmathsetmacro{\msH}{0.5*#6*\pxc{#3}}
    \pgfmathsetmacro{\msO}{0.30*\msH/2.5}
    \pgfmathtruncatemacro{\msN}{#4-1}
    \pgfmathsetmacro{\msS}{\msH/0.92}
    \pgfmathsetmacro{\msW}{2*\msH}
    \edef\msfile{#7}%
    \foreach \msi in {\msN,...,0} {
      \pgfmathsetmacro{\msd}{\msi*\msO}
      \pgfmathtruncatemacro{\msreal}{(\msi==0) ? 1 : 0}
      \edef\mshasimg{0}%
      \ifx\msfile\empty\else\ifnum\msreal=1 \edef\mshasimg{1}\fi\fi
      \begin{scope}[shift={(\msd,\msd)}]
        \fill[white] (-\msH,-\msH) rectangle (\msH,\msH);
        \fill[#5!8]  (-\msH,-\msH) rectangle (\msH,\msH);
        \ifnum\mshasimg=1
          \node[inner sep=0pt, outer sep=0pt] at (0,0)
            {\includegraphics[width=\msW cm, height=\msW cm]{\assetdir/#7}};
        \fi
        \draw[#5!60, line width=0.5pt] (-\msH,-\msH) rectangle (\msH,\msH);
        \ifnum\mshasimg=0
          \begin{scope}[scale=\msS]
            \fill[#5!35] \hull;
            \draw[#5!90, line width=0.5pt] \hull;
          \end{scope}
        \fi
      \end{scope}}
  \end{scope}}

% --- \scorebar{x}{y}{h}{colour} : one scalar, as a short upright bar ------
\newcommand{\scorebar}[4]{%
  \begin{scope}[shift={(#1,#2)}]
    \draw[black!35, line width=0.5pt] (-0.10,0) rectangle (0.10,0.62);
    \fill[#4!55] (-0.10,0) rectangle (0.10,#3);
    \draw[#4!90, line width=0.5pt] (-0.10,0) rectangle (0.10,#3);
  \end{scope}}

% =========================================================================
% Reusable glyphs.  All are TikZ pics: \pic[scale=s] at (x,y) {name};
% =========================================================================

% --- lock: marks any component whose weights we never update -------------
\tikzset{lock/.pic={
  \fill[frozgray] (-0.105,-0.105) rectangle (0.105,0.055);
  \draw[frozgray, line width=0.65pt]
    (-0.060,0.055) -- (-0.060,0.105)
    arc[start angle=180, end angle=0, radius=0.060] -- (0.060,0.055);
  \fill[white] (0,-0.022) circle[radius=0.020];
}}

% --- planes: 3 offset stacked planes = a multi-scale feature pyramid -----
\tikzset{planes/.pic={
  \foreach \i/\o in {2/0.20, 1/0.10, 0/0.00} {
    \fill[datablue!\the\numexpr 10+8*\i\relax]
      (\o-0.24,\o-0.30) rectangle (\o+0.24,\o+0.30);
    \draw[datablue!70, line width=0.5pt]
      (\o-0.24,\o-0.30) rectangle (\o+0.24,\o+0.30);
  }
}}

% --- plate: a single feature plane (thin vertical plate, O4) -------------
\tikzset{plate/.pic={
  \fill[datablue!14] (-0.10,-0.34) rectangle (0.10,0.34);
  \draw[datablue!70, line width=0.5pt] (-0.10,-0.34) rectangle (0.10,0.34);
  \draw[datablue!45, line width=0.5pt] (-0.10,0.20) -- (0.10,0.20);
  \draw[datablue!45, line width=0.5pt] (-0.10,-0.20) -- (0.10,-0.20);
}}

% --- ckpt: an isometric slab = one stored checkpoint / weight vector -----
\tikzset{ckpt/.pic={
  \fill[frozgray!18] (-0.30,-0.13) rectangle (0.24,0.13);
  \fill[frozgray!32] (-0.30,0.13) -- (-0.24,0.21) -- (0.30,0.21) -- (0.24,0.13) -- cycle;
  \fill[frozgray!26] (0.24,-0.13) -- (0.30,-0.05) -- (0.30,0.21) -- (0.24,0.13) -- cycle;
  \draw[frozgray, line width=0.5pt]
    (-0.30,-0.13) rectangle (0.24,0.13);
  \draw[frozgray, line width=0.5pt]
    (-0.30,0.13) -- (-0.24,0.21) -- (0.30,0.21) -- (0.24,0.13)
    (0.30,0.21) -- (0.30,-0.05) -- (0.24,-0.13);
}}

% --- masked: pictogram of an output instance mask (hull on light field) --
\tikzset{maskicon/.pic={
  \fill[datablue!5] (-0.30,-0.30) rectangle (0.30,0.30);
  \draw[datablue!60, line width=0.5pt] (-0.30,-0.30) rectangle (0.30,0.30);
  \begin{scope}[rotate=28]
    \fill[datablue!45] (0,0.05) ellipse[x radius=0.235, y radius=0.062];
    \draw[datablue, line width=0.5pt] (0,0.05) ellipse[x radius=0.235, y radius=0.062];
  \end{scope}
  \begin{scope}[rotate=-12]
    \fill[datablue!45] (0.02,-0.19) ellipse[x radius=0.145, y radius=0.045];
    \draw[datablue, line width=0.5pt] (0.02,-0.19) ellipse[x radius=0.145, y radius=0.045];
  \end{scope}
}}

% =========================================================================
% \sarchip was REMOVED in v3.  It drew two smooth ellipses each ~2.6% of the
% frame on a bokeh-like field; the real HRSID inshore median instance covers
% 0.028% of the chip (13.4 px on 800x800), so the glyph was ~90x too large and
% contradicted the small-target-in-clutter argument the paper makes.  Every
% call site now embeds the REAL chip raster instead (assets/chip.png,
% assets/view_{640,800,960,1088}.png; see the realpanel macro above and
% assets/manifest.json).
% =========================================================================

% =========================================================================
% \hull{}  -- schematic ship-mask contour, ~1.3 cm wide, centred at origin.
% Used by the BPR and fusion figures so every ship outline in the paper is
% literally the same curve.
% =========================================================================
\newcommand{\hull}{plot[smooth cycle, tension=0.72] coordinates
  {(0.72,0.00) (0.34,0.26) (-0.22,0.33) (-0.64,0.17) (-0.62,-0.12)
   (-0.10,-0.30) (0.42,-0.22)}}


% ---------------- T1: type ----------------
% Script sizes for the two text sizes used in figures.  Must be in the
% preamble; standalone + \input keeps it there.
\DeclareMathSizes{8}{8}{6}{6}
\DeclareMathSizes{7}{7}{6}{6}

% ---------------- T2/T3/T4: block, frame and arrow chrome ----------------
\tikzset{
  fig/.style={font=\footnotesize, line width=0.5pt,
              >={Stealth[length=1.6mm,width=1.2mm]},
              every node/.append style={align=center, inner sep=2.4pt}},
  % blocks: one radius, two weights; hue = role, dash = redundant channel
  blk/.style   ={rounded corners=1.5pt, inner sep=3pt, align=center},
  ops/.style   ={blk, draw=opsred,   fill=opsred!14,   line width=0.8pt},
  trn/.style   ={blk, draw=trngreen, fill=trngreen!14, line width=0.8pt,
                 dash pattern=on 2pt off 1.2pt},
  froz/.style  ={blk, draw=frozgray, fill=black!7,     line width=0.5pt},
  data/.style  ={blk, draw=datablue, fill=datablue!12, line width=0.5pt},
  plain/.style ={blk, draw=black!45, fill=white,       line width=0.5pt},
  warn/.style  ={blk, draw=warnred,  fill=warnred!8,   line width=0.5pt,
                 dash pattern=on 1pt off 1pt},
  % a carrier = established practice the operator rides on (fig1 v5 local
  % style, promoted here so every figure spells it the same way)
  carrier/.style={blk, draw=opsred, fill=opsred!5, line width=0.5pt,
                 dash pattern=on 2.6pt off 1.2pt},
  gate/.style  ={draw=opsred, fill=opsred!14, line width=0.8pt,
                 diamond, aspect=2.15, inner sep=1pt, align=center},
  grp/.style   ={draw=black!40, densely dotted, line width=0.5pt,
                 rounded corners=2pt, inner sep=0pt},
  grptitle/.style={font=\footnotesize\bfseries, text=black!80, inner sep=1.5pt},
  % type: exactly two sizes; sub-labels at 75% black (contrast >= 4.5:1)
  ttl/.style   ={font=\footnotesize\bfseries},
  lbl/.style   ={font=\footnotesize},
  sub/.style   ={font=\scriptsize, text=black!75},
  subk/.style  ={font=\footnotesize, text=black!80},
  qty/.style   ={font=\footnotesize, text=black!80},
  tsr/.style   ={font=\scriptsize\ttfamily, text=black!75},
  panl/.style  ={font=\footnotesize\bfseries, text=black!85, inner sep=1pt},
  badgecar/.style={badge, draw=opsred, fill=white, text=opsred, line width=0.5pt},
  % arrows: one head, 0.5 pt flow, 0.8 pt emphasis, dashed auxiliary
  flow/.style  ={-{Stealth[length=1.6mm,width=1.2mm]}, draw=flowink,
                 line width=0.5pt},
  flowt/.style ={flow, rounded corners=1.5pt},
  flowh/.style ={-{Stealth[length=1.9mm,width=1.4mm]}, draw=flowink,
                 line width=0.8pt},
  aux/.style   ={-{Stealth[length=1.4mm,width=1.1mm]}, draw=black!45,
                 line width=0.5pt, densely dashed},
  auxw/.style  ={aux, draw=warnred!75},
  rail/.style  ={draw=black!30, line width=0.5pt},
  brace/.style ={decorate, decoration={brace, amplitude=3pt, raise=1pt},
                 draw=black!50, line width=0.5pt},
  zoomray/.style={draw=opsred!55, line width=0.5pt, dash pattern=on 1.3pt off 1.1pt},
  zoombox/.style={draw=opsred, line width=0.5pt},
}

% ---------------- T7: rules inside math must obey the stroke floor --------
% LaTeX draws a \frac bar at the math extension font's default_rule_thickness,
% which SCALES WITH THE TYPE SIZE: 0.4 pt at 10 pt becomes 0.279 pt at the 7 pt
% \scriptsize used for figure sub-labels -- the single sub-0.5 pt stroke that
% verify_fig.py reports in fig1_framework.pdf (a \frac{1}{3} at (166.7,38.3)).
% \genfrac states the rule thickness explicitly, so it survives the floor.
% Measured: \frac -> 0.279 PDF pt, \safefrac -> 0.498 PDF pt (= 0.5 TeX pt).
% The alternative with no rule at all is to write $(1/3)\sum_s\theta_s$.
\newcommand{\safefrac}[2]{\genfrac{}{}{0.5pt}{}{#1}{#2}}

% ---------------- T8: no Type-3 bitmap fonts ----------------------------
% \ttfamily under [T1]{fontenc} with no Type-1 EC typewriter installed makes
% pdflatex fall back to the ec-tt BITMAP font (ectt0800.pk), which is embedded
% as a TYPE 3 font.  Measured on figm_tta.pdf: "F74  Type 3  Custom  emb=yes
% sub=no uni=no" (pdffonts), and figm_tta.log line 655 shows \ectt0800.pk.
% IEEE asks authors to "either convert all text to outlines or embed the
% fonts" and PDF eXpress flags Type 3; the tsr/ style below sets \ttfamily, so
% the substitution must be global.  lmtt is Latin Modern Typewriter: a real
% Type-1 T1 face.  A/B measured: without -> "F39 Type 3", with -> "LMMono8
% Type 1".  Only the typewriter family moves; mathptmx keeps Times for text.
\renewcommand{\ttdefault}{lmtt}

% ---------------- T6: legend row ----------------
% \legendswatch{x}{y}{style}{text} : a 4.2 x 3 mm swatch in <style> with its
% 7 pt label to the right.  Lay the swatches on one baseline y.
\newcommand{\legendswatch}[4]{%
  \draw[#3] (#1,{#2-0.15}) rectangle ({#1+0.42},{#2+0.15});
  \node[sub, anchor=west] at ({#1+0.54},#2) {#4};}

% =========================================================================
% T5: 3D tensor pics (PlotNeuralNet idiom, geometry from the tensor shape)
%
%   \pic (c2) at (x,y) {tensor={px=200, ch=256, img=feat_c2_x4.png,
%                               caption={$256@200^2$}}};
%   \draw[flow] (x-east) -- (c2-west);
%
% keys (all optional; defaults in brackets)
%   px       spatial side in pixels           -> drawn side = zoom*px/400 cm  [200]
%   ch       channel count                    -> depth = 0.04+0.055*log2(ch/256) cm,
%                                                floored at 0.04 cm           [256]
%   depth    explicit depth in cm, overrides ch (hand-sized blocks, UNet)   [ ]
%   zoom     magnification inside a zoom inset (G4)                          [1]
%   color    any xcolor expression; tints are derived from it               [datablue]
%   img      raster placed on the FRONT face, path relative to \assetdir     [ ]
%   hidden   1 = draw the three hidden edges dashed (the PlotNeuralNet look) [0]
%   caption  7 pt text set under the front face                             [ ]
% anchors (prefix = pic name)
%   -west -east          mid-points of the front face's left / right edge (arrows)
%   -westc -eastc        the same at half depth (for rails behind the block)
%   -north -south        mid-points of the front face's top / bottom edge
%   -front -back         centre of the front / back face
%   -nw -ne -se -sw      front-face corners;  -topne  far top-right corner
% Projection: front face in the page plane (so \includegraphics needs no
% transform), depth receding at 30 degrees to the upper right, identical to
% \fmslab, so the two glyph families can share one figure.
% =========================================================================
\tikzset{
  tensor/.pic={
    \tikzset{/tensor/.cd,#1}
    \colorlet{tnc}{\tncolor}
    \pgfmathsetmacro{\tnW}{\tnzoom*(\tnpx)/400}
    \pgfmathsetmacro{\tnH}{0.5*\tnW}
    \edef\tnDepthUser{\tndepth}
    \ifx\tnDepthUser\empty
      \pgfmathsetmacro{\tnD}{\tnzoom*max(0.04, 0.04+0.055*ln((\tnch)/256)/ln(2))}
    \else
      \pgfmathsetmacro{\tnD}{\tndepth}
    \fi
    \begin{scope}[x={(1cm,0cm)}, y={(0cm,1cm)}, z={(-0.866cm,-0.5cm)}]
      % one light source: top brightest, right middle, front darkest
      \fill[tnc!30] (-\tnH,\tnH,0) -- (-\tnH,\tnH,-\tnD) -- (\tnH,\tnH,-\tnD) -- (\tnH,\tnH,0) -- cycle;
      \fill[tnc!20] (\tnH,-\tnH,0) -- (\tnH,-\tnH,-\tnD) -- (\tnH,\tnH,-\tnD) -- (\tnH,\tnH,0) -- cycle;
      \fill[tnc!11] (-\tnH,-\tnH,0) -- (\tnH,-\tnH,0) -- (\tnH,\tnH,0) -- (-\tnH,\tnH,0) -- cycle;
      \edef\tnImgFile{\tnimg}
      \ifx\tnImgFile\empty\else
        \node[inner sep=0pt, outer sep=0pt] at (0,0,0)
          {\includegraphics[width=\tnW cm, height=\tnW cm]{\assetdir/\tnimg}};
      \fi
      \ifnum\tnhidden=1
        \draw[tnc!60, line width=0.5pt, densely dashed]
          (-\tnH,-\tnH,0) -- (-\tnH,-\tnH,-\tnD) -- (\tnH,-\tnH,-\tnD)
          (-\tnH,-\tnH,-\tnD) -- (-\tnH,\tnH,-\tnD);
      \fi
      \draw[tnc!75, line width=0.5pt]
        (-\tnH,-\tnH,0) -- (\tnH,-\tnH,0) -- (\tnH,\tnH,0) -- (-\tnH,\tnH,0) -- cycle
        (-\tnH,\tnH,0) -- (-\tnH,\tnH,-\tnD) -- (\tnH,\tnH,-\tnD) -- (\tnH,\tnH,0)
        (\tnH,\tnH,-\tnD) -- (\tnH,-\tnH,-\tnD) -- (\tnH,-\tnH,0);
      \coordinate (-front) at (0,0,0);
      \coordinate (-back)  at (0,0,-\tnD);
      \coordinate (-west)  at (-\tnH,0,0);
      \coordinate (-east)  at (\tnH,0,0);
      \coordinate (-westc) at (-\tnH,0,-0.5*\tnD);
      \coordinate (-eastc) at (\tnH,0,-0.5*\tnD);
      \coordinate (-north) at (0,\tnH,0);
      \coordinate (-south) at (0,-\tnH,0);
      \coordinate (-nw) at (-\tnH,\tnH,0);  \coordinate (-ne) at (\tnH,\tnH,0);
      \coordinate (-sw) at (-\tnH,-\tnH,0); \coordinate (-se) at (\tnH,-\tnH,0);
      \coordinate (-topne) at (\tnH,\tnH,-\tnD);
    \end{scope}
    \edef\tnCap{\tncaption}
    \ifx\tnCap\empty\else
      \node[sub, anchor=north, inner sep=1.5pt] at (0,-\tnH) {\tncaption};
    \fi
  },
  % a pile of k offset H x W planes (mask stack); front plane may carry a raster
  %   \pic (m) at (x,y) {tensorpile={px=800, ch=3, img=pred_masks.png}};
  % here ch = number of planes drawn.
  tensorpile/.pic={
    \tikzset{/tensor/.cd,#1}
    \colorlet{tnc}{\tncolor}
    \pgfmathsetmacro{\tpW}{\tnzoom*(\tnpx)/400}
    \pgfmathsetmacro{\tpH}{0.5*\tpW}
    \pgfmathsetmacro{\tpO}{0.06*\tnzoom}
    \pgfmathtruncatemacro{\tpN}{\tnch-1}
    \foreach \i in {\tpN,...,0} {
      \pgfmathsetmacro{\d}{\i*\tpO}
      \fill[white]  ({-\tpH+\d},{-\tpH+\d}) rectangle ({\tpH+\d},{\tpH+\d});
      \fill[tnc!8]  ({-\tpH+\d},{-\tpH+\d}) rectangle ({\tpH+\d},{\tpH+\d});
      \ifnum\i=0 \edef\tpImg{\tnimg}\ifx\tpImg\empty\else
        \node[inner sep=0pt, outer sep=0pt] at (0,0)
          {\includegraphics[width=\tpW cm, height=\tpW cm]{\assetdir/\tnimg}};\fi\fi
      \draw[tnc!60, line width=0.5pt] ({-\tpH+\d},{-\tpH+\d}) rectangle ({\tpH+\d},{\tpH+\d});
    }
    \coordinate (-west)  at (-\tpH,0);
    \coordinate (-east)  at ({\tpH+\tpN*\tpO},0);
    \coordinate (-north) at (0,{\tpH+\tpN*\tpO});
    \coordinate (-south) at (0,-\tpH);
    \coordinate (-front) at (0,0);
    \edef\tnCap{\tncaption}
    \ifx\tnCap\empty\else
      \node[sub, anchor=north, inner sep=1.5pt] at (0,-\tpH) {\tncaption};
    \fi
  },
  /tensor/.search also={/tikz},
  /tensor/.cd,
  px/.store in=\tnpx,           px=200,
  ch/.store in=\tnch,           ch=256,
  zoom/.store in=\tnzoom,       zoom=1,
  color/.store in=\tncolor,     color=datablue,
  img/.store in=\tnimg,         img=,
  depth/.store in=\tndepth,     depth=,
  hidden/.store in=\tnhidden,   hidden=0,
  caption/.store in=\tncaption, caption=,
}


% ---- the three local overrides of D1/D2 above, and one stroke-budget fix --
% tikz_style.tex draws the lock shackle at 0.65 pt, a third line weight;
% SPEC 2.2 allows 0.5 / 0.8 only, so the glyph is restated at 0.5 pt.
\tikzset{
  carrier/.append style={fill=white, dash pattern=on 2pt off 1.2pt},
  grp/.append style     ={draw=frozgray},
  lock/.pic={
    \fill[frozgray] (-0.105,-0.105) rectangle (0.105,0.055);
    \draw[frozgray, line width=0.5pt]
      (-0.060,0.055) -- (-0.060,0.105)
      arc[start angle=180, end angle=0, radius=0.060] -- (0.060,0.055);
    \fill[white] (0,-0.022) circle[radius=0.020];
  },
}

\begin{document}
\begin{tikzpicture}[fig]
\figcanvasd{4.00}

\def\yS{2.30}      % the one strip height: every block and both rasters
\def\yG{2.90}      % baseline of the three region names

% =========================================================================
%  REGIONS (H9): three dotted containers, names set outside, nesting = 1
% =========================================================================
\draw[grp] (2.35,1.75) rectangle (4.95,2.85);
\draw[grp] (5.05,1.75) rectangle (11.35,2.85);
\draw[grp] (11.40,1.75) rectangle (15.65,2.85);
\node[grptitle, anchor=south] at (3.65,\yG) {Stage A: frozen base};
\node[grptitle, anchor=south] at (8.20,\yG) {Stage B, rung 1: training-free};
\node[grptitle, anchor=south] at (13.40,\yG)
  {Stage B, rungs 2--3: trained correctors};

% =========================================================================
%  INPUT -- the real 800 x 800 model input (assets/chip.png, 1016 dpi)
% =========================================================================
\realpanel{1.15}{\yS}{800}{1}{chip.png}{datablue!60, line width=0.5pt}
\node[sub, anchor=north] at (1.15,1.22) {SAR chip\\ $800\times800$};

% =========================================================================
%  STAGE A -- the frozen base.  One lock in the whole figure (SPEC 2.3).
% =========================================================================
\node[froz, minimum width=24mm, minimum height=9mm, inner xsep=1.5pt]
  (base) at (3.65,\yS)
  {\textbf{Mask2Former--R50}\\[0.5pt]
   {\scriptsize frozen, 3 seeds}};
\fill[white] ([shift={(-0.20,-0.10)}]base.north east)
             rectangle ([shift={(0,0.08)}]base.north east);
\pic[scale=0.85] at ([shift={(-0.10,-0.01)}]base.north east) {lock};

% =========================================================================
%  STAGE B, RUNG 1 -- two carriers (established practice) feeding the one
%  operator this article contributes.
% =========================================================================
\node[carrier, minimum width=18mm, minimum height=9mm] (soup) at (6.05,\yS)
  {\textbf{Weight soup}\\[0.5pt] {\scriptsize weight space}};
\node[carrier, minimum width=18mm, minimum height=9mm] (tta) at (8.10,\yS)
  {\textbf{Wide TTA}\\[0.5pt] {\scriptsize input space}};
\node[ops, minimum width=20mm, minimum height=9mm, inner xsep=2pt]
  (vote) at (10.25,\yS)
  {\textbf{Consensus vote}\\[0.5pt] {\scriptsize output space}};

% =========================================================================
%  STAGE B, RUNGS 2-3 -- trained correctors on disjoint write fields
% =========================================================================
\node[trn, minimum width=22.5mm, minimum height=9mm, inner xsep=2pt]
  (ref) at (12.625,\yS)
  {\textbf{Boundary refiner}\\[0.5pt] {\scriptsize writes mask pixels}};
\node[trn, minimum width=15.5mm, minimum height=9mm] (rr) at (14.775,\yS)
  {\textbf{Re-ranker}\\[0.5pt] {\scriptsize writes scores}};

% =========================================================================
%  OUTPUT -- the real fused instance masks (assets/pred_masks.png, 1016 dpi)
% =========================================================================
\pic (out) at (16.85,\yS)
  {tensorpile={px=800, ch=3, color=datablue, img=pred_masks.png}};
\node[sub, anchor=north] at (16.85,1.22) {instances\\ $\{(\mathbf{m},s)\}$};

% =========================================================================
%  ONE connection idiom: 0.5 pt Stealth flow arrows, all on one baseline
% =========================================================================
\draw[flow] (2.15,\yS)   -- (base.west);
\draw[flow] (base.east)  -- (soup.west);
\draw[flow] (soup.east)  -- (tta.west);
\draw[flow] (tta.east)   -- (vote.west);
\draw[flow] (vote.east)  -- (ref.west);
\draw[flow] (ref.east)   -- (rr.west);
\draw[flow] (rr.east)    -- (out-west);

\end{tikzpicture}
\end{document}
